% !TeX program = xelatex
% Revisão editorial e normativa em 10/10/2026; dados autorais pendentes.
\documentclass[12pt,openany,oneside,a4paper,brazil,chapter=TITLE]{abntex2}
\usepackage{fontspec}
\setmainfont{TeX Gyre Termes}
\setsansfont{TeX Gyre Termes}
\setmonofont{DejaVu Sans Mono}[Scale=0.82]
\usepackage{amsmath,amssymb,booktabs,tabularx,longtable,array,verbatim,listings}
\usepackage[authoryear,round]{natbib}
\usepackage{xurl}
\setlength{\emergencystretch}{2em}
\setcitestyle{aysep={,},citesep={;}}
\let\cite\citep
\renewcommand{\bibsection}{\chapter*{Referências}\addcontentsline{toc}{chapter}{REFERÊNCIAS}}
\setlength{\bibhang}{0pt}
\setlength{\bibsep}{\baselineskip}
\renewcommand{\ABNTEXchapterfont}{\normalfont\bfseries}
\renewcommand{\ABNTEXchapterfontsize}{\normalsize}
\renewcommand{\ABNTEXsectionfont}{\normalfont\bfseries}
\renewcommand{\ABNTEXsectionfontsize}{\normalsize}
\renewcommand{\ABNTEXsubsectionfont}{\normalfont\bfseries}
\renewcommand{\ABNTEXsubsectionfontsize}{\normalsize}
\setlength{\beforechapskip}{0pt}
\setlength{\afterchapskip}{1.5\baselineskip}
\setlrmarginsandblock{3cm}{2cm}{*}
\setulmarginsandblock{3cm}{2cm}{*}
\checkandfixthelayout[fixed]
\setlength{\parindent}{1.25cm}
\setlength{\parskip}{0pt}
\OnehalfSpacing
\makepagestyle{abntsubmissao}
\makeoddhead{abntsubmissao}{}{}{\fontsize{10}{12}\selectfont\thepage}
\makeevenhead{abntsubmissao}{}{}{\fontsize{10}{12}\selectfont\thepage}
\makeoddfoot{abntsubmissao}{}{}{}
\makeevenfoot{abntsubmissao}{}{}{}
\captionnamefont{\fontsize{10}{12}\selectfont}
\captiontitlefont{\fontsize{10}{12}\selectfont}
\captiondelim{ --- }
\captionstyle{\raggedright}
\setlength{\tabcolsep}{4pt}
\lstset{basicstyle=\fontsize{10}{12}\selectfont\ttfamily,breaklines=true,
breakatwhitespace=false,columns=fullflexible,keepspaces=true,showstringspaces=false,
literate={á}{{\'a}}1 {à}{{\`a}}1 {â}{{\^a}}1 {ã}{{\~a}}1 {é}{{\'e}}1 {è}{{\`e}}1 {ê}{{\^e}}1 {í}{{\'i}}1 {ó}{{\'o}}1 {ô}{{\^o}}1 {õ}{{\~o}}1 {ú}{{\'u}}1 {ç}{{\c c}}1 {Á}{{\'A}}1 {É}{{\'E}}1 {Í}{{\'I}}1 {Ó}{{\'O}}1 {Õ}{{\~O}}1 {Ú}{{\'U}}1 {Ç}{{\c C}}1 {—}{{---}}1 {–}{{--}}1 {°}{{\textdegree}}1 {±}{{$\pm$}}1 {×}{{$\times$}}1 {→}{{$\to$}}1 {≤}{{$\le$}}1 {≥}{{$\ge$}}1 {≠}{{$\ne$}}1}
\titulo{Roteamento inspirado em atenção: formalização do núcleo matemático em Lean 4 e avaliação empírica interna}
\autor{Marcelo Claro Laranjeira}
\local{Porto Alegre}
\data{2026}
\orientador{[nome e titulação do orientador --- preencher]}
\instituicao{Pontifícia Universidade Católica do Rio Grande do Sul}
\tipotrabalho{Dissertação (Mestrado)}
\preambulo{Dissertação apresentada ao Programa de Pós-Graduação em [nome do programa --- preencher] da Pontifícia Universidade Católica do Rio Grande do Sul como requisito parcial para obtenção do título de Mestre em [denominação do título --- preencher]. Área de concentração: [preencher]. Linha de pesquisa: [preencher].}
\hypersetup{hidelinks,pdftitle={Roteamento inspirado em atenção: formalização do núcleo matemático em Lean 4 e avaliação empírica interna},pdfauthor={Marcelo Claro Laranjeira},pdfsubject={Dissertação: revisão editorial e normativa para submissão},pdfkeywords={Lean 4; verificação formal; roteamento; atenção; avaliação empírica}}
\begin{document}
\frenchspacing
\hypersetup{pageanchor=false}
\begin{capa}
\begin{center}
\imprimirinstituicao\par
Programa de Pós-Graduação em [nome do programa --- preencher]\par
\vspace{2cm}
\imprimirautor\par
\vfill
\textbf{\MakeUppercase{\imprimirtitulo}}\par
\vfill
\imprimirlocal\par
\imprimirdata
\end{center}
\end{capa}
\setcounter{page}{1}
\begin{folhaderosto}
\begin{center}
\imprimirautor\par
\vfill
\textbf{\imprimirtitulo}\par
\vfill
\end{center}
\noindent\hspace*{0.5\textwidth}\begin{minipage}{0.5\textwidth}
\SingleSpacing
\imprimirpreambulo\par\vspace{\baselineskip}
Orientação: \imprimirorientador.
\end{minipage}
\vfill
\begin{center}\imprimirlocal\par\imprimirdata\end{center}
\end{folhaderosto}
% Ficha imediatamente após o rosto; esta página não entra na contagem.
\thispagestyle{empty}\vspace*{\fill}
\begin{center}\begin{minipage}{0.8\textwidth}\SingleSpacing
\textbf{Ficha catalográfica}\par
[Inserir a ficha emitida pelo serviço de catalogação da instituição após a confirmação dos dados autorais e da versão final.]
\end{minipage}\end{center}
\vspace*{\fill}\clearpage\addtocounter{page}{-1}
\begin{folhadeaprovacao}
\begin{center}\imprimirautor\par\vspace{\baselineskip}\textbf{\imprimirtitulo}\end{center}
\vspace{\baselineskip}
\noindent\hspace*{0.5\textwidth}\begin{minipage}{0.5\textwidth}\SingleSpacing\imprimirpreambulo\end{minipage}
\vspace{2\baselineskip}
\noindent Data da aprovação: [preencher somente após a defesa].\par
\noindent Banca examinadora: [nomes, titulações, instituições e assinaturas --- preencher].\par
\vspace{2\baselineskip}
\noindent [A aprovação e as assinaturas ainda não foram registradas nesta versão.]
\end{folhadeaprovacao}
\begin{agradecimentos}
Às comunidades Lean e Mathlib, pela manutenção das ferramentas e bibliotecas utilizadas, e às iniciativas de ciência aberta que apoiam a documentação e a consulta das fontes deste trabalho.
\end{agradecimentos}
\begin{resumo}
Este trabalho investiga a formalização do núcleo matemático e a avaliação empírica interna de um mecanismo de roteamento de tarefas entre agentes de software, inspirado na ponderação por atenção. A pesquisa combina provas em Lean 4 com Mathlib, inspeção dos artefatos de execução, uma bancada sintética e testes aleatórios de propriedades numéricas. Os registros locais documentam onze lemas sobre normalização do softmax, limites de combinações convexas, existência de maximizador finito, invariância a deslocamento, pesos positivos, caso de candidato único, monotonicidade, ordenação top-\textit{k}, temperatura e limite superior da exponencial após deslocamento. As provas referem-se à aritmética real exata e não certificam integralmente a implementação em ponto flutuante. Na bancada principal, com 200 decisões distribuídas por cinco sementes, o arrependimento médio registrado foi $0{,}042$, com margem exploratória de $0{,}011$, frente a valores entre $0{,}19$ e $0{,}22$ dos comparadores simples. Ablações em um conjunto distinto de instâncias indicam dependência da confiança, construída como proxy da qualidade latente. A variação do parâmetro do oráculo apenas reescala a perda entre candidatos plenamente elegíveis. Os intervalos desconsideram o agrupamento por piscina de agentes e devem ser interpretados como resumos exploratórios. O relatório numérico registra 20\,000 vetores sem violações acima das tolerâncias examinadas. Conclui-se que os artefatos sustentam propriedades do modelo abstrato e resultados preliminares no simulador, permanecendo pendentes a certificação numérica completa, a análise estatística com agrupamento, a comparação externa e a reprodução independente.
\par\vspace{\baselineskip}
\noindent\textbf{Palavras-chave:} roteamento de agentes; atenção; Lean 4; verificação formal; avaliação empírica; reprodutibilidade.
\end{resumo}
\begin{resumo}[Abstract]
\begin{otherlanguage*}{english}
This dissertation investigates the formalization of the mathematical core and the internal empirical evaluation of a task router for software agents inspired by attention-based weighting. The study combines proofs in Lean 4 with Mathlib, inspection of execution artifacts, a synthetic benchmark, and randomized tests of numerical properties. Local records document eleven lemmas on softmax normalization, convex-combination bounds, existence of a finite maximizer, shift invariance, positive weights, the singleton case, monotonicity, top-\textit{k} ordering, temperature, and an upper bound on shifted exponentials. The proofs concern exact real arithmetic and do not fully certify the floating-point implementation. In the main benchmark, comprising 200 decisions across five seeds, recorded mean regret was $0.042$, with an exploratory margin of $0.011$, compared with values between $0.19$ and $0.22$ for simple baselines. Ablations on a distinct set of instances indicate dependence on confidence, which was constructed as a proxy for latent quality. Changing the oracle parameter only rescales regret among fully eligible candidates. The reported intervals disregard clustering within agent pools and should be interpreted as exploratory summaries. The numerical report records 20,000 vectors with no violations above the tested tolerances. The artifacts support properties of the abstract model and preliminary findings within the simulator. Full numerical certification, statistical analysis accounting for clustering, external comparison, and independent reproduction remain pending.
\par\vspace{\baselineskip}
\noindent\textbf{Keywords:} agent routing; attention; Lean 4; formal verification; empirical evaluation; reproducibility.
\end{otherlanguage*}
\end{resumo}
\pdfbookmark[0]{Lista de tabelas}{lot}
\listoftables*
\begin{siglas}
\item[ABNT] Associação Brasileira de Normas Técnicas
\item[IC] Intervalo de confiança
\item[IEEE] Institute of Electrical and Electronics Engineers
\item[JSON] JavaScript Object Notation
\item[MCP] Model Context Protocol
\item[OSF] Open Science Framework
\item[PUCRS] Pontifícia Universidade Católica do Rio Grande do Sul
\item[Q/K/V] Consulta, chave e valor (\textit{query, key, value})
\item[SDD] Desenvolvimento orientado a especificação
\item[SMT] Satisfatibilidade módulo teorias
\item[TDD] Desenvolvimento orientado a testes
\end{siglas}
\begin{simbolos}
\item[$w_i$] Peso de atenção do candidato $i$
\item[$f_i$] Escore do candidato $i$
\item[$\tau$] Temperatura positiva do softmax
\item[$q$] Qualidade latente do agente no simulador
\item[$q_w$] Peso da qualidade no oráculo sintético
\item[$u(a)$] Utilidade combinada do candidato $a$
\item[$\alpha_h$] Peso fixo da cabeça $h$
\end{simbolos}
\pdfbookmark[0]{Sumário}{toc}
\tableofcontents*
\clearpage
\textual
\pagestyle{abntsubmissao}\aliaspagestyle{chapter}{abntsubmissao}
\hypersetup{pageanchor=true}
% Módulo 01 — Introdução (R775+; problema, justificativa, estrutura)
\chapter{Introdução}\label{cap:intro}
\section{O problema}
Sistemas que delegam tarefas a agentes de software decidem sob incerteza com critérios múltiplos, e essas decisões raramente carregam garantias auditáveis: não se sabe demonstrar que os pesos somam um, que a regra escolhe o melhor pelos critérios declarados, nem que a medição de qualidade significa o que diz. A engenharia de software experimental estabeleceu há décadas que alegações exigem desenho, medição e relato disciplinados \cite{basili1986,wohlin2012,runeson2009}, e que revisões sistemáticas exigem protocolo de relato \cite{page2021}; a crise de reprodutibilidade documentada em múltiplas áreas \cite{osc2015,ioannidis2005,collberg2016} mostra o custo de evidência fraca. Este trabalho aplica esse padrão a um roteador por atenção: cada propriedade alegada ou é provada por máquina, ou é medida com incerteza declarada, ou é renunciada explicitamente.
\section{Questão e objetivos}
Que propriedades do núcleo matemático de um roteador inspirado em atenção podem ser formalizadas, e que comportamento sua implementação apresenta em uma avaliação sintética interna? Objetivo geral: formalizar e avaliar internamente um roteamento por atenção. Objetivos específicos: (OE1) verificar por máquina o núcleo matemático em Lean 4 com Mathlib \cite{moura2021,mathlib2020}; (OE2) comparar arrependimento e taxa de acerto com os comparadores, relatando intervalos exploratórios, seguindo a recomendação de estimação em vez de significância pontual \cite{cumming2014,wasserstein2016}; (OE3) auditar axiomas, alegações, referências e código, incluindo inspeção e métricas na tradição de \cite{fagan1976,mccabe1976,chidamber1994}; (OE4) declarar limites e publicar a trilha de reprodução segundo princípios FAIR \cite{wilkinson2016} e documentação de dados e modelos \cite{gebru2021,mitchell2019}.
\section{Justificativa}
Roteamento opaco não é auditável; garantias formais exigem hipóteses explícitas; medições requerem análise de incerteza. Como perspectiva educacional, trilhas auditáveis desse tipo oferecem material didático para o ensino de avaliação e validação --- hipótese de trabalho futuro, não resultado desta dissertação.
\section{Delimitação e tese defendida}
Defende-se exclusivamente otimalidade relativa à utilidade definida e evidência preliminar interna. Não se compara com plataformas existentes, não se generaliza além do simulador descrito e não se transfere propriedades de Transformers treinados. A tese metodológica é que verificação por máquina, teste numérico de propriedades, ablação e auditoria de alegações compõem um padrão de evidência auditável mesmo sem validação externa.
\section{Estrutura do trabalho}
A Seção~\ref{cap:preliminares} fixa preliminares matemáticos; a Seção~\ref{cap:fundamentacao} fundamenta atenção, correspondência Q/K/V, combinação convexa, seleção gulosa, Lean e os trabalhos relacionados; a Seção~\ref{cap:metodo} detalha método, estatística, ameaças e validade; a Seção~\ref{cap:provas} apresenta as onze provas comentadas; a Seção~\ref{cap:resultados}, os testes de propriedades e a avaliação empírica; a Seção~\ref{cap:discussao} discute revisões e riscos; a Seção~\ref{cap:arguicoes} antecipa arguições e responde nos limites do texto; a Seção~\ref{cap:conclusao} conclui; o glossário define os termos principais; os apêndices preservam código, resultados, transcrições e especificações.
\section{Cenário-guia de roteamento (saídas registradas)}
Tarefa ``provar teorema formal com código Lean'' exigindo \texttt{formal} e \texttt{codigo}, com três agentes: Ana (prova e código; carga 0,1; confiança 0,9), Bia (código e dados; carga 0,8; confiança 0,5) e Cid (as três capacidades; carga 0,4; confiança 0,7). O gate exclui Bia (falta \texttt{formal}); utilidades: Ana $0{,}9175$, Cid $0{,}8161$; pesos finais Ana $0{,}5253$ e Cid $0{,}4747$ (soma 1). Ana vence com menos capacidades porque combina menor carga, maior confiança e melhor ajuste semântico --- exatamente o tipo de compensação que os teoremas disciplinam e a bancada mede.
\par Dois casos-limite reais: (i) agentes gêmeos que diferem só na descrição recebem $0{,}5053$ contra $0{,}4947$ --- nem gêmeos empatam na semântica, e empates exatos decidem-se pelo identificador, deterministicamente; (ii) candidato único recebe peso $1{,}0$ --- o T7 em execução.

% Módulo 02a — Preliminares matemáticos e computacionais (R775+)
\chapter{Preliminares matemáticos e computacionais}\label{cap:preliminares}
\section{Informação, entropia e exponenciação}
Os valores numéricos de pesos apresentados no texto são arredondados a quatro casas decimais; a soma unitária refere-se aos valores antes desse arredondamento.
A teoria matemática da comunicação \cite{shannon1948} introduziu a entropia como medida de incerteza; a relação entre informação e suficiência foi formalizada por \citet{kullback1951}; e a conexão entre teoria da informação e mecânica estatística, com a distribuição exponencial como maximizadora de entropia sob média fixada, por \citet{jaynes1957}. É dessa linhagem que descende a ponderação exponencial normalizada: sob uma restrição adequada de média, a solução tem forma $w_i\propto e^{\beta f_i}$, com $\beta$ determinado pela restrição; o softmax usual corresponde ao caso particular $\beta=1$. A interpretação probabilística de saídas de redes feedforward, com relações ao reconhecimento estatístico de padrões, foi estabelecida por \citet{bridle1990}. Regras de escore estritamente próprias, que incentivam probabilidades calibradas, foram caracterizadas por \citet{gneiting2007}.
\section{Conjuntos convexos e ordens finitas}
Uma combinação $\sum w_ix_i$ com $w_i\geq0$ e $\sum w_i=1$ jamais escapa ao intervalo $[\min x_i,\max x_i]$: é o conteúdo de T2/T6, e a razão formal pela qual pesos devem somar um. Em conjuntos finitos não vazios sempre existe elemento maximal de qualquer função real (T3): finitude elimina a possibilidade de fuga ao infinito que invalida máximos em domínios infinitos. Listas ordenadas de forma decrescente têm a propriedade de que qualquer prefixo domina o sufixo (T9): é a anatomia do top-\textit{k}.
\section{Atenção neural e tradução automática}
A atenção produto-escalado com múltiplas cabeças foi introduzida para tradução automática \cite{vaswani2017}, com variantes efetivas exploradas em tradução neural \cite{luong2015}, pré-treino bidirecional \cite{devlin2019}, dependências longas \cite{dai2019}, benchmarks multitarefa \cite{wang2018}, compreensão de leitura em larga escala \cite{rajpurkar2016}, bases hierárquicas de imagens \cite{deng2009} e ferramental de processamento \cite{wolf2020}. Este trabalho herda apenas a forma --- ponderação normalizada de cabeças com decisão pelo máximo --- sem pesos treinados.
\section{Verificação, SMT e inspeção}
Bases axiomáticas para programação \cite{hoare1969} fundam a verificação dedutiva mecanizada em Lean \cite{moura2021} e Mathlib \cite{mathlib2020}; solucionadores SMT eficientes como o Z3 \cite{demoura2008} e computação simbólica como SymPy \cite{meurer2017} compõem o ecossistema de apoio. Inspeções de projeto e código reduzem erros antes da execução \cite{fagan1976}; complexidade ciclomática \cite{mccabe1976} e suítes de métricas orientadas a objetos \cite{chidamber1994} disciplinam a auditoria de código; processos racionais de projeto documentam decisões \cite{parnas1986}. Alocação de tarefas distribuídas segue desde o protocolo de rede de contratos \cite{smith1980}, e integração de conhecimento sob incerteza desde sistemas de quadro-negro \cite{erman1980}.
\section{Ponto flutuante e computação científica}
A aritmética de ponto flutuante exige atenção a arredondamento e casos especiais \cite{goldberg1991}; o manual de referência consolida algoritmos e análise \cite{muller2018}. A computação científica em Python apoia-se em arranjos \cite{harris2020} e algoritmos fundamentais \cite{virtanen2020}.
\section{Método empírico e estatística}
Experimentação disciplinada em engenharia de software \cite{basili1986,wohlin2012}, diretrizes de estudo de caso \cite{runeson2009} e relato sistemático \cite{page2021} estruturam o desenho; a declaração sobre valores-\textit{p} \cite{wasserstein2016}, a nova estatística por estimação \cite{cumming2014} e o bootstrap \cite{efron1979} fundamentam o relato por intervalos em vez de significância pontual. Reprodutibilidade como problema de campo \cite{osc2015,collberg2016,ioannidis2005}, dados documentados \cite{gebru2021}, modelos documentados \cite{mitchell2019} e dados geridos por FAIR \cite{wilkinson2016} justificam a trilha local de reprodução.
\section{Exemplos de cálculo}
\textit{Normalização.} Com $S=\sum_j e^{f_j}>0$, cada quociente $e^{f_i}/S$ é positivo e a soma com denominador comum reduz-se a $S/S=1$. \textit{Deslocamento.} Subtrair $c$ multiplica numerador e denominador por $1/e^c$, que cancela: a decisão independe do nível. \textit{Temperatura.} Para $\tau_1<\tau_2$ e maximizador único $i^*$, a razão de pesos cresce porque cada termo concorrente $e^{(f_j-f_{i^*})/\tau}$ decai mais rápido no $\tau$ menor --- a Tabela~\ref{tab:temp} mostra $0{,}5065\to0{,}9817$. \textit{Limite da exponencial.} Em reais exatos, deslocar pelo máximo torna cada expoente não positivo e cada exponencial pertence a $(0,1]$. T11 demonstra esse limite; em ponto flutuante, underflow e arredondamento requerem análise própria.
\section{Top-\textit{k} e complexidade}
Extrair os $k$ maiores por heap custa $O(n\log k)$ contra $O(n\log n)$ da ordenação total; o desempate por identificador torna a decisão determinística. A garantia independe do algoritmo: sob ordem decrescente, o prefixo domina o sufixo.

% Módulo 02f — Núcleo matemático: MaxEnt, cálculo softmax, temperatura, top-k (R776+)
\section{Cálculo da ponderação exponencial}
A entropia de uma distribuição, na teoria matemática da comunicação \cite{shannon1948}, mede incerteza; dada média fixada dos escores, a distribuição de máxima entropia tem forma exponencial \cite{jaynes1957} --- é a origem variacional do softmax, cuja interpretação probabilística em classificação foi estabelecida em \cite{bridle1990}. Três cálculos sustentam o uso: (i) normalização, $\sum_i w_i=1$ por construção (T1); (ii) invariância a deslocamento (T4), que autoriza a subtração do máximo; (iii) limite superior da exponencial após deslocamento por um majorante (T11). A Tabela~\ref{tab:temp} mostra a nitidez por temperatura em escores $[1,2,3]$: o peso máximo sobe $0{,}5065\to0{,}6652\to0{,}8668\to0{,}9817$ ao esfriar --- uma ilustração numérica de T10.
\begin{table}[htbp]
\fontsize{10}{12}\selectfont\SingleSpacing
\centering
\caption{Concentração por temperatura em $[1,2,3]$ (calculado pela implementação).}
\label{tab:temp}
\begin{tabular}{lccc}
\toprule
$\tau$ & $w_1$ & $w_2$ & $w_3$ \\
\midrule
2,0 & $0{,}1863$ & $0{,}3072$ & $0{,}5065$ \\
1,0 & $0{,}0900$ & $0{,}2447$ & $0{,}6652$ \\
0,5 & $0{,}0159$ & $0{,}1173$ & $0{,}8668$ \\
0,25 & $0{,}0003$ & $0{,}0180$ & $0{,}9817$ \\
\bottomrule
\end{tabular}
\fonte{Elaboração própria (2026), a partir dos artefatos e das fontes identificados no texto.}
\end{table}
\section{Ordenação e top-\textit{k}}
Ordenar $n$ utilidades custa $O(n\log n)$; extrair os $k$ maiores por heap custa $O(n\log k)$. A implementação ordena tudo e fatia o topo; com 40 agentes e latência de $2{,}17$\,ms, a seleção parcial segue como melhoria medida pendente. O desempate é determinístico (menor identificador vence), o que torna a seleção determinística quando entradas, ambiente e demais operações permanecem fixos. A garantia matemática independe do algoritmo: sob ordem pairwise decrescente, o prefixo domina o sufixo (T9).

% Módulo 02g — Curso embutido: cálculo, floats, estatística, Lean, orquestração (R776+)
\section{Derivação variacional (MaxEnt)}
Maximizar $H(w)=-\sum w_i\ln w_i$ sob $\sum w_i=1$ e média fixada, por multiplicadores de Lagrange, produz $w_i\propto e^{\beta f_i}$, com $\beta$ determinado pela média imposta: a forma exponencial não é arbitrária, é a distribuição de máxima entropia compatível com a média \cite{jaynes1957,shannon1948}. Com $\beta=1$ recupera-se o softmax; $\beta=1/\tau$ introduz a temperatura. Regras de escore próprias caracterizam quando probabilidades relatadas são honestas \cite{gneiting2007}.
\section{Limites de temperatura}
Com maximizador único $i^*$: quando $\tau\to0^+$, cada razão $e^{(f_j-f_{i^*})/\tau}\to0$ e $w_{i^*}\to1$; quando $\tau\to\infty$, todas as razões tendem a 1 e os pesos equalizam em $1/n$. T10 é a versão finita e verificada dessa intuição: esfriar aumenta estritamente o peso do favorito, sem passar ao limite.
\section{Anatomia do ponto flutuante}
IEEE-754 representa $x\approx m\times2^e$ com arredondamento ao mais próximo; operações elementares têm erro relativo limitado pela unidade de arredondamento, mas exponenciação de argumentos grandes estoura para infinito \cite{goldberg1991,muller2018}. Subtrair o máximo antes de exponenciar mantém todo expoente $\leq0$ (T11), logo cada termo fica em $(0,1]$ em aritmética real exata. Esse limite não certifica as operações IEEE-754: exponenciais pequenas podem sofrer underflow, e a invariância exata de T4 não implica igualdade bit a bit para quaisquer entradas.
\section{Intervalos, Pearson e pareamento}
Para média $\bar{x}$ de $n$ observações com desvio $s$, o IC95 $\bar{x}\pm1{,}96s/\sqrt{n}$ aproxima a incerteza sob normalidade; o relato por intervalos segue a estatística por estimação \cite{cumming2014} e a advertência contra o uso pontual de valores-\textit{p} \cite{wasserstein2016}; reamostragem bootstrap \cite{efron1979} é a alternativa quando a normalidade é suspeita. Pearson mede relação linear; parear estratégias nas mesmas decisões controla diferenças de dificuldade entre comparadores. Regras de escore próprias \cite{gneiting2007} lembram que probabilidade bem relatada é hipótese, sem presunção de validade.
\section{Táticas Lean em prosa}
\textit{Positividade} (\texttt{exp\_pos}, divisão de positivos) estabelece sinais. \textit{Reescrita} (\texttt{sum\_div}, \texttt{sum\_mul}, \texttt{exp\_sub}, \texttt{exp\_add}) move termos entre formas equivalentes. \textit{Monotonicidade} (\nolinkurl{mul_le_mul_of_nonneg_left}, \texttt{div\_lt\_div\_iff₀}) preserva desigualdades. \textit{Fechamento} (\texttt{div\_self}, \texttt{linarith}, \texttt{field\_simp}, \texttt{ring}) conclui identidades. \textit{Estrutura} (\texttt{take\_append\_drop}, \texttt{pairwise\_append}, \texttt{add\_sum\_erase}, \texttt{exists\_max\_image}, \texttt{sum\_lt\_sum}, \texttt{sum\_eq\_single\_of\_mem}) decompõe objetos finitos. Cada nome foi conferido na Mathlib pinada, nunca na memória.
\section{Orquestração: do contrato ao quadro-negro}
Alocar tarefas entre resolvedores distribuídos remonta ao protocolo de rede de contratos \cite{smith1980}; integrar conhecimento heterogêneo sob incerteza, aos sistemas de quadro-negro \cite{erman1980}. Nosso roteador instancia a linhagem com meios atuais: gates decidem elegibilidade, cabeças pontuam aspectos, softmax normaliza, máximo decide --- cada camada auditável separadamente, como documenta a trilha de pinos e hashes.

% Módulo 02 — Fundamentação teórica (R762+; Q/K/V parcial honesto)
\chapter{Fundamentação teórica}\label{cap:fundamentacao}
Antes de provar ou medir, fixamos o vocabulário: o que a ponderação exponencial garante, onde ela se inspira e onde a analogia termina. A base axiomática da programação \cite{hoare1969} sustenta a verificação dedutiva; a computação numérica e simbólica executa-se sobre arranjos e algoritmos fundamentais \cite{harris2020,virtanen2020}, com apoio simbólico e SMT \cite{meurer2017,demoura2008}.
\section{Ponderação exponencial normalizada}
A soma unitária dos pesos não demonstra calibração probabilística. A avaliação deste trabalho mede arrependimento e acerto, sem ajuste de temperatura, aprendizagem de pesos ou estimação de métricas de calibração.
Dados escores finitos $f_1,\dots,f_n$, os pesos $w_i=e^{f_i}/\sum_j e^{f_j}$ somam 1 (T1), situam-se em $(0,1)$ com ao menos dois candidatos (T5) e valem 1 no singleton (T7); deslocar escores por constante nada altera (T4), autorizando a subtração do máximo por estabilidade numérica. A calibração da temperatura, amplamente estudada como parâmetro de nitidez e calibração de classificadores \cite{komisarenko2024,xuan2025,xie2024}, aparece aqui em forma elementar, demonstrada por máquina (T10): esfriar aumenta o peso do maximizador único. A verificação de programas com aritmética de ponto flutuante exige análise própria \cite{bagnara2019,murray2026}; este trabalho documenta a ponte empírica (fuzz) e um lema de limite superior da exponencial após deslocamento (T11), sem reivindicar certificação total de floats.
\section{Correspondência Q/K/V (parcial e estrutural)}
\begin{table}[htbp]
\fontsize{10}{12}\selectfont\SingleSpacing
\centering
\caption{Analogia Transformer $\leftrightarrow$ roteador (sem transferência de propriedades).}
\label{tab:qkv}
\begin{tabularx}{\textwidth}{lXX}
\toprule
Transformer & Roteador & Natureza \\
\midrule
Consulta (Q) & Vetor da tarefa & Estrutural \\
Chave (K) & Vetor do cartão do agente & Estrutural \\
Escore QK & Similaridade + 4 cabeças & Parcial (sem produto escalado/temperatura) \\
Valor (V) & (utilidade, identificador) & Parcial (sem mistura de valores) \\
Agregação & Convexa + softmax & Demonstrada (T2, T6) \\
Decisão & Máximo / top-\textit{k} & Demonstrada (T3, T9) \\
\bottomrule
\end{tabularx}
\fonte{Elaboração própria (2026), a partir dos artefatos e das fontes identificados no texto.}
\end{table}
\section{Combinação convexa e seleção gulosa}
Com pesos não negativos somando 1, a utilidade $u(a)=\sum_h\alpha_h s_h(a)$ situa-se entre mínimo e máximo dos componentes (T2, T6): conteúdo formal do validador que rejeita somas distintas de 1. Em conjunto finito não vazio há maximizador (T3); o padrão top-\textit{k} domina o resto sob ordem pairwise (T9); o peso cresce com o próprio escore (T8). Custo documentado: ordenação completa onde bastaria seleção parcial.
\section{Assistentes de prova e orquestração multiagente}
Lean 4 \cite{moura2021} e Mathlib \cite{mathlib2020} fundamentam a verificação; o mecanismo de atenção, \cite{vaswani2017}. Benchmarks de agentes e orquestração multiagente (\textit{AgentBench} \cite{agentbench2023}, \textit{PerspectiveGap} \cite{perspectivegap2026}, recompensa de orquestração \cite{orchrm2026}, arcabouços atencionais \cite{orchestrator2025}) delimitam o contexto externo: nenhum foi executado aqui, e a comparação permanece como trabalho futuro (protocolo proposto na Seção~\ref{cap:conclusao}).

% Módulo 02b — Contexto: atenção em Transformers (didático, R773)
Situada a matemática local, olhamos a fonte da inspiração para separar o que foi herdado do que foi apenas sugerido.
\section{O mecanismo de atenção original}
No Transformer \cite{vaswani2017}, cada posição de consulta $Q$ compara-se a chaves $K$ por produto escalado, $A=\mathrm{softmax}(QK^\top/\sqrt{d})$, e agrega valores $V$ por $AV$. Três ideias importam aqui: (i) escores viram distribuição de probabilidade; (ii) a divisão por $\sqrt{d}$ e a temperatura controlam a nitidez; (iii) múltiplas cabeças capturam aspectos distintos, combinados depois. Exemplo numérico real (implementação deste trabalho): $\mathrm{softmax}([1,2,3])=[0{,}0900;\ 0{,}2447;\ 0{,}6652]$, soma $1{,}0$ --- o maior escore leva dois terços da massa. A linhagem da atenção inclui abordagens efetivas para tradução neural \cite{luong2015}, pré-treino bidirecional \cite{devlin2019}, dependências longas \cite{dai2019} e ferramental consolidado \cite{wolf2020}; benchmarks como GLUE, SQuAD e ImageNet \cite{wang2018,rajpurkar2016,deng2009} fixaram o padrão de avaliação por tarefas que este trabalho adapta em escala sintética. Nosso roteador herda a forma --- ponderação normalizada de múltiplas cabeças com decisão pelo máximo --- sem herdar pesos treinados, representações ou desempenho: cada cabeça é uma heurística explícita (similaridade, cobertura, confiança, carga), e cada peso da combinação é fixo e auditável.
\section{Temperatura e calibração}
Dividir escores por $\tau>0$ antes do softmax regula a nitidez: $\tau\to0$ concentra no máximo, $\tau\to\infty$ equaliza. A literatura usa a temperatura para calibrar classificadores \cite{komisarenko2024,xuan2025,xie2024}. Nossa contribuição elementar e verificada (T10): com maximizador único, esfriar aumenta estritamente seu peso --- sem limites, sem assintótica, só álgebra de exponenciais.
\section{Ponto flutuante e certificação}
Programas com IEEE-754 exigem análise própria de arredondamento \cite{bagnara2019,murray2026}. A evidência contempla dois níveis: (i) lema de limite superior da exponencial após deslocamento por majorante (T11), em reais exatos; (ii) teste empírico de propriedades de 20\,000 vetores contra a implementação real, com pior desvio de soma $2{,}2\times10^{-16}$ e sem violações da invariância a deslocamento acima da tolerância examinada. Certificação total de floats segue futura.

% Módulo 02c — Tutorial didático de Lean e Mathlib (R773)
Da matemática à máquina: como essas ideias viram provas conferidas por computador.
\section{Verificação formal para leigos}
Um assistente de provas é um corretor implacável: cada passo lógico é conferido pelo núcleo até axiomas explícitos. Em Lean 4 \cite{moura2021}, teoremas são programas cujo tipo é a proposição; \textit{táticas} (\texttt{simp}, \texttt{ring}, \texttt{omega}, \texttt{linarith}, \texttt{exact}) são comandos que transformam objetivos. A Mathlib \cite{mathlib2020} é a biblioteca comunitária com milhares de lemas reusados aqui (somas finitas, exponencial real, máximos). \texttt{Sorry} e \texttt{admit} são atalhos que afirmam sem provar --- bloqueados neste trabalho em duas camadas. \texttt{\#print axioms} lista os axiomas de cada teorema: todos os onze usam só \textit{propext}, \textit{Classical.choice} e \textit{Quot.sound}, o ambiente escolhido e utilizado neste trabalho, sem atalhos.
\section{Ler uma prova deste trabalho}
Cada prova combina: fatos de positividade (exponencial sempre positiva), reescrita de somas (tirar denominador comum de somatórios), monotonicidade (multiplicar desigualdade por número não negativo preserva o sentido) e linearidade (combinar desigualdades). Exemplo guiado (T2): para mostrar que a média ponderada não passa do máximo $M$, majoramos cada termo $w_ix_i$ por $w_iM$ (válido pois $w_i\geq0$), somamos as $n$ desigualdades e fatoramos $M$ para fora da soma, que vale 1 por hipótese.
A Tabela~\ref{tab:taticas} resume os recursos utilizados nas provas.
\begin{table}[htbp]
\fontsize{10}{12}\selectfont\SingleSpacing
\centering
\caption{Táticas usadas nas onze provas (guia de leitura).}
\label{tab:taticas}\begin{tabularx}{\textwidth}{XX}
\toprule
Tática/lema & Papel \\
\midrule
\texttt{Finset.sum\_pos} & positividade de somatórios \\
\texttt{Finset.sum\_div} / \texttt{sum\_mul} & reescrita de somas com divisão/multiplicação \\
\texttt{div\_self} / \texttt{div\_lt\_div\_iff₀} & quocientes: igualdade e desigualdade \\
\nolinkurl{mul_le_mul_of_nonneg_left} & monotonicidade por não negativos \\
\texttt{linarith} & combinação linear de desigualdades \\
\texttt{field\_simp} & limpeza de denominadores \\
\texttt{Real.exp\_*} & álgebra da exponencial \\
\nolinkurl{Finset.exists_max_image} & maximizador finito \\
\texttt{take\_append\_drop} / \texttt{pairwise\_append} & anatomia do top-\textit{k} \\
\bottomrule
\end{tabularx}
\fonte{Elaboração própria (2026), a partir dos artefatos e das fontes identificados no texto.}
\end{table}
\section{Táticas de uso frequente}
\textit{Introduzir e instanciar:} \texttt{intro} traz hipóteses; \texttt{exact} entrega a prova pronta; \texttt{apply} ajusta conclusões; \texttt{refine} deixa buracos honestos. \textit{Reescrever:} \texttt{rw} substitui iguais; \texttt{simp} simplifica com a base; \texttt{ring} e \texttt{omega} decidem anéis e inteiros; \texttt{decide} computa verdades decidíveis. \textit{Estruturar:} \texttt{cases} e \texttt{induction} quebram formas; \texttt{constructor} monta pares e existências; \texttt{obtain} extrai testemunhas. \textit{Fechar:} \texttt{linarith} combina lineares; \texttt{field\_simp} limpa denominadores; \texttt{trivial} e \texttt{assumption} encerram o óbvio; \texttt{contradiction} e \texttt{by\_contra} exploram o absurdo.

% Módulo 02d — Trabalhos relacionados (resumos como fonte, R773+)
Com teoria e ferramenta postas, situamos o trabalho na literatura --- declarando o que foi apenas lido.
\section{Trabalhos relacionados}
\textit{Calibração por temperatura.} A literatura registra a temperatura como parâmetro de nitidez e calibração: relações entre focal loss, temperature scaling e Scorings próprios, com o hiato de generalização gerando sobreconfiança no teste \cite{komisarenko2024}; aspectos práticos e teóricos em classificação e robustez adversarial \cite{xuan2025}; e escalonamento adaptativo para modelos de linguagem \cite{xie2024}. Nossa T10 contribui a forma elementar verificada por máquina, sem limites assintóticos. \textit{Ponto flutuante.} Aproximação correta de IEEE-754 para verificação de programas, com arredondamentos, infinitos, NaNs e zeros sinalizados \cite{bagnara2019}; certificação de robustez via Lipschitz sob execução float, com lacunas float-relevantes para garantias \cite{murray2026}. Contribuímos T11 e teste empírico de propriedades, sem reivindicar certificação total. \textit{Benchmarks de agentes.} Avaliação multidimensional de LLMs como agentes em ambientes interativos \cite{agentbench2023}; lacuna de perspectiva em prompting de orquestração multiagente \cite{perspectivegap2026}; modelagem de recompensa para orquestradores \cite{orchrm2026}; arcabouço com coordenação auto-emergente inspirada em atenção \cite{orchestrator2025} --- convergência independente digna de nota, não validação. Nenhum foi executado aqui. (Síntese limitada aos resumos; leitura integral futura.)
\section{Fichamento por obra (o que cada uma informa)}
\textit{Calibração:} relações entre perda focal, scaling e propriedade \cite{komisarenko2024} informam T10; prática e teoria em classificação e robustez \cite{xuan2025} informam a leitura da Tabela~\ref{tab:temp}; escalonamento adaptativo para linguagem \cite{xie2024} informa limites do que não fizemos. \textit{Floats:} aproximação correta de IEEE-754 \cite{bagnara2019} informa o desenho do fuzz; certificação via Lipschitz sob floats \cite{murray2026} informa o que faltaria para certificar. \textit{Benchmarks:} avaliação multidimensional \cite{agentbench2023}, lacuna de perspectiva \cite{perspectivegap2026}, recompensa de orquestração \cite{orchrm2026} e coordenação atencional \cite{orchestrator2025} informam o protocolo externo pendente. \textit{Verificação:} base axiomática \cite{hoare1969}, Lean \cite{moura2021}, Mathlib \cite{mathlib2020}, Z3 \cite{demoura2008} e SymPy \cite{meurer2017} informam o núcleo. \textit{Método:} experimentação \cite{basili1986,wohlin2012}, casos \cite{runeson2009}, relato \cite{page2021}, valores-\textit{p} \cite{wasserstein2016}, estimação \cite{cumming2014}, bootstrap \cite{efron1979}, reprodutibilidade \cite{osc2015,collberg2016,ioannidis2005}, dados e modelos \cite{gebru2021,mitchell2019}, FAIR \cite{wilkinson2016} informam desenho, relato e trilha.
\par Leitura por obra: calibração informa T10; ponto flutuante informa T11 e o fuzz; benchmarks informam o protocolo externo pendente; verificação informa o núcleo Lean; método informa desenho, relato e estatística; reprodutibilidade informa a trilha; dados e modelos documentados informam manifestos. Nenhuma obra foi executada: são referências de contexto, não resultados reivindicados.
\par Comentário teórico: cada família acima é relacionada na Tabela~\ref{tab:qkv} e no texto a uma lacuna do trabalho (temperatura a T10, ponto flutuante a T11 e ao fuzz, benchmarks ao protocolo externo pendente); registra-se que esses trabalhos foram apenas lidos em resumo, sem execução, e não integram os resultados reivindicados.
\section{Leitura por obra (resumos como fonte)}
\textit{Atenção e linguagem.} \cite{vaswani2017} informa a forma herdada; \cite{luong2015}, que a forma admite variantes; \cite{devlin2019}, o alcance da arquitetura, não usado aqui; \cite{dai2019}, limites do que não foi feito; \cite{wolf2020}, o ferramental consolidado do campo. \textit{Benchmarks e avaliação.} \cite{wang2018}, \cite{rajpurkar2016} e \cite{deng2009} fixaram o padrão de avaliação por tarefas aqui adaptado em escala sintética; \cite{agentbench2023}, \cite{perspectivegap2026}, \cite{orchrm2026} e \cite{orchestrator2025} definem o palco externo ainda não pisado. \textit{Calibração e ponto flutuante.} \cite{komisarenko2024}, \cite{xuan2025} e \cite{xie2024} informam T10 e sua fronteira; \cite{bagnara2019} e \cite{murray2026} informam T11, o fuzz e o que faltaria certificar; fundamentos em \cite{muller2018,goldberg1991,harris2020,virtanen2020,meurer2017,demoura2008}. \textit{Método e auditoria.} Desenho e relato em \cite{basili1986,wohlin2012,runeson2009,page2021}; relato por intervalos em \cite{wasserstein2016,cumming2014,efron1979,gneiting2007}; trilha em \cite{osc2015,collberg2016,ioannidis2005,gebru2021,mitchell2019,wilkinson2016}; auditoria em \cite{fagan1976,mccabe1976,chidamber1994,parnas1986}; linhagem da orquestração em \cite{smith1980,erman1980}; verificação dedutiva em \cite{hoare1969,moura2021,mathlib2020}. Leitura limitada aos resumos; nenhuma obra foi executada.
\section{Síntese: matriz de cobertura deste trabalho}
Numa matriz formal $\times$ empírico $\times$ auditável: formal coberto para o núcleo (11/11 verificados); empírico coberto internamente (bancada, ablações, fuzz, correlações); auditável coberto (axiomas, alegações, referências, pinos, sincronia); externo descoberto (sem head-to-head, sem reprodução independente, sem Q/K/V formal completo). A novidade reivindicável é composicional --- verificação, avaliação e auditoria amarradas sobre um mesmo objeto com trilha local de reprodução --- não a de cada peça isolada.

% Revisão metodológica: artefatos históricos, sem reexecução experimental.
\chapter{Método}\label{cap:metodo}
A pesquisa combina formalização matemática, avaliação em simulador e inspeção documental dos artefatos locais. As etapas e seus limites são descritos a seguir para permitir a reprodução e distinguir os registros históricos da revisão editorial.
\section{Provas formais}
O arquivo \texttt{roteamento\_atencao.lean} (Apêndice~\ref{ap:lean}) foi processado por \texttt{lake env lean} no ambiente registrado, com Lean v4.34.1 e Mathlib rev.\ v4.34.1. A auditoria \texttt{\#print axioms} registra as dependências de cada enunciado. O verificador e a interface MCP incluem controles para \texttt{sorry}, \texttt{admit} e axiomas não declarados. As propriedades formalizadas referem-se a funções sobre os reais; não foi demonstrado um refinamento completo entre o modelo abstrato e todos os componentes da implementação Python.
\section{Bancada sintética e conjuntos de instâncias}
A bancada principal utiliza cinco sementes (11, 22, 33, 44 e 55), cada uma gerando uma piscina de 40 agentes sintéticos e 40 tarefas, totalizando 200 decisões. A qualidade latente $q$ não é fornecida diretamente ao roteador, mas a confiança é construída como uma proxy informativa dessa qualidade. O oráculo atribui utilidade $\mathrm{cobertura}\times((1-q_w)+q_wq)$. Comparam-se o roteador, a seleção aleatória, a escolha de menor carga e a escolha de maior cobertura. Registram-se arrependimento, acerto, latência e desvio da soma dos pesos.
As ablações utilizam as sementes efetivas 10011, 10022, 10033, 10044 e 10055, produzindo outro conjunto de 200 decisões. Incluem média igualitária das cabeças, remoção da confiança e confiança embaralhada, anticorrelacionada ou ruidosa. O pareamento compara estratégias dentro de cada execução das ablações; não vincula essas decisões às instâncias da bancada principal. A coleta de correlações utiliza sua própria sequência de sorteios e também não corresponde integralmente às tarefas das ablações.
Na condição de confiança ruidosa, adiciona-se ruído gaussiano com desvio 0,4, limita-se o resultado a $[0,1]$ e aplica-se arredondamento inteiro, gerando confiança 0 ou 1. O tratamento combina perturbação e discretização binária.
\section{Parâmetro do oráculo}
Nas ablações, avaliam-se $q_w\in\{0{,}3;0{,}6;0{,}9\}$ sobre as mesmas instâncias. O filtro de elegibilidade exige cobertura completa; assim, o oráculo entre candidatos elegíveis reduz-se a $1-q_w+q_wq$. Para $q_w>0$, seu ranking permanece inalterado e o arrependimento é $q_w(q_{\max}-q_{\mathrm{escolhido}})$. A variação descreve a escala da perda, sem testar critérios ou ordenações diferentes do oráculo.
\section{Testes numéricos e estatística exploratória}
O relatório interno registra testes aleatórios de 20\,000 vetores, de tamanho 1--50 e escores de magnitude até $\pm1000$, examinando propriedades da implementação em ponto flutuante. A ausência de violações na amostra não certifica todas as entradas nem substitui a prova de refinamento numérico. A documentação da segunda execução mencionada na versão anterior não foi identificada de forma individualizável nesta revisão; o total aqui relatado restringe-se ao registro local identificado.
Os intervalos foram calculados por aproximação normal com desvio de divisor $n$. As 40 decisões de cada semente compartilham uma piscina de agentes. Como os intervalos não consideram esse agrupamento em cinco piscinas, são resumos exploratórios que podem subestimar a incerteza entre piscinas. Uma análise com agrupamento permanece pendente. Correlações de Pearson descrevem associações lineares; não estabelecem causalidade.
\section{Ameaças à validade e considerações éticas}
\textit{Construção:} o oráculo é arbitrário e a confiança fornece informação sobre a qualidade latente por desenho. A correlação registrada de aproximadamente $+0{,}85$ explicita essa proxy privilegiada. \textit{Validade interna:} o pareamento controla diferenças de dificuldade entre estratégias de uma mesma execução, mas não elimina todos os fatores de confusão. \textit{Validade externa:} os resultados são restritos ao simulador; não há avaliação executada em benchmarks externos nem mensuração de tarefas realizadas por agentes. \textit{Validade de conclusão:} cinco piscinas, comparadores simples e intervalos sem agrupamento limitam a interpretação inferencial.
Não foram utilizados dados de participantes humanos nos experimentos descritos. A declaração ética exigida pelo programa e eventuais dispensas institucionais permanecem sob responsabilidade do autor. O estudo não fundamenta alegações de superioridade externa ou de originalidade validada por avaliadores externos.
\section{Protocolo operacional de reprodução}
O protocolo abaixo permite a terceiro reexecutar cada artefato a partir do ambiente especificado. Os resultados esperados repetem os registros preservados nos apêndices; a execução integral da fonte atual não foi revalidada nesta revisão editorial, de modo que divergências devem ser reportadas contra os registros, não contra o texto.
\subsection{Ambiente}
Lean 4 v4.34.1, Lake 5.0.0, elan 4.2.4; Mathlib rev.\ v4.34.1 com 4760 oleans em cache; Python com as dependências do repositório; \texttt{pytest} com 9 suítes. A impressão do ambiente está transcrita nos apêndices. Sem o compilador no \texttt{PATH}, o verificador degrada corretamente para \texttt{syntax\_verified} em vez de alegar compilação.
\subsection{Verificação dos onze lemas}
Comando: \nolinkurl{lake env lean RoteamentoAtencao.lean} no projeto lake pinado (saída 0 registrada). Esperado: onze vereditos \texttt{machine\_checked} e auditoria \texttt{\#print axioms} restrita a \textit{propext}, \textit{Classical.choice} e \textit{Quot.sound} por teorema (Tabela~\ref{tab:ax}). Falhas conhecidas: cache frio de oleans pode exceder o limite interno de 300\,s do teste automatizado (observado uma vez; a repetição com cache aquecido concluiu em 224\,s) --- trata-se de custo de elaboração, não de erro de prova; qualquer \texttt{sorryAx} na auditoria reprova o ciclo.
\subsection{Bancada principal e ablações}
Bancada: sementes 11, 22, 33, 44 e 55, 40 tarefas por semente; esperado: arrependimento $0{,}042\pm0{,}011$, acerto $0{,}625\pm0{,}067$, latência $2{,}17$\,ms, soma softmax 1 até $1{,}1\times10^{-16}$ (\nolinkurl{resultado_bancada.json}, Apêndice~\ref{ap:json}). Ablações: sementes efetivas 10011--10055; esperado: valores da Tabela~\ref{tab:bench} por $q_w$, diferenças pareadas e taxas de vitória/empate da Tabela~\ref{tab:pareado}. Correlações: sequência própria de sorteios; esperado: confiança $\times$ $q$ $+0{,}85$, demais coeficientes da Seção~\ref{cap:resultados}.
\subsection{Testes numéricos e suíte}
Fuzz: sementes 7 e 99, vetores de tamanho 1--50 e magnitudes até $\pm1000$; esperado: zero violações acima de $10^{-9}$ (soma) e $10^{-12}$ (deslocamento), pior desvio $\approx2{,}2\times10^{-16}$, com transcrições preservadas. Suíte: \texttt{pytest}, 9 suítes e 63 testes aprovados no registro histórico; a revisão editorial não reexecutou essa bateria.
\subsection{Pinos, hashes e sincronia}
Pinos: \texttt{git ls-remote} nos seis repositórios com SHA-256 das licenças (Tabela~\ref{tab:pins}, verificação datada de 2026-10-09). Sincronia: \nolinkurl{scanners/diagnostic_pipeline} (SPEC-935-R767, 8 testes herméticos) amarra teoremas, testes, SPEC e axiomas; o diagnóstico preliminar da revisão apresentou 18 verificações aprovadas e duas advertências, sem falhas.
\subsection{Roteiro de inspeção das checagens}
Listas validadas, extração única, faltante/excedente/proibido com detalhe, referência cruzada, leitura que devolve \texttt{None} e veredito somado por \texttt{id}. A Seção~\ref{cap:arguicoes} percorre funções e dicionário pergunta a pergunta.
\subsection{Roteiro de inspeção de gaps e gates}
Cinco motivos fechados com agregação por \texttt{id} e plano validado antes; seis motivos de gate nominais com duplicata que elimina ambas e cópia normalizada. A Seção~\ref{cap:arguicoes} percorre motivos e dicionário pergunta a pergunta.
\subsection{Roteiro de inspeção de r212 e da partição}
Dezenove testes em seis classes (gate, confiança, normalização, determinismo, carga, explicação); partição com classificação sem quebra, duplicata que elimina ambas e vazio válido. A Seção~\ref{cap:arguicoes} percorre classes, partição e dicionário pergunta a pergunta.
\subsection{Roteiro de inspeção de \texttt{explain} em uso}
Chamadas pela suíte com teste dedicado; bancada com \texttt{route} e campos diretos por economia de formato; avaliação única estrutural; eco inseparável; função pura sem estado. A Seção~\ref{cap:arguicoes} percorre roteiros, uso e dicionário pergunta a pergunta.
\subsection{Roteiro de inspeção de \texttt{route} e \texttt{explain}}
Assinatura com índice descartado por compatibilidade, avaliação única delegada, eco de entrada na explicação, determinismo triplo e docstrings como especificação. A Seção~\ref{cap:arguicoes} percorre método, dupla e dicionário pergunta a pergunta.
\subsection{Roteiro de inspeção de verbos e números}
Oito verbos da tese com endereço por pilar; dez números com chave no apêndice (arrependimento, acerto, proxy, vetores, desvios, teoremas, desenho, DOIs, tempos); dicionário de rastro, lastro, conta, retrato e guarda. A Seção~\ref{cap:arguicoes} percorre verbos, números e dicionário pergunta a pergunta.
\subsection{Roteiro de inspeção de perfil, estatísticas e serialização}
Perfil por agente com média ou zero, estatísticas unindo cadastro e motor, serialização com chave composta sem o vetor; tese metodológica lida como índice de seções. A Seção~\ref{cap:arguicoes} percorre funções, tese e dicionário pergunta a pergunta.
\subsection{Roteiro de inspeção do abstract e da observabilidade}
Conferência dígito a dígito entre resumo e abstract, ressalvas nos dois idiomas e palavras-chave termo a termo; estatísticas de motor e handbook por cópia, sonda única, cache com despejo declarado e segurança por imutabilidade. A Seção~\ref{cap:arguicoes} percorre gênero, operação e dicionário pergunta a pergunta.
\subsection{Roteiro de inspeção do motor e das funções restantes}
Hash por texto normalizado, cosseno em $[0,1]$, caminhos LiteRT/Colibri com disciplina de dimensão e recuo, estatísticas por cópia, sonda única e cache aproximado; perfil, estatísticas, serialização e cobertura no handbook. A Seção~\ref{cap:arguicoes} percorre motor, funções e dicionário pergunta a pergunta.
\subsection{Roteiro de inspeção da episteme e do handbook}
Seis regimes com léxico bilíngue, mínimo de dois sinais, dominante/secundária, confiança por contagem, afinidade com neutro e peso $0{,}20$; inferência sobre metadados com fail-open declarado. No handbook: idempotência, enriquecimento de tags, médias móveis, filtros, composto e nove campos por resposta. A Seção~\ref{cap:arguicoes} percorre episteme, handbook e dicionário pergunta a pergunta.
\subsection{Roteiro de inspeção do handbook e da conclusão}
No handbook: registro idempotente com enriquecimento de tags, episteme com override e fail-open declarado, atualização por médias móveis, casamento com filtros, composto $0{,}6/0{,}4$ e peso brando epistêmico. Na conclusão: dois níveis, cinco contribuições com endereço, assimetria $3+1$ dos objetivos e tabela de futuros como contrato. A Seção~\ref{cap:arguicoes} percorre conclusão, handbook e dicionário pergunta a pergunta.
\subsection{Roteiro de inspeção do casamento semântico}
Dimensões distintas por módulo (64 legado, 128 semântico); hierarquia LiteRT $\to$ Colibri $\to$ hash com sonda e registro; cache de 500 com estatísticas; handbook com confiança inicial $0{,}5$; casamento por \texttt{top\_k} (20/50) com agregações distintas; singleton na importação com carga preguiçosa. A Seção~\ref{cap:arguicoes} percorre resumo, casamento e dicionário pergunta a pergunta.
\subsection{Roteiro de inspeção do incorporador e da reverificação}
No incorporador (93 linhas, dimensão 64): hash assinado por token, tokenização com exemplo, normalização L2, senoidal de amplitude $0{,}1$, tarefa com requisitos em dobro e cartão com capacidades em triplo. Na reverificação Lean 4: suíte de 4 testes sobre o arquivo de 150 linhas no projeto pinado (4 aprovados em $191$\,s aquecido), auditoria de axiomas e registro do comportamento ante cache frio. A Seção~\ref{cap:arguicoes} percorre sementes, incorporador e reverificação pergunta a pergunta.
\subsection{Roteiro de inspeção das cabeças}
Semântica (casamento \texttt{top\_k=20} com recuo legado, cosseno com neutro $0{,}5$); cobertura (fração, $1{,}0$ sem requisitos); confiança (média per-skill \texttt{top\_k=50} com recuo ao cartão/razão); carga ($1-\mathrm{carga}$, ausência livre); sanitizador \texttt{\_unit\_score} (booleanos, não-numéricos, não-finitos e fora de faixa tratados). A Seção~\ref{cap:arguicoes} percorre cada cabeça e cada condição de ablação pergunta a pergunta.
\subsection{Roteiro de inspeção dos testes e transcrições}
Nos testes do núcleo (107 linhas): registro da SPEC, existência com frontmatter e bloqueadores, par ruim/bom fail-closed, exportação com SHA recomputado, autoformalização de formato, manifesto com contagens e anti-overclaim condicional. Nas transcrições (159 linhas): índice nominativo, verboso com fecho de 63 em $248$\,s, ambiente com versões e hashes, e as duas sementes do fuzz com zero violações e piores desvios. A Seção~\ref{cap:arguicoes} percorre cada teste e cada linha pergunta a pergunta.
\subsection{Roteiro de inspeção do estudo de caso, da MCP e dos testes}
No estudo de caso (86 linhas): lista \texttt{SEL} com doze nomes e onze cartões carregados (\texttt{technical-writer} sem correspondência), extração por regex de nome/tags/descrição, carga e confiança fixas, cinco tarefas, ranking arredondado e impressão resumida. Na MCP: bloqueio por regex a \texttt{sorry}/\texttt{admit} e \texttt{axiom}, rebaixamento anti-overclaim sem compilador, exportação com SHA-256 e procedência por chamada. Nos testes: oito verificações do núcleo mais oito do pipeline, verdes em $15{,}23$\,s; bateria histórica com ressalva. A Seção~\ref{cap:arguicoes} percorre cada peça pergunta a pergunta.
\subsection{Roteiro de inspeção do roteador e do pipeline}
No roteador (349 linhas): motivos de gate (\texttt{\_hard\_gate\_reasons}), partição com regra de duplicata, caminho vazio, índice posicional fixo e descartado, utilidade por \texttt{fsum}, softmax, ordenação por ($-\mathrm{peso}$, $\mathrm{id}$) e \texttt{explain} completo; validador de pesos e cabeças com casamento e recuo legado. Na MCP: quatro funções mais procedência, em duas camadas fail-closed. No pipeline: \texttt{varrer} com leitura, extração de nomes e checagem de nomes e cobertura (8 testes herméticos). A Seção~\ref{cap:arguicoes} percorre cada peça pergunta a pergunta.
\subsection{Roteiro de inspeção dos scripts Python}
\texttt{bancada\_routing.py} (142 linhas): vocabulário de 12 capacidades e 6 palavras de ruído; \texttt{make\_pool} (40 agentes, 3--6 capacidades, $q$ latente separado, carga e confiança ruidosa); \texttt{make\_task} (2--3 requisitos, $\geq4$ elegíveis em até 50 tentativas); \texttt{oracle\_u} (cobertura vezes $0{,}4+0{,}6q$); \texttt{ci95} (média e $1{,}96s_n/\sqrt{n}$ com \texttt{pstdev}); \texttt{run\_seed} (escolhas, arrependimentos, acertos, latência, soma) e agregado com impressão. \texttt{bancada\_ablacao.py} (122 linhas): reaproveita piscina, tarefa e intervalo; \texttt{PESOS}, desempate por $(-\mathrm{peso}, \mathrm{id})$, \texttt{oracle\_qw}, semente deslocada, três ataques à confiança e dois comparadores fortes, diferenças pareadas e taxas, arquivo por tag de $q_w$. \texttt{correlacoes.py} (93 linhas): Pearson com divisor $n$ e \texttt{None} ante variância zero. \texttt{contraprova\_float.py} (63 linhas): tamanhos, magnitudes, tolerâncias, seis contadores e saída com erro ante violação. A Seção~\ref{cap:arguicoes} percorre cada função pergunta a pergunta.
\subsection{Roteiro de inspeção do código Lean}
O arquivo tem 150 linhas: cabeçalho com a SPEC e a definição de eficiência (otimalidade relativa à utilidade definida, nunca superioridade universal), \texttt{import Mathlib}, abertura de \texttt{BigOperators Finset} e onze teoremas na ordem T1--T11 com nomes estáveis (\texttt{softmax\_soma\_um} a \texttt{sem\_overflow\_apos\_max}). O revisor confere: as hipóteses de não-vazio (\texttt{Nonempty}) e não-trivialidade (\texttt{Nontrivial}) onde o texto as exige (T1, T4, T5, T7, T10); a hipótese de maximizador estrito em T10 (\texttt{hmax}) e a de unicidade em T7 (\texttt{h}); as táticas-chave por teorema (\texttt{sum\_div}/\texttt{div\_self}, \texttt{sum\_le\_sum}/\texttt{sum\_mul}, \texttt{exists\_max\_image}, \texttt{exp\_sub}/\texttt{field\_simp}, \texttt{add\_sum\_erase}, \texttt{sum\_eq\_single\_of\_mem}, \texttt{div\_lt\_div\_iff₀}, \texttt{pairwise\_append}, \texttt{sum\_lt\_sum} com testemunha, \texttt{exp\_le\_one\_iff}); e a ausência de \texttt{sorry}, \texttt{admit} e axiomas declarados, complementada pela auditoria. A Seção~\ref{cap:arguicoes} percorre essa anatomia pergunta a pergunta.
\subsection{Roteiro de inspeção manual dos JSONs}
Sem executar nada, o revisor confere: em \nolinkurl{resultado_bancada.json}, as chaves por estratégia (\texttt{regret\_medio}, \texttt{regret\_ci95}, \texttt{hit\_rate}, \texttt{hit\_ci95}), a latência (\texttt{media}, \texttt{max}), o desvio máximo da soma e as três decisões da \texttt{amostra}; em cada \nolinkurl{resultado_ablacao_q*.json}, \texttt{qw}, os arrependimentos por estratégia e, por perturbação, \texttt{diff\_pareado\_medio}, \texttt{diff\_pareado\_ci95} e \texttt{atencao\_vence\_ou\_empata}; em \nolinkurl{correlacoes.json}, os cinco coeficientes, os nulos com sua \texttt{nota} de variância zero e \texttt{n\_decisoes}; em \nolinkurl{per_seed.json}, os cinco blocos com arrependimentos e acerto; em \nolinkurl{estudo_caso_reais.json}, \texttt{n\_agentes} e os campos por tarefa (\texttt{elegiveis}, \texttt{ranking}, \texttt{soma\_pesos}, \texttt{lat\_ms}, \texttt{cabecas\_top1}). Cada número das tabelas do texto deve ser rastreável a uma dessas chaves; a Seção~\ref{cap:arguicoes} percorre essa rastreabilidade item a item.
\subsection{Desenvolvimento (registros históricos)}
Os registros de desenvolvimento vinculam as funcionalidades às especificações SPEC-935 e às transcrições de testes preservadas nos apêndices. O protocolo documenta etapas RED--GREEN, auditorias de axiomas e correções de falhas. Esses registros históricos não equivalem à reexecução integral na data desta revisão. As versões, os hashes e as saídas devem acompanhar o depósito final dos artefatos, ainda pendente.

\section{Interpretação das métricas}
O arrependimento é a diferença entre a maior utilidade do oráculo e a utilidade do candidato escolhido; zero indica uma escolha ótima segundo esse oráculo. A taxa de acerto é a fração de decisões que atingem esse máximo. A taxa de vitória ou empate nas ablações é a fração de decisões em que a atenção apresenta arrependimento menor ou igual ao comparador. Pearson $r$ mede associação linear, sem demonstrar causalidade ou independência quando próximo de zero. A semente fixa a sequência pseudoaleatória utilizada em cada coleta.
\section{Intervalos registrados e seus limites}
O relato por estimação é apoiado por \citet{cumming2014,wasserstein2016}. Nos scripts deste trabalho, a margem registrada é $1{,}96s_n/\sqrt{n}$, com $s_n$ calculado com divisor $n$. Com $n=200$ e margem $0{,}0105$, obtém-se $s_n\approx0{,}0758$. Essa identidade explica o cálculo preservado, sem validar a independência das decisões. A interpretação comparativa requer considerar o desenho, a dependência entre observações e o intervalo da diferença pareada. A não sobreposição de intervalos isolados não substitui essa análise.
Como cada piscina gera 40 decisões, a análise futura deve considerar o agrupamento por semente. Reamostragem \cite{efron1979} é uma possibilidade, desde que respeite a unidade de agrupamento; cinco piscinas ainda limitam a precisão. Não foi efetuado recálculo inferencial nesta revisão.
\section{Pareamento e calibração probabilística}
Na ablação com $q_w=0{,}6$, a diferença de arrependimento entre média igualitária e atenção é aproximadamente $+0{,}035$, com margem exploratória $0{,}011$. O sinal positivo descreve vantagem na amostra pareada registrada; não demonstra robustez fora desse desenho. A normalização de pesos não é calibração probabilística. Regras de escore próprias \cite{gneiting2007} oferecem fundamentos para outra análise, mas não foram calculadas curvas de calibração, erros de calibração ou escores probabilísticos nesta avaliação.

% Módulo 04 — Provas formais detalhadas (R762+; 11 teoremas)
\chapter{Provas formais}\label{cap:provas}
A formalização abrange o núcleo matemático abstrato. Não foi demonstrado um refinamento completo entre esses objetos e a implementação Python, incluindo elegibilidade, representações semânticas, desempate, arredondamento e execução dos agentes. Os registros de auditoria dos onze enunciados são apresentados na Tabela~\ref{tab:ax}.
Os enunciados e suas hipóteses são resumidos a seguir: normalização e limitação convexa são o conteúdo formal da exigência de soma unitária dos pesos, prática padrão em ponderação exponencial desde \cite{vaswani2017} e objeto de calibração na literatura \cite{komisarenko2024,xuan2025}; maximizador finito e seleção gulosa formalizam a regra de decisão; invariância, pertinência, singleton, monotonicidade, top-\textit{k}, nitidez e overflow cobrem estabilidade, casos-limite e decisão, cada qual demonstrado em Lean sobre Mathlib \cite{moura2021,mathlib2020}:
\begin{equation*}
\sum_i\left(e^{f_i}/\sum_j e^{f_j}\right)=1 \tag{T1}
\end{equation*}
\begin{equation*}
m\leq\sum_{i\in s}w_ix_i\leq M,\quad \sum_{i\in s}w_i=1,\ w_i\geq0,\ m\leq x_i\leq M \tag{T2/T6}
\end{equation*}
\begin{equation*}
\exists\,\text{melhor}\in C,\ \forall a\in C,\ u(a)\leq u(\text{melhor}) \tag{T3}
\end{equation*}
\begin{equation*}
e^{f_i-c}/\sum_j e^{f_j-c}=e^{f_i}/\sum_j e^{f_j} \tag{T4}
\end{equation*}
\begin{equation*}
0<e^{f_i}/\sum_j e^{f_j}<1\ \text{(2+ candidatos)};\quad =1\ \text{(singleton)} \tag{T5/T7}
\end{equation*}
\begin{equation*}
x\mapsto e^x/(e^x+C)\ \text{estritamente crescente para }C>0 \tag{T8}
\end{equation*}
\begin{equation*}
l\ \text{pairwise-decrescente}\Rightarrow\forall x\in\text{drop}_k,\forall y\in\text{take}_k,\ y\geq x \tag{T9}
\end{equation*}
\begin{equation*}
0<\tau_1<\tau_2\Rightarrow w_{i^*}(\tau_1)>w_{i^*}(\tau_2)\ \text{(maximizador único; ao menos dois candidatos)} \tag{T10}
\end{equation*}
\begin{equation*}
f_i\leq M\Rightarrow e^{f_i-M}\leq1 \tag{T11}
\end{equation*}
Ingredientes de prova, por teorema: (T1) $\sum_i(e^{f_i}/\sum_j e^{f_j})=1$ em tipo finito não vazio (positividade da exponencial, \texttt{Finset.sum\_pos} sobre o universo, reescrita \texttt{Finset.sum\_div}, \texttt{div\_self}). (T2/T6) $m\leq\sum w_ix_i\leq M$ sob pesos convexos (monotonicidade da multiplicação por não negativos, reescrita \texttt{Finset.sum\_mul}, hipótese de soma 1). (T3) $\exists$ maximizador finito (máximo de imagem, \nolinkurl{Finset.exists_max_image}). (T4) invariância a $f_i-c$ (subtração da exponencial, soma de quocientes, simplificação de corpos que fecha sozinha). (T5) pesos em $(0,1)$ sob não-trivialidade (divisão, soma do resto positivo via \texttt{add\_sum\_erase}, linear). (T7) peso 1 no singleton (soma de um termo, \texttt{sum\_eq\_single\_of\_mem} com elemento explícito --- cuja ausência gerou \texttt{sorryAx} na primeira versão, capturado pela auditoria). (T8) monotonicidade estrita no próprio escore (divisão de desigualdades na forma atual \texttt{div\_lt\_div\_iff₀}, linear; nome consultado no Mathlib pinado, não na memória). (T9) take-\textit{k} domina o resto sob pairwise (concatenação take/drop, \texttt{pairwise\_append}). (T10) nitidez por temperatura: com maximizador único, esfriar aumenta seu peso (fatoração de somas, divisão de desigualdades, soma estrita com testemunha, exigindo não-trivialidade --- sem a qual a tese é falsa). (T11) limite superior da exponencial após deslocamento: expoentes deslocados $\leq 0$ (caracterização da exponencial real). Auditoria: todos com axiomas restritos a \textit{propext}, \textit{Classical.choice}, \textit{Quot.sound} (Tabela~\ref{tab:ax}), sem \texttt{sorryAx}.
\begin{table}[htbp]
\fontsize{10}{12}\selectfont\SingleSpacing
\centering
\caption{Auditoria de axiomas (\texttt{\#print axioms}).}
\label{tab:ax}
\begin{tabularx}{\textwidth}{XX}
\toprule
Teorema & Axiomas \\
\midrule
T1 \texttt{softmax\_soma\_um} & propext, Classical.choice, Quot.sound \\
T2 \nolinkurl{combinacao_convexa_limitada} & propext, Classical.choice, Quot.sound \\
T3 \nolinkurl{roteamento_guloso_otimo} & propext, Classical.choice, Quot.sound \\
T4 \texttt{softmax\_invariante\_max} & propext, Classical.choice, Quot.sound \\
T5 \texttt{softmax\_mem\_Ioo} & propext, Classical.choice, Quot.sound \\
T6 \nolinkurl{combinacao_convexa_limitada_baixo} & propext, Classical.choice, Quot.sound \\
T7 \nolinkurl{softmax_peso_total_singleton} & propext, Classical.choice, Quot.sound \\
T8 \nolinkurl{peso_cresce_com_escore} & propext, Classical.choice, Quot.sound \\
T9 \texttt{topk\_dominam\_resto} & propext, Classical.choice, Quot.sound \\
T10 \nolinkurl{temperatura_aguca_maximo} & propext, Classical.choice, Quot.sound \\
T11 \texttt{sem\_overflow\_apos\_max} & propext, Classical.choice, Quot.sound \\
\bottomrule
\end{tabularx}
\fonte{Elaboração própria (2026), a partir dos artefatos e das fontes identificados no texto.}
\end{table}

% Módulo 04b — Demonstrações integrais por teorema (R775+; tratado R776+)
\section{Demonstrações integrais, teorema por teorema}
Cada subseção traz enunciado, derivação completa, instância numérica real e necessidade das hipóteses; os nomes de lemas entre parênteses são os usados no código Lean do Apêndice~\ref{ap:lean}.
\subsection{T1 --- soma unitária}
Enunciado: $\sum_i(e^{f_i}/\sum_j e^{f_j})=1$ em tipo finito não vazio. Derivação: cada termo exponencial é positivo (\texttt{exp\_pos}), logo a soma $S$ sobre o universo não vazio é positiva (\texttt{sum\_pos}) e não nula; reescreva a soma de quocientes como quociente das somas (\texttt{sum\_div}); obtém-se $S/S=1$ (\texttt{div\_self}). Instância: $f=[1,2,3]$ dá pesos $[0{,}0900;\ 0{,}2447;\ 0{,}6652]$ somando $1{,}0$. Sem o não-vazio, o denominador seria soma vazia ($=0$) e a fração indefinida.
\subsection{T2/T6 --- limitação bilateral convexa}
Enunciados: $m\leq\sum w_ix_i\leq M$ sob $w_i\geq0$, $\sum w_i=1$, $m\leq x_i\leq M$. Derivação: termo a termo, $w_im\leq w_ix_i\leq w_iM$ (\nolinkurl{mul_le_mul_of_nonneg_left}); somar preserva cada lado (\texttt{sum\_le\_sum}); fatorar a constante (\texttt{sum\_mul}) e usar soma $=1$. Instância: $w=[0{,}3;0{,}7]$, $x=[2,8]$: $6{,}2\in[2,8]$. Contraexemplo sem a hipótese: pesos $[2,-1]$ somando 1 com valores $[0,10]$ devolvem $-10<\min$ --- por isso o validador rejeita pesos fora de $[0,1]$.
\subsection{T3 --- maximizador finito}
Enunciado: existe maximizador em conjunto finito não vazio. Derivação: invocação direta do máximo de imagem (\texttt{exists\_max\_image}); finitude impede fuga ao infinito, não-vazio garante testemunha. Sem finitude ($\mathbb{N}$ com identidade), não há máximo.
\subsection{T4 --- invariância a deslocamento}
Enunciado: subtrair $c$ nada altera. Derivação: $e^{f_j-c}=e^{f_j}/e^c$ (\texttt{exp\_sub}); soma de quocientes vira quociente das somas; simplificação de corpos fecha sozinha (\texttt{field\_simp}), pois $e^c\neq0$. Instância: $[1,2,3]$ e $[-4,-3,-2]$ produzem os mesmos pesos em aritmética real; os valores arredondados do exemplo coincidem.
\subsection{T5 --- pertinência a $(0,1)$}
Enunciado: cada peso em $(0,1)$ com ao menos dois candidatos. Derivação: positividade por divisão de positivos; escreva $S=e^{f_i}+R$ com $R>0$ (resto via \texttt{add\_sum\_erase}, não vazio por testemunha da não-trivialidade); então $e^{f_i}/S<1$ por divisão e linearidade. Instância: os três pesos de $[1,2,3]$ estão em $(0,1)$. Sem dois candidatos, o singleton dá peso 1 --- caso de T7.
\subsection{T7 --- peso unitário no singleton}
Enunciado: candidato único recebe peso 1. Derivação: a soma colapsa no único termo (\texttt{sum\_eq\_single\_of\_mem} com elemento explícito --- cuja ausência gerou \texttt{sorryAx}, capturado pela auditoria) e $e/e=1$. Instância: $[5]\to[1{,}0]$, também observada no roteador real.
\subsection{T8 --- monotonicidade estrita}
Enunciado: $x\mapsto e^x/(e^x+C)$ estritamente crescente para $C>0$. Derivação: reduza a produtos cruzados (divisão de desigualdades na forma vigente, consultada no fonte pinado) e conclua por linearidade com $e^aC<e^bC$. Instância: $C=1$ dá $0{,}7311<0{,}8808$. Sem $C>0$, a fração degenera.
\subsection{T9 --- dominância top-\textit{k}}
Enunciado: numa lista pairwise-decrescente, o prefixo domina o resto. Derivação: reescreva como take $++$ drop, decomponha a ordem da concatenação e aplique a componente cruzada. É a anatomia do padrão ordenar-e-fatiar, independente do algoritmo de ordenação.
\subsection{T10 --- nitidez por temperatura}
Enunciado: com maximizador único, esfriar aumenta seu peso. Derivação: fatore $A(\tau_1)S(\tau_2)$ contra $A(\tau_2)S(\tau_1)$; cada termo com $j\neq i^*$ é estrito pela identidade das diferenças positivas; some com testemunha estrita e conclua por divisão de desigualdades. Instância: $0{,}6652\to0{,}8668$. O enunciado formal adota maximizador único e pelo menos dois candidatos como hipóteses suficientes. Máximos empatados podem ganhar massa quando existem escores inferiores; se todos os escores forem iguais, a temperatura não altera os pesos.
\subsection{T11 --- limite superior da exponencial após deslocamento}
Enunciado: se $f_i\leq M$, então $e^{f_i-M}\leq1$ em reais exatos. Derivação: monotonicidade da exponencial e $e^0=1$. A prova não certifica subtração, exponenciação, soma ou divisão IEEE-754. Instâncias: $e^{-2}\approx0{,}1353$ e $e^{0}=1{,}0$. Sem majorante, nada limita os expoentes.
\par Mapa de dependências: T1 sustenta T5 e T7; T2/T6 sustentam a combinação das cabeças; T4 sustenta a implementação estável usada em T10; T3 sustenta a decisão por máximo ao lado de T8; T9 descreve a ordenação do top-\textit{k}, sem constituir sozinho prova do algoritmo concreto.
\section{Necessidade das hipóteses (contraexemplos)}
Reunidas: T1 exige não-vazio; T2 exige não-negatividade e soma 1; T3 exige finitude e não-vazio; T5 exige dois candidatos; T8 exige $C>0$; T10 exige maximizador único e não-trivialidade; T11 exige majorante. As hipóteses delimitam os enunciados formalizados. Os contraexemplos apresentados mostram falhas específicas quando algumas delas são removidas; não demonstram a necessidade individual de todas as hipóteses.

% Módulo 05 — Avaliação empírica e testes numéricos (R762+; números integrais)
\chapter{Avaliação empírica e testes numéricos}\label{cap:resultados}
Esta seção apresenta os resultados históricos da bancada e das ablações, distinguidos pelos conjuntos de instâncias descritos na Seção~\ref{cap:metodo}.
\section{Fuzz em ponto flutuante}
O relatório interno registra testes aleatórios de 20\,000 vetores, sem violações acima das tolerâncias examinadas: $10^{-9}$ para a soma e $10^{-12}$ para o deslocamento. O maior desvio registrado da soma foi aproximadamente $2{,}2\times10^{-16}$ e o do deslocamento foi $0{,}0$. A preservação de argmax, os limites, o caso de candidato único e o tratamento do conjunto vazio também foram examinados. A ausência de violações nessa amostra não demonstra correção para todas as entradas. Testes por entradas aleatórias têm antecedentes na avaliação de utilitários Unix \cite{miller1990}.
\section{Bancada principal e variação do parâmetro do oráculo}
Bancada principal (200 decisões): atenção $0{,}042\pm0{,}011$ (acerto $0{,}625\pm0{,}067$) contra $0{,}19$--$0{,}22$; latência $2{,}17$\,ms; soma softmax 1 até $1{,}1\times10^{-16}$ Os valores da bancada principal constam de \nolinkurl{resultado_bancada.json}, reproduzido no Apêndice~\ref{ap:json}.
A Tabela~\ref{tab:bench} apresenta as médias das ablações para cada valor do parâmetro.
\begin{table}[htbp]
\fontsize{10}{12}\selectfont\SingleSpacing
\centering
\caption{Ablações: arrependimento por parâmetro do oráculo e diferença em relação à atenção (200 decisões por configuração).}
\label{tab:bench}
\begin{tabular}{lccc}
\toprule
$q_w$ & Atenção & Pesos iguais (dif.) & Sem confiança (dif.) \\
\midrule
0,3 & $0{,}0146$ & $0{,}0319$ (+0,017) & $0{,}1072$ (+0,093) \\
0,6 & $0{,}0292$ & $0{,}0638$ (+0,035) & $0{,}2144$ (+0,185) \\
0,9 & $0{,}0438$ & $0{,}0957$ (+0,052) & $0{,}3216$ (+0,278) \\
\bottomrule
\end{tabular}
\fonte{Elaboração própria (2026), a partir dos artefatos e das fontes identificados no texto.}
\end{table}
Outras condições de ablação ($q_w=0{,}6$): confiança embaralhada $0{,}2243$ (+0,195), anticorrelacionada $0{,}3604$ (+0,331), ruidosa $0{,}1515$ (+0,122); vitória/empate da atenção $\geq 0{,}92$ em todos os pareamentos. Os registros internos relatam reexecução determinística; uma reprodução independente permanece pendente. A Tabela~\ref{tab:pareado} reúne as diferenças pareadas e as taxas de vitória ou empate.
\begin{table}[htbp]
\fontsize{10}{12}\selectfont\SingleSpacing
\centering
\caption{Diferenças pareadas por decisão vs.\ atenção ($q_w=0{,}6$) e taxa de vitória/empate.}
\label{tab:pareado}
\begin{tabular}{lcc}
\toprule
Estratégia & Diferença $\pm$ margem exploratória & Atenção vence/empata \\
\midrule
Média igualitária & $+0{,}035\pm0{,}011$ & $0{,}99$ \\
Sem confiança & $+0{,}185\pm0{,}024$ & $0{,}98$ \\
Confiança ruidosa & $+0{,}122\pm0{,}022$ & $0{,}93$ \\
Confiança embaralhada & $+0{,}195\pm0{,}025$ & $0{,}92$ \\
Confiança anticorrelacionada & $+0{,}331\pm0{,}019$ & $0{,}97$ \\
\bottomrule
\end{tabular}
\fonte{Elaboração própria (2026), a partir dos artefatos e das fontes identificados no texto.}
\end{table}
\begin{table}[htbp]
\fontsize{10}{12}\selectfont\SingleSpacing
\centering
\caption{Arrependimento por estratégia nos três valores do parâmetro do oráculo (200 decisões cada).}
\label{tab:mundos}
\begin{tabular}{lccc}
\toprule
Estratégia & $q_w=0{,}3$ & $q_w=0{,}6$ & $q_w=0{,}9$ \\
\midrule
Atenção & $0{,}0146$ & $0{,}0292$ & $0{,}0438$ \\
Média igualitária & $0{,}0319$ & $0{,}0638$ & $0{,}0957$ \\
Sem confiança & $0{,}1072$ & $0{,}2144$ & $0{,}3216$ \\
Confiança ruidosa & $0{,}0757$ & $0{,}1515$ & $0{,}2272$ \\
Confiança embaralhada & $0{,}1121$ & $0{,}2243$ & $0{,}3364$ \\
Confiança anticorrelacionada & $0{,}1802$ & $0{,}3604$ & $0{,}5406$ \\
\bottomrule
\end{tabular}
\fonte{Elaboração própria (2026), a partir dos artefatos e das fontes identificados no texto.}
\end{table}
O ranking das estratégias permanece estável na Tabela~\ref{tab:mundos} porque, entre elegíveis com cobertura unitária, o arrependimento é proporcional a $q_w$. As configurações não representam avaliações independentes de oráculos distintos. As diferenças pareadas e suas margens descrevem as instâncias das ablações, sem correção para agrupamento.
\begin{table}[htbp]
\fontsize{10}{12}\selectfont\SingleSpacing
\centering
\caption{Arrependimento por semente e estratégia (40 decisões cada; acerto da atenção entre parênteses).}
\label{tab:seeds}
\begin{tabular}{lcccc}
\toprule
Semente & Atenção & Aleatório & Menor carga & Maior cobertura \\
\midrule
11 & $0{,}0752$ (0,45) & $0{,}1889$ & $0{,}3157$ & $0{,}2319$ \\
22 & $0{,}0172$ (0,78) & $0{,}2463$ & $0{,}2276$ & $0{,}2310$ \\
33 & $0{,}0337$ (0,55) & $0{,}1719$ & $0{,}2447$ & $0{,}2033$ \\
44 & $0{,}0716$ (0,60) & $0{,}1992$ & $0{,}2038$ & $0{,}2309$ \\
55 & $0{,}0129$ (0,75) & $0{,}1350$ & $0{,}0918$ & $0{,}1618$ \\
\bottomrule
\end{tabular}
\fonte{Elaboração própria (2026), a partir dos artefatos e das fontes identificados no texto.}
\end{table}
A atenção obteve o menor arrependimento nas 5 sementes avaliadas (Tabela~\ref{tab:seeds}).
\section{Estudo de caso completo: Ana contra Cid}
Na tarefa-guia, as cabeças mediram: semântica Ana $0{,}8417$ contra Cid $0{,}7703$; capacidade $1{,}0$ para ambas (o gate já havia decidido); confiança $0{,}9$ contra $0{,}7$; carga livre $0{,}9$ contra $0{,}6$. Utilidades: Ana $0{,}30\cdot0{,}8417+0{,}35\cdot1{,}0+0{,}25\cdot0{,}9+0{,}10\cdot0{,}9=0{,}9175$; Cid $0{,}30\cdot0{,}7703+0{,}35\cdot1{,}0+0{,}25\cdot0{,}7+0{,}10\cdot0{,}6=0{,}8161$. Pesos finais $0{,}5253$ contra $0{,}4747$. Note-se a capacidade constante: com cobertura total dos elegíveis, a decisão correu nas outras três cabeças --- exatamente o fenômeno quantificado nas correlações.
\section{Correlações de Pearson (200 decisões)}
Arrependimento $\times$ elegíveis $+0{,}08$; arrependimento $\times$ dispersão $-0{,}16$; acerto $\times$ elegíveis $-0{,}19$; acerto $\times$ dispersão $+0{,}13$; confiança $\times$ $q$ $+0{,}85$ (checagem de desenho: proxy forte, fonte da circularidade parcial); cabeças quase descorrelacionadas (semântica $\times$ confiança $+0{,}05$); capacidade constante após o filtro de elegibilidade (o gate decide cobertura; a escolha corre nas outras cabeças). A transcrição histórica preservada registra 63 testes aprovados em nove suítes. A revisão editorial não reexecutou essa bateria. O diagnóstico preliminar da revisão apresentou 18 verificações aprovadas e duas advertências, sem falhas.

% Módulo 05b — Demonstração de roteamento sobre cartões do catálogo (R776+)
\section{Demonstração de roteamento sobre cartões do catálogo}
Foram utilizados onze cartões extraídos do catálogo local para cinco descrições de tarefas construídas com seu vocabulário de tags. A demonstração avalia elegibilidade e ranking; não executa os agentes nem mede o sucesso na realização das tarefas. Quatro decisões foram singleton (peso $1{,}0$ --- T7 em operação); a quinta opôs \texttt{reviewer} a \texttt{code-reviewer} em tarefa de revisão: ambas com cobertura, confiança e carga idênticas ($1{,}0$; $0{,}5$; $1{,}0$), a semântica decidiu $0{,}5589$ e o placar fechou $0{,}505$ contra $0{,}495$ --- decisão puramente semântica, à margem mínima. A primeira decisão custou $27{,}8$\,ms (carga fria do casamento semântico) e as seguintes $0{,}4$--$0{,}5$\,ms; soma dos pesos $1{,}0$ em todas. Tarefas fora do vocabulário retornam vazio por gate --- resultado compatível com o contrato do filtro; a cobertura do vocabulário do catálogo ainda requer avaliação.

\section{Interpretação das estratégias}
A estratégia de atenção utiliza quatro componentes com pesos fixos declarados, sem aprendizagem desses pesos. No simulador, a confiança é construída a partir da qualidade latente com ruído, constituindo uma proxy informativa por desenho. A média igualitária compara outra ponderação dos mesmos componentes; a remoção da confiança avalia a decisão sem essa informação.
A seleção aleatória sorteia entre elegíveis. A escolha de menor carga utiliza a disponibilidade, que não compõe diretamente a qualidade latente do oráculo. A escolha de maior cobertura encontra empate entre candidatos plenamente elegíveis, pois o filtro já exigiu cobertura completa. Seu resultado não mede discriminação adicional por cobertura.
A remoção da confiança e suas perturbações produziram maior arrependimento nas amostras das ablações. Esse resultado é compatível com a dependência da decisão em relação à proxy. Não demonstra equivalência estatística com o acaso, relação dose--resposta ou ausência de influência das sementes. A condição ruidosa inclui discretização binária, conforme descrito na Seção~\ref{cap:metodo}.

\chapter{Discussão}\label{cap:discussao}
Os resultados articulam duas formas de evidência com alcances diferentes: as provas documentam propriedades do modelo matemático; a bancada descreve o comportamento de uma implementação em instâncias sintéticas. A composição torna explícitos os pressupostos da decisão, mas não fornece uma certificação integral da execução nem uma comparação de desempenho em tarefas reais.
\section{Alcance das garantias formais}
Os onze lemas são propriedades matemáticas elementares, formalizadas sobre o objeto estudado. Não se reivindica originalidade desses fatos. T3 estabelece existência de maximizador finito, enquanto T9 descreve a dominância de um prefixo de lista já ordenada. A correspondência entre essas propriedades e o código exige demonstrar a preservação de hipóteses pelo filtro, pelos escores, pelo desempate e pela aritmética utilizada. T11 limita a exponencial deslocada em reais exatos; erros de arredondamento, underflow, entradas não finitas e operações intermediárias permanecem fora dessa prova.
\section{Dependência do desenho sintético}
A vantagem observada depende da confiança, construída como proxy da qualidade latente. Sua remoção aumenta o arrependimento no conjunto de ablações, sem estabelecer equivalência com a seleção aleatória. Mesmo confiança igual à qualidade latente não garantiria que a utilidade combinada preservasse o ranking do oráculo, pois semântica e carga também influenciam a escolha. Essa configuração não foi avaliada.
Os três valores de $q_w$ reescalam a perda sem mudar o ranking do oráculo entre elegíveis. Sua estabilidade não constitui evidência independente de robustez a outros critérios de utilidade. Os intervalos preservados ignoram a dependência dentro das piscinas; uma análise que considere agrupamento pode alterar as margens e a interpretação comparativa.
\section{Riscos residuais e avaliação interna}
\begin{table}[htbp]
\fontsize{10}{12}\selectfont\SingleSpacing\centering
\caption{Limitações identificadas e requisitos de avaliação futura.}\label{tab:riscos}
\begin{tabularx}{\textwidth}{XX}\toprule
Limitação & Requisito futuro \\\midrule
Oráculo arbitrário e proxy privilegiada & Outros critérios de qualidade e confiança sem acesso ao gerador latente \\
Cinco piscinas e intervalos sem agrupamento & Mais piscinas independentes e análise na unidade correta \\
Modelo real e implementação em ponto flutuante & Refinamento formal e tratamento de casos numéricos especiais \\
Comparadores internos & Benchmark externo executado com orçamento e critérios comuns \\
Artefatos locais & Depósito versionado e reprodução independente \\
\bottomrule\end{tabularx}
\fonte{Elaboração própria (2026), a partir dos artefatos e das fontes identificados no texto.}
\end{table}
Duas revisões internas automatizadas examinaram a bancada e os teoremas. Não constituem revisão cega ou validação externa. As objeções sobre proxy, comparadores, analogia Q/K/V e lacuna numérica são compatíveis com as limitações descritas na Tabela~\ref{tab:riscos}. As ablações acrescentam comparações entre ponderações e perturbações, mas não eliminam essas objeções.
\section{Relação com a literatura e perspectivas}
A literatura de calibração por temperatura \cite{komisarenko2024,xuan2025,xie2024} aborda objetivos que não foram medidos aqui: a normalização por softmax não demonstra calibração probabilística. Os estudos de aritmética flutuante \cite{bagnara2019,murray2026} indicam a necessidade de garantias próprias para a execução numérica. Os benchmarks de agentes \cite{agentbench2023,perspectivegap2026,orchrm2026,orchestrator2025} fornecem contextos para comparação futura, sem validar os resultados locais. A interpretação educacional das trilhas auditáveis é uma hipótese de aplicação futura, sem avaliação com estudantes nesta pesquisa.

% Módulo 06b — Arguições antecipadas e respostas (R778 aprofundado em R779; suplemento de matéria)
\chapter{Arguições antecipadas e respostas}\label{cap:arguicoes}
Este capítulo antecipa perguntas plausíveis de banca e as responde estritamente nos limites do que foi demonstrado, medido ou declarado nos capítulos anteriores. Nenhuma resposta introduz dados novos: os números citados repetem os registros da Seção~\ref{cap:resultados} e do Método (Seção~\ref{cap:metodo}); as provas citadas remetem à Seção~\ref{cap:provas}; as derivações desta seção reescrevem, passo a passo, cálculos já apresentados, sem acrescentar instâncias inéditas. Quando a pergunta aponta além do alcance, a resposta registra a lacuna e o trabalho futuro correspondente, conforme a Tabela~\ref{tab:riscos} e a Tabela~\ref{tab:futuro}. O último bloco ensina a ler cada tabela e apêndice, para que o examinador confira cada número na fonte.
\section{Provas formais}
\textit{1. O que exatamente foi provado por máquina?} Onze enunciados sobre funções reais (T1--T11), listados na Seção~\ref{cap:provas} com auditoria de axiomas na Tabela~\ref{tab:ax}: soma unitária, limitação bilateral convexa, maximizador finito, invariância a deslocamento, pertinência a $(0,1)$, caso singleton, monotonicidade, dominância top-\textit{k}, nitidez por temperatura e limite superior da exponencial após deslocamento. Todos verificados por \texttt{lake env lean} com Lean v4.34.1 e Mathlib rev.\ v4.34.1, restritos a \textit{propext}, \textit{Classical.choice} e \textit{Quot.sound}, sem \texttt{sorryAx}. Cada linha da Tabela~\ref{tab:ax} associa o nome do lema no código (por exemplo, \texttt{softmax\_soma\_um}, \nolinkurl{roteamento_guloso_otimo}, \texttt{sem\_overflow\_apos\_max}) aos três axiomas; a uniformidade das onze linhas é o conteúdo auditável da alegação ``sem atalhos''.
\textit{2. Por que esses fatos elementares merecem formalização?} Porque cada propriedade alegada do roteador ou é provada por máquina, ou é medida com incerteza declarada, ou é renunciada (Seção~\ref{cap:intro}). A formalização fixa a semântica dos pesos e da decisão pelo máximo, impedindo que o texto use esses fatos como intuição não auditada. O leitor cético pode percorrer a Seção~\ref{cap:provas} teorema por teorema: cada subseção da Seção~de demonstrações integrais traz enunciado, derivação completa, instância numérica real e a necessidade das hipóteses, com os nomes de lemas entre parênteses idênticos aos do Apêndice~\ref{ap:lean}.
\textit{3. T3 garante que o roteador escolhe o melhor agente?} Não. T3 estabelece existência de maximizador em conjunto finito não vazio, por invocação direta do máximo de imagem (\nolinkurl{Finset.exists_max_image}): a finitude impede fuga ao infinito e o não-vazio garante testemunha; sem finitude ($\mathbb{N}$ com identidade), não há máximo. Que o ``melhor'' formal corresponda ao melhor agente real exigiria preservar hipóteses pelo filtro, pelos escores, pelo desempate e pela aritmética (Seção~\ref{cap:discussao}), o que não foi demonstrado. A resposta honesta tem, portanto, duas camadas: dentro da matemática, o máximo existe; entre a matemática e o código, a ponte é trabalho futuro (refinamento, Tabela~\ref{tab:futuro}).
\textit{4. Quais hipóteses de T10 são essenciais e por quê?} Maximizador único e ao menos dois candidatos, como hipóteses suficientes. A derivação compara $A(\tau_1)S(\tau_2)$ contra $A(\tau_2)S(\tau_1)$ expandindo em somas: o termo $j=i^*$ é igual dos dois lados e cada outro termo é estrito pela identidade $(a_1+b_2)-(a_2+b_1)=(a-b)(1/\tau_1-1/\tau_2)>0$ com $a-b=f_{i^*}-f_j>0$; somar com a testemunha estrita decide por divisão de desigualdades. Sem não-trivialidade a tese é falsa: no singleton, esfriar nada altera --- foi exatamente essa falsificação que o compilador farejou ao exigir \texttt{Nontrivial}. Máximos empatados podem ganhar massa quando existem escores inferiores; se todos os escores forem iguais, a temperatura não altera os pesos (Seção~\ref{cap:provas}). A instância registrada, em $[1,2,3]$, vai de $0{,}6652$ ($\tau=1{,}0$) a $0{,}8668$ ($\tau=0{,}5$); a Tabela~\ref{tab:temp} estende a leitura a quatro temperaturas ($0{,}5065\to0{,}6652\to0{,}8668\to0{,}9817$).
\textit{5. T10 passa ao limite $\tau\to0$ ou $\tau\to\infty$?} Não. T10 é a versão finita: com $0<\tau_1<\tau_2$, esfriar aumenta estritamente o peso do maximizador único. Os limites assintóticos ($w_{i^*}\to1$ quando $\tau\to0^+$; equalização em $1/n$ quando $\tau\to\infty$) são intuição não formalizada: com maximizador único $i^*$, cada razão $e^{(f_j-f_{i^*})/\tau}\to0$ no esfriamento e todas as razões tendem a 1 no aquecimento. A dissertação descreve essa intuição como tal, sem reivindicar prova (seção de limites de temperatura e Seção~\ref{cap:provas}); teoremas de limite constam como pacote futuro na Tabela~\ref{tab:futuro}.
\textit{6. O que T11 garante contra overflow?} Em reais exatos: se $f_i\leq M$, então $e^{f_i-M}\leq1$, por monotonicidade da exponencial e $e^0=1$; instâncias $e^{-2}\approx0{,}1353$ e $e^{0}=1{,}0$. O encadeamento prático é: subtrair o máximo antes de exponenciar mantém todo expoente $\leq0$ (T11), logo cada termo fica em $(0,1]$ em aritmética real exata. Não certifica subtração, exponenciação, soma ou divisão IEEE-754; underflow, arredondamento, entradas não finitas e operações intermediárias permanecem fora da prova (Seções~\ref{cap:provas} e~\ref{cap:discussao}). Sem majorante, nada limita os expoentes --- a hipótese é, portanto, load-bearing, não burocrática.
\textit{7. A invariância de T4 vale bit a bit no código?} Não foi demonstrado. A derivação formal: $e^{f_j-c}=e^{f_j}/e^c$ (\texttt{exp\_sub}); a soma de quocientes vira quociente das somas; a simplificação de corpos fecha sozinha (\texttt{field\_simp}), pois $e^c\neq0$. Em $[1,2,3]$ e $[-4,-3,-2]$, os pesos em aritmética real coincidem e os valores arredondados do exemplo coincidem. Mas a invariância exata não implica igualdade bit a bit para quaisquer entradas em ponto flutuante: o fuzz examina a propriedade empiricamente dentro da tolerância de $10^{-12}$ para o deslocamento, sem certificar todas as entradas (Seção~\ref{cap:resultados}).
\textit{8. O que o episódio do \texttt{sorryAx} em T7 ensina?} Que erro silencioso de aplicação (\texttt{sum\_eq\_single\_of\_mem} sem elemento explícito) gera axioma espúrio capturável só pela auditoria de axiomas: a soma colapsa no único termo e $e/e=1$, mas sem o elemento explícito o termo de prova apoiou-se em \texttt{sorryAx}. A correção exigiu o elemento explícito; a auditoria final registra os três axiomas em todos os teoremas, inclusive T7 (Tabela~\ref{tab:ax}, Apêndice~\ref{ap:sessao}). A lição metodológica está incorporada ao pipeline: o verificador bloqueia \texttt{sorry} e \texttt{admit} em duas camadas, e a auditoria captura o que o bloqueio sintático não vê.
\textit{9. Por que o compilador exigiu \texttt{Nontrivial} em T10?} Porque o singleton falsificaria a tese de nitidez: com um só candidato, seu peso é sempre 1 e ``esfriar aumenta o peso'' é falso. A exigência veio da mensagem \nolinkurl{failed to synthesize Nontrivial}, não de decisão do autor: a soma estrita com testemunha precisa de ao menos dois termos para existir testemunha estrita. É o segundo caso, ao lado do T7, em que a máquina forçou o enunciado honesto (Apêndice~\ref{ap:sessao}).
\textit{10. Por que Lean 4 com Mathlib, e não outra ferramenta?} Foi o ambiente escolhido e utilizado neste trabalho: kernel verificador com auditoria de axiomas por teorema (\texttt{\#print axioms}) e biblioteca matemática madura acessível no ambiente, com somas finitas, exponencial real e máximos reusados nas provas (Seções~\ref{cap:fundamentacao} e~\ref{cap:metodo}). Z3 e SymPy compõem o ecossistema de apoio, não o núcleo de prova. A Tabela~\ref{tab:taticas} e a prosa de táticas documentam que cada nome foi conferido na Mathlib pinada (rev.\ v4.34.1), nunca na memória --- inclusive a forma vigente \texttt{div\_lt\_div\_iff₀}, após a forma antiga falhar com \texttt{Unknown identifier}.
\textit{11. Qual a distância entre as provas e o código Python?} Um refinamento completo não foi demonstrado: elegibilidade, representações semânticas, desempate, arredondamento e execução dos agentes estão fora dos enunciados. As provas referem-se ao modelo abstrato; o mapa de dependências registra honestamente os papéis (T1 sustenta T5 e T7; T2/T6 sustentam a combinação das cabeças; T4 sustenta a implementação estável usada em T10; T3 sustenta a decisão por máximo ao lado de T8; T9 descreve a ordenação do top-\textit{k}, sem constituir sozinho prova do algoritmo concreto). Refinamento formal com estruturas é o primeiro pacote da Tabela~\ref{tab:futuro}, com aceite em \texttt{machine-checked} (Seções~\ref{cap:metodo},~\ref{cap:discussao} e~\ref{cap:conclusao}).
\section{Desenho empírico}
\textit{12. O oráculo sintético é arbitrário?} Sim, declarado como tal. A utilidade $\mathrm{cobertura}\times((1-q_w)+q_wq)$ foi estipulada pelo desenho; a confiança é proxy informativa da qualidade latente por construção, com correlação registrada de aproximadamente $+0{,}85$ (Seção~\ref{cap:metodo}). A bancada principal usa cinco sementes (11, 22, 33, 44 e 55), cada uma com piscina de 40 agentes e 40 tarefas, totalizando 200 decisões; comparam-se roteador, seleção aleatória, menor carga e maior cobertura, registrando arrependimento, acerto, latência e desvio da soma dos pesos. Resultados valem dentro desse desenho; validade externa é explicitamente renunciada (Seção~de ameaças à validade).
\textit{13. A correlação $+0{,}85$ entre confiança e $q$ não torna tudo circular?} Explicita a proxy privilegiada: a vantagem do roteador depende de informação sobre a qualidade latente fornecida por desenho, e a correlação de Pearson ($r=+0{,}85$) quantifica exatamente essa dependência, ao lado da independência entre cabeças (semântica $\times$ confiança $+0{,}05$) e da capacidade constante após o filtro. A remoção da confiança aumenta o arrependimento nas ablações (de $0{,}0292$ para $0{,}2144$ em $q_w=0{,}6$) sem estabelecer equivalência com a seleção aleatória; mesmo confiança igual à qualidade latente não garantiria preservação do ranking do oráculo, pois semântica e carga também influenciam a escolha --- configuração não avaliada (Seções~\ref{cap:resultados} e~\ref{cap:discussao}).
\textit{14. O que a variação de $q_w$ prova?} Apenas reescala: entre elegíveis com cobertura unitária, o oráculo reduz-se a $1-q_w+q_wq$, o ranking permanece inalterado para $q_w>0$ e o arrependimento é $q_w(q_{\max}-q_{\mathrm{escolhido}})$. Os números confirmam a proporção exata: atenção $0{,}0146\to0{,}0292\to0{,}0438$ (razão $1{:}2{:}3$, como $0{,}3{:}0{,}6{:}0{,}9$); sem confiança $0{,}1072\to0{,}2144\to0{,}3216$; anticorrelacionada $0{,}1802\to0{,}3604\to0{,}5406$ (Tabela~\ref{tab:mundos}). A estabilidade do ranking entre valores não constitui evidência independente de robustez a outros critérios de utilidade; os três valores não testam critérios ou ordenações diferentes do oráculo (Seção~\ref{cap:metodo}).
\textit{15. Bancada principal e ablações usam as mesmas instâncias?} Não. A bancada usa as sementes 11, 22, 33, 44 e 55; as ablações, as sementes efetivas 10011--10055, outro conjunto de 200 decisões com média igualitária, remoção e perturbações da confiança; as correlações usam sequência própria de sorteios. O pareamento compara estratégias dentro de cada execução das ablações, sem vincular decisões entre conjuntos; as diferenças pareadas e suas margens descrevem as instâncias das ablações, sem correção para agrupamento (Seção~\ref{cap:metodo}, Tabela~\ref{tab:pareado}).
\textit{16. Os intervalos são intervalos de confiança válidos?} São resumos exploratórios. A margem registrada é $1{,}96s_n/\sqrt{n}$ com $s_n$ de divisor $n$: com $n=200$ e margem $0{,}0105$, obtém-se $s_n\approx0{,}0758$ ($0{,}0105\times\sqrt{200}/1{,}96$), identidade que explica o cálculo preservado sem validar a independência das decisões. Como cada piscina gera 40 decisões que compartilham agentes, e os intervalos não consideram esse agrupamento em cinco piscinas, podem subestimar a incerteza entre piscinas; a não sobreposição de intervalos isolados não substitui a análise da diferença pareada. Reamostragem que respeite a unidade de agrupamento é possibilidade futura, mas cinco piscinas ainda limitariam a precisão; nenhum recálculo inferencial foi efetuado nesta revisão (Seções~\ref{cap:metodo},~de estatística e~\ref{cap:resultados}).
\textit{17. Por que só cinco piscinas e comparadores simples?} Limite do desenho executado: cinco piscinas; na bancada, seleção aleatória, menor carga e maior cobertura; nas ablações, média igualitária e quatro perturbações da confiança. A Tabela~\ref{tab:seeds} mostra a heterogeneidade entre sementes (atenção de $0{,}0129$ na semente 55 a $0{,}0752$ na semente 11; menor carga de $0{,}0918$ a $0{,}3157$), o que reforça que cinco piscinas limitam a interpretação inferencial. Benchmark externo executado com orçamento e critérios comuns, mais piscinas independentes e análise na unidade correta são requisitos futuros (Tabela~\ref{tab:riscos}, Tabela~\ref{tab:futuro}).
\textit{18. A condição ruidosa é só ruído gaussiano?} Não: adiciona-se ruído gaussiano com desvio $0{,}4$, limita-se o resultado a $[0,1]$ e aplica-se arredondamento inteiro, gerando confiança 0 ou 1 --- o tratamento combina perturbação e discretização binária. Em $q_w=0{,}6$, essa condição registra $0{,}1515$ ($+0{,}122$ sobre a atenção), com atenção vencendo ou empatando em $0{,}93$ das decisões pareadas (Tabela~\ref{tab:pareado}). A leitura honesta: a discretização binária destrói granularidade, mas a comparação permanece dentro do desenho das ablações (Seção~\ref{cap:metodo}).
\textit{19. O que decide a demonstração sobre cartões (05b)?} Elegibilidade e ranking sobre onze cartões do catálogo local e cinco descrições de tarefas construídas com seu vocabulário de tags; não executa os agentes nem mede sucesso na realização das tarefas. Quatro decisões foram singleton (peso $1{,}0$ --- T7 em operação); a quinta opôs \texttt{reviewer} a \texttt{code-reviewer} com cobertura, confiança e carga idênticas ($1{,}0$; $0{,}5$; $1{,}0$), e a semântica decidiu $0{,}5589$ com placar $0{,}505$ contra $0{,}495$ --- decisão puramente semântica, à margem mínima. A primeira decisão custou $27{,}8$\,ms (carga fria do casamento semântico) e as seguintes $0{,}4$--$0{,}5$\,ms; soma dos pesos $1{,}0$ em todas. Tarefas fora do vocabulário retornam vazio por gate --- resultado compatível com o contrato do filtro; a cobertura do vocabulário do catálogo ainda requer avaliação (Seção~\ref{cap:resultados}).
\textit{20. Como ler o caso Ana contra Cid?} Como verificação aritmética completa de uma instância-guia, conferível à mão. Cabeças medidas: semântica Ana $0{,}8417$ contra Cid $0{,}7703$; capacidade $1{,}0$ para ambas (o gate já havia decidido); confiança $0{,}9$ contra $0{,}7$; carga livre $0{,}9$ contra $0{,}6$. Utilidades: Ana $0{,}30\cdot0{,}8417+0{,}35\cdot1{,}0+0{,}25\cdot0{,}9+0{,}10\cdot0{,}9=0{,}9175$; Cid $0{,}30\cdot0{,}7703+0{,}35\cdot1{,}0+0{,}25\cdot0{,}7+0{,}10\cdot0{,}6=0{,}8161$. Pesos finais $0{,}5253$ contra $0{,}4747$ (conferência: $e^{0{,}9175}\approx2{,}5033$, $e^{0{,}8161}\approx2{,}2616$, soma $\approx4{,}7649$, quociente $\approx0{,}5253$). Com cobertura total dos elegíveis, a decisão correu nas outras três cabeças --- exatamente o fenômeno quantificado nas correlações, onde a capacidade é constante após o filtro (Seção~\ref{cap:resultados}). Dois casos-limite reais complementam: gêmeos que diferem só na descrição recebem $0{,}5053$ contra $0{,}4947$ (empates exatos decidem-se pelo identificador, deterministicamente) e candidato único recebe peso $1{,}0$ (Seção~\ref{cap:intro}).
\textit{21. A atenção venceu em todas as sementes?} Obteve o menor arrependimento nas cinco sementes avaliadas (Tabela~\ref{tab:seeds}): $0{,}0752$ (acerto $0{,}45$) na 11; $0{,}0172$ ($0{,}78$) na 22; $0{,}0337$ ($0{,}55$) na 33; $0{,}0716$ ($0{,}60$) na 44; $0{,}0129$ ($0{,}75$) na 55 --- contra faixas de $0{,}1350$--$0{,}2463$ (aleatório), $0{,}0918$--$0{,}3157$ (menor carga) e $0{,}1618$--$0{,}2319$ (maior cobertura). A variação entre sementes é dado registrado, não inferência sobre outras piscinas; a margem agregada $0{,}042\pm0{,}011$ (acerto $0{,}625\pm0{,}067$) resume as 200 decisões com os limites da Seção~\ref{cap:metodo}.
\section{Testes numéricos}
\textit{22. O que os 20\,000 vetores demonstram?} Que, na amostra examinada (tamanhos 1--50, magnitudes até $\pm1000$), não houve violações acima das tolerâncias ($10^{-9}$ para a soma, $10^{-12}$ para o deslocamento), com pior desvio de soma $\approx2{,}2\times10^{-16}$ --- ordem de grandeza do épsilon de máquina em dupla precisão. Examinei ainda preservação de argmax, limites, caso de candidato único e tratamento do conjunto vazio. Ausência de violações na amostra não certifica todas as entradas nem substitui prova de refinamento numérico; o relatório usa ``testes aleatórios de propriedades'' e ``tentativas de refutação'', nunca ``contraprovas'' (Seções~\ref{cap:metodo} e~\ref{cap:resultados}).
\textit{23. Houve segunda execução do fuzz?} A documentação de uma segunda execução mencionada em versão anterior não foi identificada de forma individualizável nesta revisão; o total relatado restringe-se ao registro local identificado, com transcrições das sementes 7 e 99 preservadas nos apêndices. Relatar apenas o registro identificável --- em vez de somar execuções não individualizáveis --- é a aplicação direta do padrão ``medido com incerteza declarada ou renunciado'' (Seção~\ref{cap:metodo}).
\textit{24. Por que citar Miller et al.\ (1990) para o fuzz?} Como antecedente de testes por entradas aleatórias na avaliação de utilitários Unix, não como validação dos resultados locais. A citação ancora o método no campo (experimentação disciplinada desde \cite{basili1986,wohlin2012}), sem transferir autoridade: o que sustenta os números é a transcrição preservada, não a linhagem (Seção~\ref{cap:resultados}).
\section{Literatura e originalidade}
\textit{25. As obras relacionadas foram lidas integralmente?} Não; a síntese limita-se aos resumos, com leitura integral futura declarada, e a seção de leitura por obra registra, obra por obra, apenas o que cada uma informa (forma herdada, variantes, alcance não usado, limites do não feito, palco externo não pisado, fronteiras de T10/T11, desenho do fuzz, protocolo externo pendente). Nenhuma obra foi executada: são referências de contexto, não resultados reivindicados (seções de trabalhos relacionados e leitura por obra).
\textit{26. Qual a relação com a literatura de calibração por temperatura?} Complementar, não concorrente: a literatura aprende temperaturas e relaciona perdas próprias \cite{komisarenko2024,xuan2025,xie2024}; aqui a temperatura é teorema elementar sem aprendizagem. A distinção é load-bearing porque a normalização por softmax não demonstra calibração probabilística --- correspondência entre probabilidades previstas e frequências observadas --- objetivo que não foi medido aqui: não foram calculadas curvas de calibração, erros de calibração ou escores probabilísticos, embora regras de escore próprias \cite{gneiting2007} ofereçam fundamentos para outra análise (Seções~\ref{cap:fundamentacao},~de estatística e~\ref{cap:discussao}).
\textit{27. A convergência com o arcabouço atencional valida o trabalho?} Não: convergência independente digna de nota, não validação. O palco externo (avaliação multidimensional \cite{agentbench2023}, lacuna de perspectiva \cite{perspectivegap2026}, recompensa de orquestração \cite{orchrm2026}, coordenação atencional \cite{orchestrator2025}) define o protocolo futuro, sem validar os resultados locais; nenhum foi executado aqui (seções de trabalhos relacionados e~\ref{cap:discussao}).
\textit{28. Onde está a novidade reivindicável?} Na composição: verificação por máquina, avaliação interna e auditoria amarradas sobre um mesmo objeto com trilha local de reprodução (pinos, hashes, sincronia entre teoremas, testes, SPEC e axiomas via \nolinkurl{scanners/diagnostic_pipeline}) --- não em cada peça isolada, e sem reivindicar originalidade dos fatos matemáticos elementares, declarados como tal na Discussão (seção de síntese e Seção~\ref{cap:discussao}).
\section{Limites, ética e reprodução}
\textit{29. Há participantes humanos ou dados sensíveis?} Não foram utilizados dados de participantes humanos nos experimentos descritos. A declaração ética exigida pelo programa e eventuais dispensas institucionais permanecem sob responsabilidade do autor; o estudo tampouco fundamenta alegações de superioridade externa ou de originalidade validada por avaliadores externos (Seção~\ref{cap:metodo}).
\textit{30. O que falta para reprodução independente?} Depósito público versionado com hashes e saídas: hoje os artefatos estão organizados localmente (pinos verificados por \texttt{git ls-remote} em 2026-10-09 com SHA-256 das licenças na Tabela~\ref{tab:pins}; transcrições de testes, JSONs de resultados e SPECs nos Apêndices~\ref{ap:codigo}--\ref{ap:specs}), mas a reprodução integral da fonte atual não foi revalidada nesta revisão editorial e o depósito público com identificação persistente está pendente. O objetivo de disponibilização pública (OE4) está, por isso, parcialmente atendido (Seções~\ref{cap:metodo} e~\ref{cap:conclusao}).
\textit{31. As duas revisões internas equivalem a revisão por pares?} Não. Duas revisões internas automatizadas examinaram bancada e teoremas; não constituem revisão cega nem validação externa. As objeções sobre proxy, comparadores, analogia Q/K/V e lacuna numérica são compatíveis com as limitações da Tabela~\ref{tab:riscos}: as ablações acrescentam comparações entre ponderações e perturbações, mas não eliminam essas objeções (Seção~\ref{cap:discussao}).
\textit{32. A trilha auditável tem aplicação educacional comprovada?} Não avaliada com estudantes nesta pesquisa; é hipótese de aplicação futura, com material didático potencial para ensino de avaliação e validação: as derivações passo a passo, as instâncias conferíveis à mão (como Ana contra Cid) e o diário de erros reais oferecem matéria-prima didática, mas sem avaliação com estudantes nada se conclui sobre eficácia de ensino (Seções~\ref{cap:intro} e~\ref{cap:discussao}).
\section{Leitura guiada das tabelas e dos apêndices}
\textit{33. Como ler a Tabela~\ref{tab:ax}?} Linha por linha: cada teorema aparece com seu nome no código e os três axiomas auditados por \texttt{\#print axioms}. A uniformidade das onze linhas é a evidência; a exceção histórica (T7 com \texttt{sorryAx} na primeira versão) está documentada e corrigida. Quem duvidar pode reproduzir: o Apêndice~\ref{ap:lean} contém o código integral e o Apêndice~\ref{ap:reproducao} informa o ambiente (Lean 4 v4.34.1, Lake 5.0.0, elan 4.2.4, Mathlib rev.\ v4.34.1, 4760 oleans) e o comando.
\textit{34. Como ler a Tabela~\ref{tab:pins}?} Como identidade criptográfica de terceiros sem executar terceiros: cada repositório (leanprover/lean4, mathlib4, batteries, aesop, Z3Prover/z3, sympy/sympy) comparece com os 8 primeiros dígitos do commit HEAD verificado por \texttt{git ls-remote} e os 8 primeiros dígitos do SHA-256 da licença. A tabela fixa ``quais terceiros''; não afirma que foram auditados por dentro.
\textit{35. Como ler as Tabelas~\ref{tab:temp} e~\ref{tab:seeds}?} A primeira, coluna por coluna: ao esfriar $\tau$ de $2{,}0$ a $0{,}25$ em escores $[1,2,3]$, o peso máximo sobe $0{,}5065\to0{,}6652\to0{,}8668\to0{,}9817$ enquanto os demais murcham --- ilustração numérica de T10, não prova adicional. A segunda, linha por linha: cada semente tem sua piscina e sua dificuldade (a atenção varia de $0{,}0129$ a $0{,}0752$); a regularidade a notar é que a coluna da atenção é a menor em todas as cinco linhas, e que os comparadores trocam de ordem entre sementes (menor carga é o pior na 11 e o melhor comparador na 55).
\textit{36. Como ler as Tabelas~\ref{tab:bench} e~\ref{tab:pareado}?} A primeira mostra o reescalonamento: cada coluna multiplica a anterior pela razão dos $q_w$, e os parênteses dão a diferença contra a atenção na mesma coluna. A segunda traduz para decisões: diferenças pareadas com margens exploratórias (por exemplo, $+0{,}185\pm0{,}024$ sem confiança) e a fração em que a atenção vence ou empata ($0{,}92$--$0{,}99$). A leitura correta é amostral e pareada; a incorreta, tratar margens como significância entre mundos distintos.
\textit{37. Como ler as tabelas do Apêndice~\ref{ap:sessao}?} A de micro-provas como cobertura do vocabulário: quinze táticas ou combinações, cada qual com sua miniatura provada e veredito \texttt{machine-checked}, fechando a lacuna do vocabulário estático que causava \texttt{syntax\_error} indevido. A do diário como taxonomia de falhas reais: notação, assinatura, vocabulário, hipótese, nome e o caso silencioso --- cada linha com mensagem verídica do compilador e correção aplicada, nunca tentativa cega.
\textit{38. O que há nos demais apêndices?} Código integral dos experimentos e do pipeline (Apêndice~\ref{ap:codigo}), resultados em JSON (Apêndice~\ref{ap:json}, incluindo \nolinkurl{resultado_bancada.json} fonte dos valores da bancada), transcrições da suíte de testes e do ambiente (Apêndice~\ref{ap:testes}), histórico de auditoria com os ciclos R747--R778, e as quatro SPECs formais (Apêndice~\ref{ap:specs}). São artefatos históricos preservados; a execução integral da fonte atual não foi revalidada nesta revisão editorial.
\textit{39. Como os objetivos específicos se distribuem nos capítulos?} OE1 (verificação por máquina) nas Seções~\ref{cap:provas} e de demonstrações; OE2 (comparação interna com intervalos exploratórios) na Seção~\ref{cap:resultados} com método na Seção~\ref{cap:metodo}; OE3 (auditoria de axiomas, alegações, referências e código) nas auditorias, tabelas de pinos, micro-provas e referências; OE4 (limites declarados e trilha pública) nas Seções~\ref{cap:metodo},~\ref{cap:discussao} e~\ref{cap:conclusao} --- este, parcialmente atendido, com depósito pendente.
\textit{40. O que muda se o leitor só tiver o PDF?} Nada no verificável pelo texto: cada número remete a tabela ou apêndice contidos no próprio documento; cada prova remete a enunciado e auditoria impressos; cada limite está escrito. O que o PDF não entrega é a execução: sem o depósito público, reproduzir exige o ambiente descrito e os artefatos locais. Essa é a fronteira exata entre ``auditável no papel'' e ``reproduzido por terceiro''.
\section{Segunda rodada: objeções externas simuladas}
As objeções abaixo simulam revisores externos a partir da literatura citada; as respostas mantêm-se nos limites do texto, indicando onde a objeção procede e o que a contempla no futuro.
\textit{41. Por que não comparar com AgentBench ou outro benchmark?} Porque nenhum foi executado aqui: a avaliação é interna ao simulador, e a comparação permanece como trabalho futuro com suíte compartilhada, orçamento próprio e critérios comuns (Tabela~\ref{tab:futuro}). Executar AgentBench \cite{agentbench2023} exigiria harness, orçamento e protocolo que o desenho atual não prevê; alegar comparabilidade sem executá-lo seria overclaim. A objeção procede como direção, não como refutação dos resultados internos.
\textit{42. Por que pesos fixos $0{,}30/0{,}35/0{,}25/0{,}10$ em vez de aprendidos?} Porque o objeto de estudo é a decisão auditável, não o aprendizado: cada cabeça é heurística explícita (similaridade, cobertura, confiança, carga) e cada peso da combinação é fixo e declarado, como mostra a conta de Ana contra Cid na Seção~\ref{cap:resultados}. Pesos aprendidos exigiriam dados de treino, validação e análise de generalização --- fora do desenho. A média igualitária nas ablações compara outra ponderação dos mesmos componentes, sem aprender nada.
\textit{43. O desempate por identificador não é arbitrário?} É arbitrário no conteúdo e determinístico no contrato: empates exatos decidem-se pelo menor identificador, de modo que a seleção é determinística quando entradas, ambiente e demais operações permanecem fixos (seção de ordenação). A arbitrariedade é declarada em vez de escondida sob pseudo-aleatoriedade não registrada; deu-se o caso real dos gêmeos ($0{,}5053$ contra $0{,}4947$), em que nem gêmeos empataram na semântica.
\textit{44. Por que $n=200$ e só cinco piscinas?} Limite do desenho executado e registrado: cinco sementes gerando 40 decisões cada. A Tabela~\ref{tab:seeds} mostra que a dificuldade varia entre piscinas (atenção de $0{,}0129$ a $0{,}0752$), de modo que generalizar para outras piscinas seria extrapolação. Mais piscinas independentes e análise na unidade correta são requisitos futuros explícitos (Tabela~\ref{tab:riscos}).
\textit{45. Por que desvio-padrão com divisor $n$ em vez de $n-1$?} Porque os scripts preservados calculam $s_n$ com divisor $n$, e o texto relata o cálculo preservado em vez de recalculá-lo: com $n=200$ e margem $0{,}0105$, $s_n\approx0{,}0758$. A diferença entre divisores em $n=200$ é de $0{,}25\%$ no desvio --- irrelevante diante da incerteza entre piscinas, que é a limitação dominante e declarada. Nenhum recálculo inferencial foi efetuado nesta revisão (seção de estatística).
\textit{46. Pearson em variáveis assimétricas ou discretizadas faz sentido?} Como medida descritiva de associação linear, com os limites declarados: Pearson $r$ mede associação linear, sem demonstrar causalidade ou independência quando próximo de zero (seção de métricas). Os usos no texto são dois e modestos: quantificar a proxy ($r=+0{,}85$, checagem de desenho) e a independência aproximada entre cabeças ($+0{,}05$). Nenhum teste de significância é construído sobre esses coeficientes.
\textit{47. As latências são comparáveis entre si?} Não diretamente, e o texto não as compara: $2{,}17$\,ms é a latência registrada na bancada principal; $27{,}8$\,ms é a primeira decisão da demonstração sobre cartões (carga fria do casamento semântico); $0{,}4$--$0{,}5$\,ms são as decisões seguintes aquecidas. São três contextos de medição distintos, cada qual reportado no seu: a leitura correta é custo operacional documentado por contexto, não ranking de velocidade.
\textit{48. E se nenhum candidato for elegível?} O gate retorna vazio, por contrato do filtro --- comportamento observado com tarefas fora do vocabulário na demonstração sobre cartões, registrado como compatível com o contrato, não como falha. A cobertura do vocabulário do catálogo ainda requer avaliação (Seção~\ref{cap:resultados}).
\textit{49. O singleton aparece na bancada sintética?} O texto não registra a frequência de singleton na bancada; o caso é coberto formalmente por T7 e observado em operação nas quatro decisões singleton da demonstração sobre cartões (peso $1{,}0$) e no caso-limite do candidato único. A cobertura formal (T7) e a observação operacional (05b) são as duas evidências reivindicadas --- nem mais, nem menos.
\textit{50. O que falta exatamente para o Q/K/V formal completo?} Estruturas para consulta, chave e valor com produto escalado e mistura de valores, mais o refinamento que liga essas estruturas ao código: hoje a correspondência é parcial e estrutural (Tabela~\ref{tab:qkv}), com agregação e decisão demonstradas (T2, T6, T3, T9) mas sem produto escalado, temperatura aprendida ou mistura de valores. É o primeiro pacote da Tabela~\ref{tab:futuro}, com aceite em \texttt{machine-checked}.
\section{Terceira rodada: por objetivo específico}
\textit{51. OE1: qual a cadeia de evidência da verificação por máquina?} Arquivo \texttt{roteamento\_atencao.lean} (Apêndice~\ref{ap:lean}) $\to$ processamento por \texttt{lake env lean} no ambiente pinado $\to$ onze vereditos \texttt{machine\_checked} $\to$ auditoria \texttt{\#print axioms} por teorema (Tabela~\ref{tab:ax}) $\to$ bloqueio de atalhos em duas camadas (verificador e MCP). Cada elo é inspecionável: código, comando, ambiente, axiomas e sessão de erros (Apêndice~\ref{ap:sessao}).
\textit{52. OE2: por que estimação em vez de valor-\textit{p}?} Seguindo a recomendação de estimação em vez de significância pontual \cite{cumming2014,wasserstein2016}: o relato apresenta arrependimento e acerto com margens exploratórias e diferenças pareadas, sem nenhum teste de hipótese pontual. A escolha é coerente com o desenho (simulador interno, cinco piscinas) e evita a maquinaria inferencial que os dados não sustentam.
\textit{53. OE3: qual o inventário da auditoria?} Axiomas (Tabela~\ref{tab:ax}, três por teorema); alegações (triagem por regex com decisão humana documentada); referências (51 citadas = 51 entradas, 41 DOIs, sem entradas estranhas); código (inspeção e métricas na tradição de \cite{fagan1976,mccabe1976,chidamber1994}, pinos e SHAs na Tabela~\ref{tab:pins}, sincronia via pipeline R767). Cada item tem artefato e seção correspondentes; nenhum é alegação sem rastro.
\textit{54. OE4: como fica o mapeamento FAIR, item a item?} Localizável: caminhos, versões, hashes e transcrições identificados nos apêndices; parcialmente. Acessível: artefatos organizados localmente, depósito público pendente --- daí o ``parcialmente atendido''. Interoperável: JSONs de resultados, SPECs em Markdown, código Python e Lean em formatos abertos. Reutilizável: licenças e pinos registrados, protocolo de sincronia documentado. O mapeamento honesto tem um ``pendente'' explícito no segundo item, sem o qual o OE4 não se declara cumprido.
\textit{55. Como cada objetivo poderia ser falsificado?} OE1: um \texttt{sorryAx} na auditoria, um axioma além dos três, ou um veredito abaixo de \texttt{machine\_checked} --- todos verificáveis por reexecução. OE2: intervalos que, refeitos com agrupamento, invertam ordenações --- possibilidade registrada, análise pendente. OE3: uma citação sem entrada, um DOI inválido, um pino divergente --- todos conferíveis nos apêndices. OE4: já parcialmente não cumprido por declaração própria. A falsificabilidade de cada OE é o conteúdo científico desta resposta.
\textit{56. O que cada objetivo deliberadamente não cobre?} OE1 não cobre a implementação em ponto flutuante nem a execução dos agentes. OE2 não cobre benchmarks externos nem tarefas reais. OE3 não cobre revisão cega nem validação externa. OE4 não cobre o depósito público. As quatro renúncias estão escritas nas seções de método, discussão e conclusão; repeti-las aqui é o antídoto contra leitura expansiva.
\textit{57. Onde está cada objetivo no documento?} OE1 nas Seções~\ref{cap:provas} e de demonstrações, com código no Apêndice~\ref{ap:lean}; OE2 na Seção~\ref{cap:resultados} com método na Seção~\ref{cap:metodo}; OE3 nas auditorias, tabelas e apêndices de código, resultados, testes e SPECs; OE4 nas Seções~\ref{cap:metodo},~\ref{cap:discussao} e~\ref{cap:conclusao}. A resposta anterior (Q39) detalha a distribuição; esta fixa a correspondência biunívoca objetivo $\leftrightarrow$ seções.
\textit{58. Quais são os critérios de aceite dos pacotes futuros?} Os da Tabela~\ref{tab:futuro}, por pacote: Q/K/V formal com \texttt{machine-checked}; erro certificado com cotas verificadas; temperatura limite com teoremas de limite; benchmark externo com head-to-head executado; reprodução independente com mesmos números; seleção parcial com mesmos top-\textit{k} e menor latência. Cada aceite é observável por terceiro --- nenhum depende de julgamento privado.
\section{Quarta rodada: leitura dos registros brutos}
\textit{59. O que contém exatamente \nolinkurl{resultado_bancada.json}?} As chaves \texttt{seeds} (``11,22,33,44,55''), \texttt{n\_decisoes} (200), um bloco por estratégia com \texttt{regret\_medio}, \texttt{regret\_ci95}, \texttt{hit\_rate} e \texttt{hit\_ci95}, o bloco \texttt{atencao\_lat\_ms} com \texttt{media} e \texttt{max}, \texttt{softmax\_desvio\_max} e \texttt{amostra} com três decisões exemplo. Valores da atenção: $0{,}0421$, $0{,}0105$, $0{,}625$, $0{,}0671$ --- os quatro números que o texto arredonda para $0{,}042\pm0{,}011$ e $0{,}625\pm0{,}067$. Baselines: aleatório $0{,}1882$; menor carga $0{,}2167$; maior cobertura $0{,}2118$ --- a faixa $0{,}19$--$0{,}22$ do texto. Latência média $2{,}1652$\,ms com máximo $31{,}24$\,ms (primeira decisão da piscina); desvio máximo da soma $1{,}11\times10^{-16}$.
\textit{60. O que mostra a primeira decisão da amostra?} Semente 11, tarefa 0, quatro elegíveis: o oráculo marca $0{,}8823$; a atenção escolhe \texttt{ag\_00} com arrependimento $0$ e acerto $1$, em $31{,}24$\,ms, com soma softmax $1{,}0$; os três comparadores escolhem \texttt{ag\_07}, todos com arrependimento $0{,}0774$ e acerto $0$. É uma decisão individual preservada, não uma estatística: serve para auditar o formato do registro (escolhas nomeadas, utilidade do oráculo, latência, soma), não para inferir desempenho.
\textit{61. O que contém \nolinkurl{resultado_ablacao_q0p6.json}?} \texttt{n\_decisoes} (200), \texttt{qw} ($0{,}6$) e o bloco \texttt{estrategias}: a atenção ($0{,}0292$) e, por perturbação, \texttt{regret\_medio}, \texttt{regret\_ci95}, \texttt{diff\_pareado\_medio}, \texttt{diff\_pareado\_ci95} e \texttt{atencao\_vence\_ou\_empata} --- por exemplo, embaralhada $0{,}2243$ com diferença pareada $0{,}1950$ e taxa $0{,}92$; anticorrelacionada $0{,}3604$ com diferença $0{,}3312$. A estrutura do arquivo é a evidência de que ``pareado'' significa por decisão, não por agregado: cada diferença foi calculada dentro da mesma instância.
\textit{62. Como ler os nulos em \nolinkurl{correlacoes.json}?} Como informação, não como ausência: o arquivo traz \texttt{n\_decisoes} (200), uma \texttt{nota} explicativa e os blocos \texttt{pearson} e \texttt{correlacao\_cabecas}. Os valores exatos: arrependimento $\times$ elegíveis $0{,}0808$; $\times$ dispersão $-0{,}1611$; acerto $\times$ elegíveis $-0{,}1900$; $\times$ dispersão $0{,}1324$; confiança $\times$ $q$ $0{,}8538$; semântica $\times$ confiança $0{,}0513$. Os nulos (\texttt{semantic\_x\_capability}, \texttt{capability\_x\_confidence}, \texttt{capability\_x\_load}) têm a causa registrada na nota: variância zero, porque a capacidade é $1{,}0$ para todo elegível --- o gate já decidiu. O nulo documentado com causa é mais informativo que um zero espúrio.
\textit{63. O que há em \nolinkurl{per_seed.json} e \nolinkurl{estudo_caso_reais.json}?} O primeiro, um bloco por semente com os cinco arrependimentos e o acerto da atenção (a fonte da Tabela~\ref{tab:seeds}; semente 11: atenção $0{,}0752$, acerto $0{,}45$, $n=40$). O segundo, \texttt{n\_agentes} (11) e a lista de tarefas com \texttt{descricao}, \texttt{requeridas}, \texttt{elegiveis}, \texttt{excluidos}, \texttt{ranking}, \texttt{soma\_pesos}, \texttt{lat\_ms} e \texttt{cabecas\_top1} --- por exemplo, a primeira tarefa (refatoração com testes primeiro) com um elegível (\texttt{coder}, peso $1{,}0$, latência $27{,}79$\,ms) e dez excluídos. A granularidade por tarefa é o que permite conferir cada frase da demonstração sobre cartões.
\section{Quinta rodada: por pacote futuro}
\textit{64. Q/K/V formal: o que muda se for feito?} A Tabela~\ref{tab:qkv} ganha duas linhas demonstradas (escore QK com produto escalado; mistura de valores) e o refinamento liga as estruturas ao código; o aceite é \texttt{machine-checked}. Sem ele, a analogia permanece parcial e estrutural, como declarado.
\textit{65. Erro certificado: o que significa ``cotas verificadas''?} Cotas de erro de arredondamento e underflow demonstradas por máquina (na direção de \cite{bagnara2019,murray2026}), substituindo o par atual (lema exato + fuzz) por garantia. O aceite não é ``fuzz maior'', e sim teoremas sobre a execução numérica.
\textit{66. Temperatura limite: por que teoremas de limite?} Porque T10 é finito e a intuição assintótica ($w_{i^*}\to1$; equalização em $1/n$) está descrita sem prova. O pacote formalizaria os dois limites, fechando a lacuna entre a seção de limites e as provas.
\textit{67. Benchmark externo: o que conta como head-to-head executado?} Mesma suíte, mesmos critérios, orçamento próprio e protocolo pré-registrado, com os comparadores externos executados de fato --- não citados. Só então qualquer frase comparativa com plataformas passa a ser permitida; até lá, a renúncia permanece.
\textit{68. Reprodução independente: o que conta como ``mesmos números''?} Terceiro com os pinos reproduzindo os valores das tabelas dentro das tolerâncias e margens registradas, com divergências reportadas contra os registros. O depósito público versionado é pré-requisito, hoje pendente.
\textit{69. Seleção parcial: por que está na lista?} Porque a implementação ordena tudo ($O(n\log n)$) onde bastaria heap ($O(n\log k)$), com latência documentada de $2{,}17$\,ms para $n=40$. O aceite combina correção (mesmos top-\textit{k}) e desempenho (menor latência medida) --- melhoria de engenharia com critério duplo, não troca de algoritmo sem prova.
\section{Sexta rodada: por família da literatura}
\textit{70. Calibração: o que foi tomado e o que ficou de fora?} Tomada: a temperatura como parâmetro de nitidez e a fronteira do que T10 formaliza (versão finita, sem aprendizagem). De fora: aprendizagem de temperaturas, perdas próprias e qualquer métrica de calibração probabilística --- objetivos que a literatura mede e este trabalho declara não medir.
\textit{71. Ponto flutuante: idem?} Tomados: o desenho do fuzz (propriedades testáveis) e a consciência do que faltaria (aproximação correta, certificação via Lipschitz). De fora: a certificação total, explicitamente renunciada; T11 e o fuzz são o par reivindicado, nem mais.
\textit{72. Benchmarks: idem?} Tomados: o padrão de avaliação por tarefas (adaptado em miniatura sintética) e o protocolo externo pendente informado por eles. De fora: a execução de qualquer benchmark --- palco descrito, não pisado.
\textit{73. Verificação: idem?} Tomados: base axiomática, Lean, Mathlib, Z3 e SymPy como ecossistema, inspeção e métricas como disciplina de auditoria. De fora: qualquer alegação de que o núcleo verifica o sistema inteiro --- o refinamento é futuro declarado.
\textit{74. Método e reprodutibilidade: idem?} Tomados: experimentação disciplinada, relato sistemático e por estimação, dados e modelos documentados, FAIR como horizonte. De fora: a execução integral revalidada na data da revisão e o depósito público --- o ``parcialmente'' do OE4.
\textit{75. Linhagem da orquestração: por que contrato e quadro-negro?} Porque situam o roteador numa linhagem (alocação por contratos \cite{smith1980}; integração sob incerteza \cite{erman1980}) sem reivindicar herança de resultados: gates, cabeças, softmax e máximo são meios atuais instanciando temas antigos, cada camada auditável separadamente.
\section{Sétima rodada: tabelas restantes}
\textit{76. Como ler a Tabela~\ref{tab:qkv}?} Linha por linha, pela terceira coluna: Q, K e escore são estruturais ou parciais (sem produto escalado/temperatura); valor é parcial (sem mistura); agregação e decisão são demonstradas (T2/T6, T3/T9). O título da tabela (``sem transferência de propriedades'') é a tese da tabela: analogia que organiza, não que prova.
\textit{77. Como ler a Tabela~\ref{tab:riscos}?} Como contrato de honestidade em cinco linhas: cada limitação (oráculo arbitrário, cinco piscinas, modelo real vs floats, comparadores internos, artefatos locais) casa-se com um requisito futuro observável. A tabela é citada pela Discussão como o lugar onde as objeções ``são compatíveis'' --- isto é, previstas, não refutadas.
\textit{78. Como ler a Tabela~\ref{tab:futuro}?} Como seis promessas testáveis: cada pacote tem método e aceite observável por terceiro (ver Q58). A tabela transforma ``trabalhos futuros'' vagos em compromissos verificáveis --- o inverso de uma conclusão que empurra problemas para depois.
\textit{79. Qual a relação entre a Tabela~\ref{tab:taticas} e a de micro-provas?} A primeira é guia de leitura (tática $\to$ papel nas onze provas); a segunda é evidência de execução (tática $\to$ miniatura provada com veredito). Uma organiza a compreensão, a outra fecha a lacuna do vocabulário estático; juntas, respondem ``como ler'' e ``como crer''.
\textit{80. O que este marco registra?} Naquele ponto, oitenta itens sem nenhum dado novo: a demonstração de que um trabalho restrito pode ser defendido extensivamente quando cada número tem endereço e cada limite está escrito --- e a base sobre a qual as rodadas seguintes prosseguiram (o capítulo segue até o item 138).
\section{Oitava rodada: leitura cerrada por capítulo}
\textit{81. O que a Introdução promete e onde cada promessa é paga?} Três formas de evidência (prova por máquina, medida com incerteza, renúncia explícita): a primeira paga-se na Seção~\ref{cap:provas}, a segunda na Seção~\ref{cap:resultados} com método na Seção~\ref{cap:metodo}, a terceira nas Seções~\ref{cap:discussao} e~\ref{cap:conclusao}. A questão de pesquisa tem duas cláusulas (formalizar o núcleo; comportamento em avaliação sintética interna) e a conclusão responde nas mesmas duas cláusulas, sem introduzir nada novo.
\textit{82. Qual é exatamente a tese defendida?} Metodológica e restrita: verificação por máquina, teste numérico de propriedades, ablação e auditoria de alegações compõem um padrão de evidência auditável mesmo sem validação externa --- com otimalidade relativa à utilidade definida e evidência preliminar interna como teto. Não se defende superioridade sobre plataformas, generalização além do simulador, nem transferência de propriedades de Transformers treinados (Seção~\ref{cap:intro}).
\textit{83. Os casos-limite da Introdução são os mesmos dos resultados?} Sim, por desenho: gêmeos ($0{,}5053$ contra $0{,}4947$, desempate determinístico por identificador) e candidato único (peso $1{,}0$, T7 em execução) aparecem na Introdução como intuição e reaparecem na Seção~\ref{cap:resultados} como registros --- o primeiro na demonstração sobre cartões, o segundo em quatro decisões singleton. A repetição é rastreabilidade, não enchimento.
\textit{84. O que há de novo nas preliminares além de citações?} Três fixações operacionais: valores numéricos com quatro casas e soma unitária referida aos valores antes do arredondamento; a forma MaxEnt $w_i\propto e^{\beta f_i}$ com $\beta$ determinado pela restrição (softmax em $\beta=1$, temperatura em $\beta=1/\tau$); e o conteúdo formal de T2/T6, T3 e T9 em uma frase cada. O resto é linhagem declarada (Shannon, Kullback, Jaynes, Bridle, Gneiting) com função argumentativa explícita por parágrafo.
\textit{85. A seção de atenção neural transfere resultados de Transformers?} Explicitamente não: herda-se a forma (ponderação normalizada de múltiplas cabeças com decisão pelo máximo), sem pesos treinados, representações ou desempenho; cada cabeça é heurística explícita e cada peso é fixo e auditável. A frase de contenção está escrita para impedir exatamente essa leitura expansiva.
\textit{86. Para que servem as seções de verificação/SMT e de ponto flutuante nas preliminares?} A primeira situa Lean/Mathlib na verificação dedutiva (Hoare) com Z3 e SymPy como apoio, e inspeção/métricas como disciplina de auditoria; a segunda fixa que floats exigem análise própria e nomeia os dois níveis de evidência do trabalho (T11 + fuzz). Ambas preparam objeções antes que a banca as formule.
\textit{87. A correspondência Q/K/V é formal?} Parcial e estrutural, conforme a terceira coluna da Tabela~\ref{tab:qkv}: só agregação e decisão são demonstradas. A seção não contém nenhum teorema novo sobre Q/K/V --- contém a tabela e a renúncia, o que é o ponto.
\textit{88. Como o código Lean trata o maximizador único de T10?} Com hipótese explícita e estrita: \texttt{hmax : ∀ j, j ≠ iest → f j < f iest}, além de \texttt{Nontrivial}. A prova extrai testemunha (\texttt{exists\_ne iest}), trata o termo $j=i^*$ por comutatividade (\texttt{mul\_comm}) e os demais pela desigualdade estrita de diferenças, somando com testemunha (\texttt{Finset.sum\_lt\_sum}) e concluindo por divisão de desigualdades. Cada passo tem correspondente na derivação em prosa da Seção~de demonstrações.
\textit{89. E o singleton de T7 no código?} Com hipótese de unicidade explícita (\texttt{h : ∀ j, j = i}): a soma colapsa por \texttt{sum\_eq\_single\_of\_mem} com elemento explícito e conclui por \texttt{div\_self}. Sem o elemento explícito, a primeira versão gerou \texttt{sorryAx} --- o episódio documentado na Q8.
\textit{90. Como T9 lida com take e drop no código?} Reescrevendo a lista como \texttt{take k ++ drop k} (\texttt{take\_append\_drop}), decompondo a ordem da concatenação (\texttt{pairwise\_append}) e aplicando a componente cruzada (\texttt{hp.2.2}). É a anatomia do padrão ordenar-e-fatiar, independente do algoritmo de ordenação --- por isso a garantia vale para a implementação sem provar o sort.
\textit{91. Onde aparece a forma vigente de divisão de desigualdades?} Em T8 (\texttt{div\_lt\_div\_iff₀} com denominadores positivos) e no fecho de T10. O nome com subscrito é o detalhe que quebrou a primeira versão (\texttt{Unknown identifier}) e foi corrigido por consulta ao fonte pinado --- micro-história registrada no diário de erros.
\textit{92. O que a seção de estratégias (05c) acrescenta às tabelas?} A semântica das comparações: pesos fixos sem aprendizagem; média igualitária como outra ponderação; remoção como avaliação sem a informação; aleatória, menor carga e maior cobertura com o porquê de cada uma (disponibilidade fora do oráculo; empate entre plenamente elegíveis). Sem essa seção, os números das tabelas seriam comparações cegas; com ela, cada diferença tem interpretação declarada e limitada.
\textit{93. Por que a demonstração sobre cartões pertence aos resultados?} Porque é evidência de outro tipo: não estatística, mas operacional --- elegibilidade e ranking sobre vocabulário real do catálogo, com latências quente/fria, soma unitária e o caso de gate vazio. Complementa a bancada (estatística sintética) sem pretender substituí-la.
\textit{94. As correlações sustentam alguma decisão de desenho?} Uma, negativa: capacidade constante após o filtro confirma que o gate decide cobertura e a escolha corre nas outras cabeças --- por isso a discussão de Ana contra Cid. E uma advertência: $+0{,}85$ marca a proxy privilegiada, impedindo leitura ingênua da vantagem. Correlação como checagem de desenho, não como descoberta.
\textit{95. O que a Discussão faz que os capítulos não fizeram?} Compõe os alcances: provas documentam o modelo, bancada descreve a implementação em instâncias sintéticas; explicita pressupostos sem fornecer certificação integral. E recebe as objeções (proxy, comparadores, analogia, lacuna numérica) declarando-as compatíveis com a Tabela~\ref{tab:riscos} --- a Discussão é onde o trabalho se deixa atacar sem quebrar.
\textit{96. A literatura de temperatura e floats é usada contra ou a favor?} Como fronteira nos dois casos: calibração por temperatura aborda objetivos não medidos aqui (contra leitura expansiva da normalização); aritmética flutuante indica garantias próprias necessárias à execução (a favor do programa futuro, contra qualquer pretensão atual). Benchmarks de agentes: contexto futuro, sem validar o local.
\textit{97. As cinco contribuições C1--C5 são verificáveis uma a uma?} C1 (onze lemas + axiomas): Tabela~\ref{tab:ax} e Apêndice~\ref{ap:lean}. C2 (pipeline lake-env com bloqueio): Apêndice~\ref{ap:sessao} e SPECs. C3 (bancada com intervalos, pareamento e escala do oráculo): Tabelas~\ref{tab:bench}--\ref{tab:seeds} e JSONs. C4 (fuzz + auditorias): transcrições e seções de auditoria. C5 (trilha com pinos, hashes, sincronia): Tabela~\ref{tab:pins} e pipeline R767. Cada item tem endereço.
\textit{98. Por que a Conclusão repete os limites?} Porque responde à questão nos mesmos dois níveis sem introduzir nada novo, e porque OE4 parcial (artefatos locais, depósito pendente) precisa constar no fecho: uma conclusão que silencasse o pendente transformaria omissão em overclaim retroativo.
\textit{99. O que significa ``otimalidade relativa'' no fecho?} Exatamente T3 + oráculo declarado: ótimo em relação à utilidade definida, nunca superioridade universal --- a glosa do cabeçalho do próprio arquivo Lean, que define eficiência como otimalidade relativa à utilidade definida. A Conclusão herda a definição do código, não o contrário.
\textit{100. O que este segundo marco registra?} Cem itens sem nenhum dado novo, na mesma disciplina do marco anterior --- e a base sobre a qual as rodadas seguintes (dicionário operacional em diante) prosseguiram, até os 595 itens em sessenta e nove blocos. O que o capítulo não faz, em nenhuma rodada: introduzir dados, provar teoremas novos, executar benchmarks, validar externamente ou substituir a leitura dos capítulos.
\section{Nona rodada: dicionário operacional}
\textit{101. Arrependimento.} Diferença entre a maior utilidade do oráculo e a utilidade da escolha, por decisão; zero indica escolha ótima segundo esse oráculo. Operacionalmente: coluna \texttt{atencao\_regret} vs \texttt{oracle\_best} na amostra; média \texttt{regret\_medio} nos agregados.
\textit{102. Margem exploratória.} $1{,}96s_n/\sqrt{n}$ com $s_n$ de divisor $n$; coluna \texttt{regret\_ci95} nos JSONs. Lê-se como resumo da amostra com os limites da Seção~\ref{cap:metodo}, nunca como garantia frequentista sob agrupamento.
\textit{103. Diferença pareada.} Média, por decisão, do arrependimento do comparador menos o da atenção na mesma instância; chaves \texttt{diff\_pareado\_medio} e \texttt{diff\_pareado\_ci95}. É a única comparação que controla dificuldade entre estratégias.
\textit{104. Taxa de vitória ou empate.} Fração de decisões com arrependimento da atenção menor ou igual ao do comparador; chave \texttt{atencao\_vence\_ou\_empata} ($0{,}92$--$0{,}99$). Empate conta a favor por definição declarada.
\textit{105. Proxy.} Variável construída a partir da qualidade latente com ruído, informativa por desenho; quantificada por $r=+0{,}85$. Proxy não é medida independente: é o canal pelo qual o desenho alimenta o roteador.
\textit{106. Piscina.} Conjunto de 40 agentes e 40 tarefas gerado por uma semente; unidade de agrupamento ignorada pelos intervalos. Cinco piscinas: 11, 22, 33, 44, 55 na bancada; 10011--10055 nas ablações.
\textit{107. Elegível.} Candidato com cobertura completa das capacidades exigidas; o filtro decide elegibilidade e o máximo decide entre elegíveis. Fora do vocabulário, o conjunto é vazio por gate.
\textit{108. Singleton operacional.} Decisão com um só elegível: peso $1{,}0$ (T7 em operação), quatro ocorrências na demonstração sobre cartões. Não confundir com o singleton formal (hipótese \texttt{h} no código).
\textit{109. Gate.} Filtro de cobertura que precede a ponderação; já decidido quando a capacidade aparece constante ($1{,}0$) nas correlações. Gate e softmax são camadas separadas, auditáveis separadamente.
\textit{110. Machine-checked vs syntax-verified.} O primeiro, termo de prova conferido pelo núcleo no projeto lake (único veredito que sustenta alegação de prova); o segundo, análise sintática sem compilador (sem alegar compilação). A distinção é o coração do verificador em duas camadas.
\section{Décima rodada: anatomia dos scripts}
\textit{111. Como é feito um agente sintético?} Em \texttt{make\_pool}: identificador \texttt{ag\_00}--\texttt{ag\_39}, 3 a 6 capacidades sorteadas de um vocabulário de 12 (\texttt{formal}, \texttt{revisao}, \texttt{dados}, \texttt{codigo}, \texttt{geo}, \texttt{bio}, \texttt{nlp}, \texttt{visao}, \texttt{ml}, \texttt{qa}, \texttt{latex}, \texttt{etica}), qualidade latente uniforme em $[0,1)$ guardada fora do cartão, carga uniforme em $[0,0{,}9]$ e confiança $q+\mathrm{gauss}(0;0{,}15)$ limitada a $[0,1]$. O roteador recebe o cartão sem a qualidade latente: a separação entre \texttt{pool} e \texttt{latent} é o mecanismo da não-circularidade.
\textit{112. Por que a tarefa exige ao menos quatro elegíveis?} \texttt{make\_task} sorteia 2--3 requisitos e só aceita a tarefa com $\geq4$ elegíveis, tentando até 50 vezes (senão \texttt{RuntimeError}): com menos candidatos, o ranking não exercita escolha, e o singleton já tem seu caso dedicado. A descrição soma duas palavras de ruído (\texttt{análise}, \texttt{síntese}, \texttt{relatório}, \texttt{urgente}, \texttt{rotina}, \texttt{projeto}) para que a cabeça semântica tenha texto além dos requisitos.
\textit{113. Como os comparadores escolhem no código?} Aleatório: \texttt{rng.choice} entre elegíveis (usa o mesmo gerador, logo é reproduzível por semente). Menor carga: mínimo de \texttt{load}. Maior cobertura: máxima interseção com os requisitos, com \texttt{rng.choice} no desempate. Acerto: igualdade exata entre a utilidade do escolhido e a melhor do oráculo --- critério estrito, sem tolerância.
\textit{114. O que cada linha de decisão registra?} Semente, tarefa, número de elegíveis, melhor do oráculo; por estratégia, escolha, arrependimento e acerto; mais latência da atenção por \texttt{perf\_counter} e soma dos pesos do ranking. O agregado usa \texttt{ci95} (média por \texttt{fmean}, margem $1{,}96s_n/\sqrt{n}$ com \texttt{pstdev}) e imprime o resumo no console além de gravar o JSON com as três primeiras linhas em \texttt{amostra}.
\textit{115. Como as ablações atacam a confiança?} Reaproveitando \texttt{make\_pool}, \texttt{make\_task} e \texttt{ci95} do script principal, com semente deslocada ($10\,000+$ seed): embaralhada permuta os escores (quebra a proxy); anticorrelacionada usa $1-v$ (inversão adversarial); ruidosa soma $\mathrm{gauss}(0;0{,}4)$, limita e arredonda (discretização binária). Os dois comparadores fortes operam sobre as cabeças reais: média igualitária das quatro e renormalização sem confiança (pesos$/0{,}75$). O desempate é determinístico por $(-\mathrm{peso}, \mathrm{id})$.
\textit{116. De onde vêm 10011--10055?} De $10\,000+$ semente, deslocamento que separa a sequência das ablações da bancada principal mantendo determinismo: mesmos números de semente-base (11--55), outro sorteio. É o mecanismo em código da afirmação ``outro conjunto de 200 decisões''.
\textit{117. Como as correlações são calculadas?} Reexecutando as sementes deslocadas e coletando por decisão: arrependimento, acerto, elegíveis, dispersão (melhor menos pior do oráculo), escores das quatro cabeças por elegível e pares (confiança, $q$). Pearson com divisor $n$ em covariância e desvios; desvio zero devolve \texttt{None}, impresso como ``nula (constante)''. Sem pacotes externos: só \texttt{statistics} e \texttt{math}.
\textit{118. Qual é a lógica adversarial do fuzz?} ``Tenta REFUTAR'': 20\,000 vetores com tamanhos em $[1,2,3,5,8,13,25,50]$ e magnitudes em $[1,10,100,1000]$, seis contadores de violação (soma, limites, invariância, argmax, singleton, vazio), piores desvios rastreados e saída com erro (exit 1) ante qualquer violação acima das tolerâncias ($10^{-9}$ soma, $10^{-12}$ deslocamento). O impresso \texttt{CONTRAPROVA\_OK}/\texttt{CONTRAPROVA\_FALHOU} é veredito de tentativa de refutação --- por isso o texto usa ``tentativas de refutação'', nunca ``contraprovas'' como prova.
\textit{119. Como o \texttt{\_softmax} do código espelha T1/T4/T11?} Vazio devolve vazio; senão subtrai o máximo, exponencia, soma por \texttt{fsum} e divide, com guarda \texttt{or 1.0}. A subtração é T4 em operação; o denominador positivo comum é T1; os expoentes $\leq0$ são T11. Três teoremas, cinco linhas.
\textit{120. O que o validador de pesos exige?} Exatamente as quatro cabeças, cada peso finito em $[0,1]$, soma 1 dentro de $10^{-12}$ --- caso contrário \texttt{ValueError}. É o conteúdo formal de T2/T6 em código defensivo: o validador que rejeita pesos fora de $[0,1]$. A cabeça semântica tenta o casamento por habilidades e recua para o comportamento legado se indisponível.
\section{Décima primeira rodada: por apêndice}
\textit{121. Apêndice Lean: o que conferir?} As 150 linhas na ordem T1--T11, o cabeçalho com SPEC e definição de eficiência, as hipóteses de não-vazio, não-trivialidade, maximizador estrito e unicidade, e as táticas-chave listadas no roteiro de inspeção --- mais a ausência de \texttt{sorry}, \texttt{admit} e axiomas declarados.
\textit{122. Trilha de reprodução: o que ela fixa?} Ambiente (versões e oleans), sementes das duas coletas, comandos com saída registrada, pipeline de sincronia e a ressalva de não-revalidação integral --- além dos pinos, da identificação do autor (ORCID, contato) e da pendência do depósito público.
\textit{123. Código dos experimentos: quantos arquivos?} Nove listagens: o Lean, os quatro da bancada (principal, ablações, correlações, fuzz), o pipeline de sincronia, o roteador, a interface MCP e o teste do núcleo. A décima rodada deste capítulo percorre os seis primeiros; os demais têm sua sessão própria.
\textit{124. Resultados em JSON: quantos arquivos?} Seis: bancada principal, três ablações por $q_w$, correlações e por-semente. A quarta rodada percorre cada chave; o roteiro do Método ensina a conferência sem executar nada.
\textit{125. Transcrições: o que preservam?} A suíte de testes, a saída verbosa do \texttt{pytest} e a impressão do ambiente --- o ``papel'' da execução histórica, com a ressalva explícita de que a revisão editorial não reexecutou a bateria.
\textit{126. Histórico de auditoria: como lê-lo?} Como série temporal de ciclos R747--R782 com escores internos e objetos registrados, excluídos homônimos de outra linhagem pelo filtro de escore e tema; escores são autoavaliações, nunca validação externa.
\textit{127. SPECs: quais e para quê?} Quatro: R755 (teoremas), R754 (lake-env), R747 (núcleo de verificação) e R767 (pipeline de sincronia). Vinculam funcionalidades a especificações no protocolo RED--GREEN.
\textit{128. Sessão do verificador: por que em tabelas?} Porque a micro-prova por tática e o diário de erros pedem leitura tabular: quinze linhas de cobertura de vocabulário e seis de falhas reais com mensagens verídicas --- a consolidação pedida para o antigo módulo 04c.
\section{Décima segunda rodada: dicionário operacional II}
\textit{129. Semente efetiva.} $10\,000+$ semente-base nas ablações e correlações; separa sequências mantendo determinismo e rastreabilidade (10011--10055).
\textit{130. Utilidade do oráculo.} Cobertura vezes $((1-q_w)+q_wq)$; na bancada principal $q_w=0{,}6$ em código ($0{,}4+0{,}6q$). Declarada e arbitrária por desenho.
\textit{131. Cabeça.} Função de escore por aspecto (semântica, cobertura, confiança, carga) combinada pelos pesos fixos $0{,}30/0{,}35/0{,}25/0{,}10$; quase descorrelacionadas salvo a proxy.
\textit{132. Renormalização.} Divisão dos pesos restantes por $0{,}75$ ao remover a confiança ($1-0{,}25$); mantém soma unitária sem a cabeça ausente.
\textit{133. Testemunha.} Elemento cuja existência torna estrita uma desigualdade de somas: \texttt{exists\_ne} em T10, elemento explícito em T7. Sem testemunha, só desigualdade fraca.
\textit{134. Majorante.} Valor $M$ que limita escores por cima em T11; sem ele, nada limita os expoentes. Hipótese load-bearing, não burocrática.
\textit{135. Reescala.} Multiplicação do arrependimento por $q_w$ entre elegíveis com cobertura unitária; explica a razão $1{:}2{:}3$ e a estabilidade do ranking.
\textit{136. Amostra.} Três primeiras decisões preservadas no JSON da bancada; formato auditável, não base inferencial.
\textit{137. Tag de arquivo.} Sufixo \texttt{\_q0p3}/\texttt{\_q0p6}/\texttt{\_q0p9} que identifica a ablação pelo parâmetro; confere-se o $q_w$ interno contra o nome.
\textit{138. Falsificador.} Observação que derrubaria uma alegação: \texttt{sorryAx} na auditoria (OE1), inversão sob agrupamento (OE2), citação sem entrada (OE3), depósito ausente assumido como feito (OE4). A lista fecha o capítulo onde ele começou: cada número com endereço, cada limite escrito.
\section{Décima terceira rodada: anatomia do roteador, da MCP e do pipeline}
\textit{139. O que o gate rejeita, exatamente?} Cartão fora de \texttt{available} (\texttt{status\_not\_available}); capacidades inválidas ou ausentes; identificador ausente ou vazio; identificador duplicado (remove ambas as ocorrências e registra \texttt{duplicate\_agent\_id}); e capacidades requeridas faltantes, listadas nominalmente (\texttt{missing\_capabilities:...}). Cada exclusão carrega seus motivos em dicionário --- o gate explica, não apenas filtra.
\textit{140. O que acontece sem elegíveis?} \texttt{\_evaluate} devolve ranking vazio com cabeças vazias, utilidade vazia e pesos declarados, sem quebrar: é o caminho de código do gate vazio observado na demonstração sobre cartões. A robustez ao vazio é comportamento testado, não acidente.
\textit{141. Por que o índice posicional é fixo e descartado?} O incorporador recebe \texttt{positional\_index=0} com comentário explícito (o índice global não pode alterar uma decisão idêntica), e \texttt{route} descarta o parâmetro por compatibilidade de API. Duas decisões idênticas produzem o mesmo ranking: determinismo por construção, declarado no código.
\textit{142. Como utilidade vira ranking?} Soma ponderada por \texttt{fsum} dos quatro escores vezes os pesos validados; softmax sobre as utilidades; ordenação por ($-\mathrm{peso}$, $\mathrm{id}$). O desempate pelo identificador é o determinismo dos gêmeos ($0{,}5053$ contra $0{,}4947$); a soma unitária dos pesos finais é T1 em operação, conferida por linha (\texttt{softmax\_sum}) na bancada.
\textit{143. Para que serve \texttt{explain}?} Expor tudo: tarefa, requisitos, elegíveis, excluídos com motivos, escores por cabeça, utilidade, pesos e ranking --- a auditoria por decisão em um dicionário. É o instrumento que torna cada linha da bancada inspecionável sem reexecutar nada.
\textit{144. Quais são as quatro funções da MCP Lean?} \texttt{lean\_verify} (veredito fail-closed com projeto lake opcional), \texttt{lean\_export} (exportação com hash), \texttt{lean\_autoformalize} (rascunho informal $\to$ formal por domínio) e \texttt{lean\_saturate} (saturação por regras), mais procedência registrada. As duas camadas (verificador direto + MCP) são o bloqueio a \texttt{sorry}, \texttt{admit} e axiomas não declarados.
\textit{145. O que o pipeline de sincronia verifica?} \texttt{DiagnosticPipeline.varrer(plano)} percorre o plano lendo artefatos (\texttt{\_ler}), extrai nomes por regex (\texttt{\_nomes}) e checa nomes (\texttt{\_checar\_nomes}) e cobertura (\texttt{\_checar\_cobertura}) entre teoremas, testes, SPEC e axiomas. Oito testes herméticos; o diagnóstico preliminar registrou 18 verificações e duas advertências, sem falhas. É a amarração auditável da C5.
\textit{146. O que os oito testes do núcleo verificam?} Registro da SPEC R747, existência dos agentes e skills Lean, \texttt{lean\_verify} fail-closed, exportação com hash, autoformalização, manifesto do plugin e ausência de overclaim sem compilador --- o contrato do núcleo em forma executável, verde em $15{,}23$\,s com os do pipeline.
\textit{147. Onde mora o ``sem pesos treinados'' no código?} Em \texttt{HEAD\_WEIGHTS} fixo ($0{,}30/0{,}35/0{,}25/0{,}10$) validado na construção: qualquer alteração que não some 1 dentro de $10^{-12}$ levanta \texttt{ValueError} antes de qualquer decisão. Aprendizagem exigiria outro objeto; este recusa-se a aprender por construção.
\textit{148. Qual a fronteira entre roteador e simulador?} O roteador nunca vê qualidade latente: recebe cartões, descrições e requisitos, e devolve ranking. O simulador (piscina latente, oráculo, arrependimento) mora nos scripts da bancada. A separação em módulos é o mecanismo da não-circularidade --- e o motivo pelo qual T3 não escolhe ``o melhor agente'', só o maximizador declarado.
\section{Décima quarta rodada: por SPEC e por transcrição}
\textit{149. SPEC R747: o que especifica?} O núcleo de verificação total: cinco agentes, duas skills, MCP com quatro funções e plugin, com bloqueio fail-closed a atalhos --- oito testes verdes como aceite.
\textit{150. SPEC R754: o que especifica?} O ambiente lake com Mathlib e o veredito \texttt{machine\_checked} com degradação honesta (tática-Mathlib fora de projeto, sintaxe sem compilador): sem projeto, o máximo é \texttt{syntax\_verified}; com projeto, compilação real com limite e timeout fail-closed.
\textit{151. SPEC R755: o que especifica?} Os teoremas do núcleo matemático verificados por máquina com auditoria de axiomas --- onze vereditos e três axiomas por teorema, sem \texttt{sorryAx}.
\textit{152. SPEC R767: o que especifica?} O pipeline de sincronia com oito testes herméticos, amarrando teoremas, testes, SPEC e axiomas contra deriva silenciosa --- a que ``já mordeu uma vez'' na história do projeto.
\textit{153. O que a transcrição da suíte preserva?} Nomes de teste, vereditos e tempos: 63 aprovados em nove suítes no registro histórico, com a ressalva de não-reexecução editorial. É o papel da bateria, não a bateria.
\textit{154. E a saída verbosa do pytest?} O detalhe por teste (o que passou, em que ordem, em quanto tempo), permitindo conferir que ``63 em nove'' não é arredondamento nem resumo --- cada teste é nomeável.
\textit{155. E a impressão do ambiente?} Versões de Lean, Lake, elan, Mathlib, oleans em cache e dependências Python: o retrato do ambiente em que os registros foram produzidos, sem o qual ``reproduzir'' seria vago.
\textit{156. E as transcrições do fuzz?} As duas sementes (7 e 99) com contadores zerados e piores desvios: o papel das tentativas de refutação, restrito ao registro identificável.
\section{Décima quinta rodada: dicionário operacional III}
\textit{157. Motivo de gate.} Cadeia que explica uma exclusão (\texttt{missing\_capabilities:formal}, \texttt{duplicate\_agent\_id}...); auditoria por decisão, não só filtro.
\textit{158. Partição.} Divisão dos cartões em elegíveis e excluídos com motivos; duplicata remove ambas as ocorrências --- regra anti-ambiguidade.
\textit{159. Utilidade.} \texttt{fsum} dos escores vezes pesos por candidato; base do ranking e do arrependimento via oráculo.
\textit{160. Ranking.} Lista ordenada (id, peso) por ($-\mathrm{peso}$, $\mathrm{id}$); soma unitária por T1; primeira posição decide.
\textit{161. Explicação.} Dicionário completo de uma decisão via \texttt{explain}; inspeção sem execução.
\textit{162. Veredito.} \texttt{machine\_checked}, \texttt{syntax\_verified}, \texttt{syntax\_error} ou inválidos fail-closed; cada um com significado e limite distintos.
\textit{163. Procedência.} Registro de origem de cada artefato e veredito (função \texttt{\_provenance} na MCP); quem afirma o quê, rastreável.
\textit{164. Sincronia.} Amarração automatizada entre teoremas, testes, SPEC e axiomas; o que impede deriva silenciosa entre texto e código.
\textit{165. Transcrição.} Texto preservado de uma execução histórica (testes, fuzz, ambiente); papel com ressalva de data, não reexecução.
\textit{166. Aceite.} Condição observável por terceiro que declara um pacote cumprido (Tabela~\ref{tab:futuro}); sem aceite observável, sem pacote declarado.
\section{Décima sexta rodada: estudo de caso e pipeline em código}
\textit{167. Por que doze nomes e onze cartões?} A lista \texttt{SEL} tem doze nomes, mas só onze cartões carregam: \texttt{technical-writer} não tem correspondência no catálogo (nome ou arquivo ausente) e o filtro o ignora. O JSON registra \texttt{n\_agentes} 11 e o texto diz ``onze cartões'' --- o número que vale é o carregado, não o selecionado. A discrepância está documentada aqui para que ninguém cite doze.
\textit{168. Como o cartão nasce do catálogo?} Por duas expressões regulares sobre o Markdown: \texttt{name} e \texttt{tags} em colchetes, mais \texttt{description} entre aspas; capacidades viram lista, carga $0{,}0$ e confiança $0{,}5$ fixas para todos. Cartão sem nome na lista é ignorado; sem tags, ganha lista vazia (e o gate o exclui na prática).
\textit{169. Por que carga e confiança fixas?} Para isolar a semântica: com cobertura, confiança e carga idênticas entre contendores, só a cabeça semântica decide --- exatamente o desenho da quinta tarefa (\texttt{reviewer} vs \texttt{code-reviewer}). Fixar é controlar, declarado no código.
\textit{170. O que cada tarefa do estudo registra?} Descrição, requisitos, elegíveis, contagem de excluídos, ranking com pesos em quatro casas, soma em dez, latência em duas e as quatro cabeças do vencedor em quatro casas --- ou estruturas vazias se o ranking vier vazio. A impressão no console resume (elegíveis, topo, soma, latência) e o JSON guarda o integral.
\textit{171. O que ``hermético'' significa no pipeline?} Sem rede, subprocesso, LLM ou terceiros além da stdlib: só \texttt{os}, \texttt{re} e tipos. Todo resultado carrega rótulo \texttt{candidato\_a\_inspecao} e \texttt{exige\_validacao\_externa: True} --- o pipeline não substitui compilação, testes ou validação externa, por declaração própria no cabeçalho.
\textit{172. Quais são os cinco motivos de gap?} \texttt{arquivo\_ausente} (ilegível ou inexistente, no alvo ou na referência), \texttt{faltante} (esperado não achado), \texttt{excedente} (achado não esperado), \texttt{proibido} (padrão banido presente) e \texttt{cobertura} (nome sem correspondente na referência). Cinco motivos, cada um com \texttt{id} e detalhe: taxonomia fechada, sem ``outros''.
\textit{173. Como \texttt{varrer} valida o próprio plano?} Recusando plano malformado com \texttt{ValueError} (não-objeto, conjuntos vazios, tipo fora de \texttt{nomes}/\texttt{cobertura}, arquivo ou regex ausentes, referência inválida): o plano é inspecionado antes dos artefatos. Contam-se checagens, passam-se as sem gap, e o veredito agrega por \texttt{id}.
\textit{174. O que \texttt{\_nomes} garante?} Extração deduplicada do primeiro grupo de captura, com regex inválida ou sem grupo rejeitadas por exceção: nomes únicos na ordem de aparição, nunca multiconjunto. A deduplicação importa porque ``excedente'' compara conjuntos.
\textit{175. Para que servem esperados e proibidos juntos?} Para fechar nos dois sentidos: \texttt{esperados} cobra presença, a comparação com achados cobra ausência de extras, e \texttt{proibidos} veta padrões (como atalhos de prova) mesmo quando ninguém os esperava. É a forma executável do bloqueio em duas camadas.
\textit{176. O que o pipeline não faz?} Não compila, não testa, não valida externamente e não julga mérito: detecta dessincronia nominal entre artefatos declarados. O gap real pego na estreia (R767) prova utilidade, não onisciência.
\section{Décima sétima rodada: MCP e testes}
\textit{177. Como o bloqueio a \texttt{sorry} funciona?} Por regex antes de qualquer compilação: detectado \texttt{sorry} ou \texttt{admit}, devolve \texttt{is\_valid: False} com \texttt{incomplete\_sorry}, nome do teorema, táticas extraídas e log \texttt{bloqueado\_sorry\_admit} --- sem tocar no compilador. Fail-closed na primeira linha.
\textit{178. E o bloqueio a \texttt{axiom}?} Análogo, com \texttt{axiom\_nao\_declarado} e log \texttt{bloqueado\_axiom}: axioma exige declaração explícita e auditoria, nunca passagem silenciosa. As duas regexes são a primeira camada; a auditoria de \texttt{\#print axioms} é a segunda (a que pegou o \texttt{sorryAx} de T7).
\textit{179. O que é o rebaixamento anti-overclaim?} Sem compilador no ambiente, mesmo um veredito interno \texttt{machine\_checked} é rebaixado a \texttt{syntax\_verified} com nota explícita no log (``R747: sem compilador, rebaixado''). A regra está em código, não em manual: o sistema recusa-se a alegar o que não executou.
\textit{180. O que \texttt{lean\_export} custodia?} Nome, enunciado, táticas e imports (padrão: tática e reais básicos) formatados em fonte Lean com SHA-256 hex do fonte e procedência: identidade criptográfica do teorema exportado, não prova.
\textit{181. O que a procedência registra?} Presença de compilador, binário Lean, dica de versão da Mathlib e motor (verificador R444 + autoformalizador R445 via R747), mais diretório do projeto e via (\texttt{lake-env} ou \texttt{bare-lean}) por chamada. Cada veredito diz onde e como foi produzido.
\textit{182. O que cada um dos oito testes do núcleo verifica?} Registro da SPEC R747; existência dos agentes Lean; existência das skills Lean; \texttt{lean\_verify} fail-closed; exportação com hash; autoformalização; manifesto do plugin; e ausência de overclaim sem compilador --- o último testando exatamente o rebaixamento da Q179. Contrato executável, verde em $15{,}23$\,s com os do pipeline.
\textit{183. De onde vêm ``63 em nove''?} Do registro histórico da bateria completa (nove suítes, 63 aprovados), transcrito nos apêndices com a ressalva de não-reexecução editorial. Os 16 verificados nesta revisão (núcleo + pipeline) são subconjunto reexecutado, não a bateria.
\textit{184. Por que testar ``sem overclaim sem compilador''?} Porque é a propriedade que distingue este verificador: na ausência do compilador, o veredito correto é \texttt{syntax\_verified}, e o teste fixa esse comportamento contra regressão futura --- inclusive contra a tentação de ``melhorar'' o veredito sem compilador.
\textit{185. O que autoformalização e saturação fazem aqui?} Rascunho (informal $\to$ formal por domínio) e saturação por regras: utilitários de apoio registrados na MCP, sem alegação de correção das saídas --- ferramentas, não oráculos.
\textit{186. Qual a fronteira da MCP?} Ela opera vereditos, hashes e rascunhos; não compila por si (delega ao verificador/lake), não prova teoremas novos e não valida mérito matemático. Fronteira escrita nas docstrings e nos testes.
\section{Décima oitava rodada: dicionário operacional IV}
\textit{187. Cartão.} Dicionário com id, status, capacidades, descrição, carga e confiança; unidade de decisão do roteador e do gate.
\textit{188. Tag.} Vocábulo do catálogo que compõe capacidades e requisitos; fora do vocabulário, gate vazio.
\textit{189. Elegíveis e excluídos.} Partição com motivos; \texttt{eligible} ordenado pelo ranking, \texttt{excluded} mapeado a razões.
\textit{190. Ranking impresso.} Linha de console por tarefa (elegíveis, topo, soma, latência); resumo legível do JSON integral.
\textit{191. Hermético.} Sem rede, subprocesso, LLM ou terceiros além da stdlib; rótulo e flag em todo resultado do pipeline.
\textit{192. Gap.} Dessincronia nominal entre artefatos declarados, com motivo fechado e detalhe; o que o pipeline detecta.
\textit{193. Motivo.} Uma das cinco classes (\texttt{arquivo\_ausente}, \texttt{faltante}, \texttt{excedente}, \texttt{proibido}, \texttt{cobertura}); sem classe, sem gap.
\textit{194. Rótulo.} \texttt{candidato\_a\_inspecao} em todo resultado: inspeção requerida, nunca aprovação.
\textit{195. Degradação honesta.} Rebaixamento automático de veredito ante ambiente incompleto (sem compilador $\to$ \texttt{syntax\_verified}); testada contra regressão.
\textit{196. Integral.} Registro completo preservado (JSON, código, transcrição) contra o qual divergências se reportam; o texto remete, o apêndice contém.
\section{Décima nona rodada: por teste do núcleo}
\textit{197. O que \texttt{test\_spec\_r747\_registered} verifica?} Que a SPEC-935-R747 está registrada no SpecRegistry: sem registro, o ciclo SDD não existe. É o teste da própria especificação antes de qualquer funcionalidade.
\textit{198. O que o teste dos agentes exige de cada arquivo?} Existência dos cinco arquivos, frontmatter com \texttt{name:} e menção a bloqueadores (fail-closed, fail\_closed ou bloqueador, case-insensitive): presença, identidade e contrato em um só teste.
\textit{199. E o das skills?} Existência dos dois arquivos e menção a \texttt{sorry} e \texttt{admit}: a skill que não nomeia os atalhos não pode bloqueá-los. Exigência mínima, verificável por leitura.
\textit{200. Como o fail-closed é testado?} Com dois códigos: um teorema com \texttt{sorry} (deve invalidar com \texttt{incomplete\_sorry}/\texttt{axiom\_nao\_declarado}/\texttt{bloqueado\_sorry\_admit}) e o mesmo teorema com \texttt{simp} (deve validar com \texttt{syntax\_verified} ou \texttt{machine\_checked}). Par ruim/bom, sem ambiguidade.
\textit{201. Como a exportação é testada?} Exportando \texttt{soma\_zero} (\texttt{(n : Nat) : n + 0 = n}, [\texttt{simp}]): o fonte deve conter o teorema, o SHA deve ter 64 hexadecimais e deve igualar o recomputado por \texttt{hashlib} --- custódia verificada por recomputação independente no próprio teste.
\textit{202. E a autoformalização?} Com a frase ``para todo x natural, x + 0 = x'' no domínio álgebra: a saída deve conter \texttt{theorem} e \texttt{Mathlib}. Teste de formato, não de mérito matemático --- como declarado.
\textit{203. O que o manifesto garante?} \texttt{PLUGIN.md} e \texttt{lean-nucleo.json} existentes, com 5 agentes, 2 skills e chave \texttt{mcp}: a contagem que amarra código, agentes e interface. Cinco, dois e MCP --- os números do núcleo.
\textit{204. Como se testa ausência de overclaim?} Condicionalmente: só se não houver compilador, o teorema com \texttt{simp} deve voltar \texttt{syntax\_verified} (nunca \texttt{machine\_checked}) com \texttt{has\_compiler=False} na procedência. Com compilador presente, o teste nada exige --- honestidade condicional em código.
\section{Vigésima rodada: por transcrição}
\textit{205. O que \texttt{transcricao\_testes.txt} lista?} Sessenta e cinco linhas de identificadores \texttt{arquivo::classe::teste}: os oito do núcleo, os quatro do lake-env (incluindo prova errada que deve falhar e \texttt{sorry} bloqueado no projeto), os do pipeline e os demais até 63 coletados em $9{,}17$\,s. A lista é o índice nominativo da bateria.
\textit{206. O que o verboso acrescenta?} Cabeçalho da sessão (linux, Python 3.14.4, pytest 9.1.1), progresso por teste com percentual e o fecho: 63 aprovados em $248{,}00$\,s. Duração e plataforma fazem parte do registro --- reprodução sem ambiente é ficção.
\textit{207. O que a impressão do ambiente fixa?} Toolchain com hashes (Lean 4.34.1 com commit, Lake 5.0.0, elan 4.2.4), Python 3.14.4, pacotes (numpy, pytest, sympy, z3-solver com versões) e máquina (núcleos e memória). Quinze linhas que respondem ``onde isso rodou''.
\textit{208. O que as duas transcrições do fuzz registram?} Três linhas cada, idênticas nos números: 20\,000 vetores, zero violações nos seis contadores, pior soma $2{,}22\times10^{-16}$, pior deslocamento $0{,}00\times10^{0}$ e \texttt{CONTRAPROVA\_OK}. O deslocamento zero registrado autoriza citá-lo; antes desta leitura, o texto citava só a soma.
\textit{209. Quais são os testes do pipeline?} Nomes que dizem tudo: nome faltante, token proibido, cobertura faltante, plano malformado, arquivo ausente que vira gap e gap real no repositório --- cada motivo de gap com seu teste, mais o caso real da estreia.
\textit{210. O que os quatro testes do lake-env cobrem?} Prova correta verificada no projeto, prova errada que deve falhar, \texttt{sorry} bloqueado no projeto e registro da SPEC R754: o quadrilátero da verificação com compilador (acerto, erro detectado, atalho bloqueado, especificação).
\textit{211. Como ler ``63 em $248$\,s''?} Como custo histórico datado: a bateria completa levou quatro minutos na máquina impressa. Os 16 reexecutados nesta revisão (núcleo + pipeline) levaram $15{,}23$\,s; o teste lake, $224$\,s aquecido. Três tempos, três escopos, sem extrapolação entre eles.
\textit{212. Por que preservar transcrições em vez de só números?} Porque transcrição mostra formato (linhas por teste, percentuais, fecho) e permite recontagem independente: 65 linhas de índice, 73 de verboso, 15 de ambiente, 3+3 de fuzz --- 159 linhas de papel auditável.
\section{Vigésima primeira rodada: dicionário operacional V}
\textit{213. SPEC.} Especificação formal registrada que um teste referencia; sem registro, sem ciclo.
\textit{214. Frontmatter.} Cabeçalho com \texttt{name:} que dá identidade ao agente; o teste o exige por leitura.
\textit{215. Hermético (teste).} Sem rede, subprocesso ou terceiros; o que os oito do pipeline são.
\textit{216. Faltante.} Esperado não achado; motivo de gap com teste próprio.
\textit{217. Excedente.} Achado não esperado; o que impede deriva silenciosa por adição.
\textit{218. Proibido.} Padrão banido presente; o veto executável a atalhos.
\textit{219. Rebaixamento.} Troca automática de veredito ante ambiente incompleto, com nota no log; testada contra regressão.
\textit{220. Sementes 7 e 99.} As duas do fuzz, com transcrições idênticas nos números; a forma mais barata de interrogar floats, em duplicata.
\textit{221. Veredito \texttt{CONTRAPROVA\_OK}.} Zero violações acima das tolerâncias nas duas sementes; tentativa de refutação sobrevivida, não prova.
\textit{222. Papel.} As 159 linhas de transcrições contra as quais qualquer reexecução se reporta; o que esta rodada percorreu sem sair do texto.
\section{Vigésima segunda rodada: anatomia das cabeças}
\textit{223. Como a cabeça semântica pontua?} Tentando o casamento por habilidades (\texttt{top\_k=20}, melhor escore por agente); sem correspondência para um agente, recua ao legado naquele agente; ante exceção, recua em tudo. O legado é cosseno em $[0,1]$ entre vetor da tarefa e vetor do cartão. Três níveis (matcher, parcial, legado) com recuo declarado em cada falha.
\textit{224. O que o cosseno faz com vetores degenerados?} Devolve $0{,}5$ neutro ante normas nulas ou não finitas, fixa o cosseno em $[-1,1]$ e mapeia para $[0,1]$. O neutro $0{,}5$ é a posição declarada ante ausência de informação direcional --- nem otimismo, nem pessimismo.
\textit{225. Como a cobertura é calculada?} Fração das capacidades exigidas presentes no cartão; sem requisitos, $1{,}0$. É a cabeça que o gate já pré-selecionou: entre elegíveis com cobertura total, ela é constante --- por isso some das decisões e aparece constante nas correlações.
\textit{226. Como a confiança é calculada?} Tentando a média das confianças por habilidade (\texttt{top\_k=50}); sem correspondência, o valor do cartão com padrão $0{,}5$; sem matcher, o livro-razão legado. Na bancada, com o caminho legado, vale o escore do cartão ($q$ ruidosa) --- a proxy informativa por construção.
\textit{227. Como a carga vira escore?} Capacidade livre: $1-\mathrm{carga}$, com ausência equivalendo a livre. Cartão sem carga declarada pontua $1{,}0$ --- decisão de desenho documentada na docstring, não acidente.
\textit{228. O que \texttt{\_unit\_score} rejeita?} Booleanos (vão ao padrão), não-numéricos, não-finitos e fora de faixa (fixados em $[0,1]$). É o sanitizador que mantém todas as cabeças no intervalo que T2/T6 exigem --- a ponte entre código defensivo e hipótese formal.
\textit{229. As cabeças conversam entre si?} Quase não: semântica $\times$ confiança $+0{,}05$, semântica $\times$ carga $+0{,}02$, confiança $\times$ carga $+0{,}00$; capacidade constante. A independência aproximada justifica tratá-las como aspectos distintos na combinação --- sem alegar independência estatística formal.
\textit{230. De onde vêm os números de Ana e Cid?} Das quatro cabeças medidas na tarefa-guia: semântica $0{,}8417$/$0{,}7703$, capacidade $1{,}0$/$1{,}0$, confiança $0{,}9$/$0{,}7$, carga livre $0{,}9$/$0{,}6$ --- combinadas pelos pesos fixos em $0{,}9175$/$0{,}8161$ e normalizadas em $0{,}5253$/$0{,}4747$. Cada entrada é rastreável a uma cabeça nomeada.
\textit{231. Por que quatro cabeças e não três ou cinco?} Decisão de desenho declarada e validada (o validador exige exatamente as quatro): similaridade, cobertura, confiança e disponibilidade cobrem ajuste, pré-requisito, qualidade estimada e custo de ocupação. Outro número exigiria outro validador e outra justificação --- ausentes aqui.
\textit{232. O que as cabeças não fazem?} Não aprendem (pesos fixos), não transferem representações de Transformers treinados e não executam agentes: pontuam cartões. Cada ``não'' está escrito para conter a analogia Q/K/V nos limites da Tabela~\ref{tab:qkv}.
\section{Vigésima terceira rodada: por condição de ablação}
\textit{233. Embaralhada: o que quebra exatamente?} A ligação entre escore e agente: permutando os valores, a confiança vira ruído semântico com mesma marginal. Arrependimento $0{,}2243$ em $q_w=0{,}6$ ($+0{,}195$ pareado, taxa $0{,}92$): sem o vínculo com $q$, a decisão perde quase toda a vantagem --- medida da dependência da proxy, não do acaso.
\textit{234. Anticorrelacionada: por que é adversarial?} Porque $1-v$ inverte a informação: quem o gerador marcou bom parece ruim. Arrependimento $0{,}3604$ ($+0{,}331$, taxa $0{,}97$), o pior registrado: a decisão segue a proxy até contra a verdade. É o espelho da vantagem, não um comparador justo.
\textit{235. Ruidosa: o que a discretização muda?} Tudo na granularidade: após $\mathrm{gauss}(0;0{,}4)$, limitação e arredondamento, restam 0 ou 1. Arrependimento $0{,}1515$ ($+0{,}122$, taxa $0{,}93$): pior que a informativa, melhor que a embaralhada --- a hierarquia informa o valor da granularidade, dentro do desenho.
\textit{236. Média igualitária: por que é o comparador forte?} Porque usa os mesmos componentes com outra ponderação (média das quatro): isola o efeito dos pesos $0{,}30/0{,}35/0{,}25/0{,}10$ contra a ponderação neutra. Diferença $+0{,}035\pm0{,}011$ com taxa $0{,}99$: os pesos declarados vencem o neutro na amostra pareada, sem aprendizagem.
\textit{237. Sem confiança: o que a renormalização faz?} Divide os três pesos restantes por $0{,}75$, preservando soma unitária: a decisão sem a informação privilegiada. Arrependimento $0{,}2144$ ($+0{,}185$, taxa $0{,}98$): o custo da proxy em pontos de arrependimento, na amostra.
\textit{238. Como ler as três colunas de $q_w$ juntas?} Como uma só curva reescalada: cada condição multiplica-se por $1{:}2{:}3$ entre colunas, e as diferenças pareadas escalam junto. A tabela tem, portanto, uma dimensão real (condição) e uma nominal (escala) --- lê-la como três experimentos independentes é o erro que o texto proíbe.
\textit{239. O que as taxas $0{,}92$--$0{,}99$ significam e não significam?} Fração de decisões em que a atenção empata ou vence o comparador na mesma instância; não significam probabilidade de vencer fora da amostra nem teste de hipótese. Leitura amostral e pareada, como insiste cada legenda.
\textit{240. Qual condição falta?} Confiança igual à qualidade latente ($v=q$ sem ruído): isolaria ruído de viés da proxy. Não avaliada, declarada como tal --- a lacuna nomeada impede que o leitor a presuma coberta.
\section{Vigésima quarta rodada: dicionário operacional VI}
\textit{241. Casamento.} Busca por habilidades com \texttt{top\_k} amplo e melhor escore por agente; recuo legado ante falha.
\textit{242. Recuo legado.} Comportamento anterior (hashing/cosseno, livro-razão) usado ante indisponibilidade ou exceção; declarado, não escondido.
\textit{243. Neutro.} $0{,}5$ ante ausência de informação (normas degeneradas, cartão sem escore); posição declarada.
\textit{244. Fração de cobertura.} Interseção sobre requisitos; $1{,}0$ sem requisitos; constante entre elegíveis.
\textit{245. Livro-razão.} Fonte legada de confiança por cartão; o que vale na bancada.
\textit{246. Média per-skill.} Média das confianças das habilidades correspondentes; o que valeria com matcher.
\textit{247. Capacidade livre.} $1-\mathrm{carga}$; ausência equivale a livre, por docstring.
\textit{248. Fixação.} Valores idênticos entre contendores para isolar uma cabeça; técnica de controle do desenho.
\textit{249. Hierarquia.} Ordenação das condições por arrependimento (informativa $<$ ruidosa $<$ embaralhada $<$ anticorrelacionada); informa granularidade, dentro do desenho.
\textit{250. Lacuna nomeada.} Condição não avaliada declarada como tal ($v=q$); o que impede presunção de cobertura.
\section{Vigésima quinta rodada: reverificação Lean 4 em código}
\textit{251. O que foi reexecutado?} A suíte \texttt{test\_r755\_teorema\_roteamento.py} (4 testes) com cache aquecido: 4 aprovados em $191{,}40$\,s. A suíte elabora o arquivo de 150 linhas no projeto lake pinado, confere os vereditos e a auditoria de axiomas --- a mesma cadeia que produziu os 11 \texttt{machine\_checked} originais.
\textit{252. O que \texttt{machine\_checked} cobre desta vez?} Os onze teoremas do arquivo, sem \texttt{sorry}, \texttt{admit} ou axiomas declarados no fonte: a elaboração confere cada termo de prova e a auditoria confirma os três axiomas por teorema. Revalidação do mesmo código no mesmo ambiente, não prova nova.
\textit{253. E se o cache estiver frio?} A elaboração pode exceder o limite interno de 300\,s do teste (observado uma vez; a repetição concluiu em $224$\,s e, depois, em $191$\,s): custo de elaboração da Mathlib, não erro de prova. O protocolo registra o comportamento para que timeout não seja lido como refutação.
\textit{254. A reverificação valida o histórico?} Não: valida o estado atual do arquivo e do ambiente. O histórico (R755 e seguintes) permanece registro de execuções passadas com suas transcrições; a integridade entre passado e presente é garantida pela ausência de edição no fonte desde então, não por esta execução.
\textit{255. Por que Lean 4 e não conferência manual?} Porque o núcleo do assistente confere cada passo até axiomas explícitos, e \texttt{\#print axioms} lista o resíduo de confiança por teorema: três axiomas padrão, auditados um a um. Conferência manual dos mesmos 150 linhas não produziria esse resíduo mensurável.
\textit{256. O que os axiomas significam aqui?} O preço lógico das provas: \textit{propext}, \textit{Classical.choice} e \textit{Quot.sound} em todos os teoremas, sem minimalidade demonstrada (glossário). Usar o teorema é aceitar esse preço declarado --- nunca zero axiomas, nunca axioma escondido.
\textit{257. Táticas novas apareceram na reexecução?} Não: o vocabulário confinado na Mathlib pinada (\texttt{sum\_div}, \texttt{div\_self}, \texttt{sum\_le\_sum}, \texttt{exists\_max\_image}, \texttt{exp\_sub}, \texttt{field\_simp}, \texttt{add\_sum\_erase}, \texttt{sum\_eq\_single\_of\_mem}, \texttt{div\_lt\_div\_iff₀}, \texttt{pairwise\_append}, \texttt{sum\_lt\_sum}, \texttt{exp\_le\_one\_iff} e fechadores) elaborou sem \texttt{Unknown identifier}. Estabilidade de vocabulário entre revisões.
\textit{258. Qual a fronteira desta reverificação?} Ela cobre o modelo abstrato em reais exatos; não recertifica floats, agentes ou benchmarks. Cada reverificação futura repete a mesma fronteira --- por isso o protocolo a agenda sem prometer mais.
\section{Vigésima sexta rodada: narrativas por semente}
\textit{259. Semente 11: a mais difícil.} Atenção $0{,}0752$ (acerto $0{,}45$); aleatório $0{,}1889$; menor carga $0{,}3157$ (pior comparador aqui); maior cobertura $0{,}2319$. Mesmo na piscina mais dura, a atenção mantém menos da metade do arrependimento do segundo melhor.
\textit{260. Semente 22: a mais fácil para a atenção.} $0{,}0172$ com acerto $0{,}775$; comparadores entre $0{,}2276$ e $0{,}2463$, embaralhados entre si. Facilidade não é folga: a distância aos comparadores é a maior em termos relativos.
\textit{261. Semente 33: a mediana.} $0{,}0337$ (acerto $0{,}55$); aleatório $0{,}1719$; menor carga $0{,}2447$; maior cobertura $0{,}2033$. Perfil próximo às médias agregadas --- a piscina ``típica'' do conjunto.
\textit{262. Semente 44: a disputa interna.} $0{,}0716$ (acerto $0{,}60$); aleatório $0{,}1992$, menor carga $0{,}2038$, maior cobertura $0{,}2309$ --- os três comparadores empatados entre si em torno de $0{,}20$--$0{,}23$, todos a $\approx3\times$ a atenção.
\textit{263. Semente 55: a melhor.} $0{,}0129$ com acerto $0{,}75$; menor carga $0{,}0918$ (seu melhor desempenho, ainda $7\times$ pior); aleatório $0{,}1350$; maior cobertura $0{,}1618$. Cinco histórias, uma regularidade: a coluna da atenção é a menor em todas as linhas.
\section{Vigésima sétima rodada: o incorporador}
\textit{264. Por que dimensão 64?} Constante \texttt{D\_MODEL} do incorporador: vetores de 64 posições somando hashes de tokens, mais codificação posicional e normalização. Dimensão declarada, sem sintonia registrada --- hiperparâmetro fixo, não resultado.
\textit{265. Como um token vira vetor?} SHA-256 do token, 8 pares de bytes definindo índice (módulo 64) e sinal ($\pm1$): vetor esparso assinado determinístico. Mesmo token, mesmo vetor, em qualquer máquina --- determinismo por hash, sem modelo treinado.
\textit{266. Como o texto vira tokens?} Sublinhados viram espaços, camelCase é quebrado, e a expressão multilíngue extrai alfanuméricos em minúsculas: \texttt{cloudsql\_postgres\_ad} rende três tokens. Tokenização documentada com exemplo no código.
\textit{267. O que a ``LayerNorm simplificada'' faz?} Normalização L2 (divide pela norma, com guarda contra zero): vetores unitários para o cosseno. ``Simplificada'' no nome impede leitura como normalização de camada treinada.
\textit{268. Para que a codificação senoidal de amplitude $0{,}1$?} Marca de posição à Vaswani com amplitude pequena: o índice posicional entra sem dominar o conteúdo. Na bancada, o índice é sempre 0 (decisões idênticas não mudam) --- a marca existe, neutralizada por desenho.
\textit{269. Por que pesos textuais $\times2$ e $\times3$?} Requisitos repetidos duas vezes na tarefa, capacidades três vezes no cartão: ênfase declarada no vocabulário controlado sobre a prosa livre. Pesos de texto, não de cabeça --- não confundir com os $0{,}30/0{,}35/0{,}25/0{,}10$.
\textit{270. Qual o papel do incorporador na bancada?} Fornecer os vetores do caminho legado da cabeça semântica (matcher indisponível): determinístico por hash, sem embeddings neurais. Os escores semânticos medidos ($0{,}8417$ vs $0{,}7703$ em Ana/Cid) vêm desse caminho --- heurística explícita, não representação treinada.
\textit{271. O que o incorporador não faz?} Não aprende, não consulta modelos externos e não garante qualidade semântica: conta tokens com hash. A honestidade do caminho legado é exatamente essa modéstia declarada.
\section{Vigésima oitava rodada: dicionário operacional VII}
\textit{272. Dimensão.} 64 posições do espaço de hashing; hiperparâmetro fixo do incorporador.
\textit{273. Hash assinado.} SHA-256 fatiado em 8 pares índice/sinal; determinismo sem treino.
\textit{274. Token.} Unidade pós-quebra de snake/camelCase em minúsculas; o que se soma no vetor.
\textit{275. Posição senoidal.} Marca de índice com amplitude $0{,}1$; neutralizada em 0 na bancada.
\textit{276. Peso textual.} Repetição $\times2$/$\times3$ de requisitos/capacidades; ênfase declarada, não peso de cabeça.
\textit{277. Narrativa de semente.} Leitura linha a linha de uma piscina (valores, acertos, distâncias); dado, não inferência.
\textit{278. Revalidação.} Reexecução da mesma cadeia no mesmo ambiente (4 em $191$\,s); confirma presente, não passado.
\textit{279. Resíduo.} Os três axiomas por teorema; preço lógico declarado e auditado.
\textit{280. Modéstia.} Posição do caminho legado (conta tokens, sem garantias); a honestidade como especificação.
\textit{281. Fecho provisório.} 281 itens em trinta e três blocos sem nenhum dado novo; a contagem continua enquanto houver peça endereçável --- e para quando o autor mandar parar.
\section{Vigésima nona rodada: o resumo frase a frase}
\textit{282. ``Investiga a formalização e a avaliação interna''.} A primeira frase do resumo fixa o objeto duplo (núcleo matemático + implementação no simulador) e o adjetivo ``interna'' que governa todo o resto: nada fora do simulador é reivindicado adiante.
\textit{283. ``Combina provas, inspeção, bancada e testes aleatórios''.} Os quatro métodos mapeados: Seção~\ref{cap:provas}, apêndices de código, Seção~\ref{cap:resultados} e fuzz --- cada um com seção e artefato próprios, sem método órfão.
\textit{284. ``Onze lemas''.} A enumeração nominal (normalização, limites, maximizador, invariância, positividade, singleton, monotonicidade, top-\textit{k}, temperatura, limite da exponencial) com endereço por teorema na Seção~\ref{cap:provas} e auditoria na Tabela~\ref{tab:ax}.
\textit{285. ``Aritmética real exata, sem certificação integral''.} A ressalva que acompanha cada citação de prova no resumo: o escopo (reais) e a renúncia (floats) na mesma frase, impedindo leitura expansiva do restante.
\textit{286. ``$0{,}042\pm0{,}011$ contra $0{,}19$--$0{,}22$''.} Os quatro decimais do JSON ($0{,}0421$/$0{,}0105$) e a faixa dos três baselines ($0{,}1882$, $0{,}2167$, $0{,}2118$): cada dígito do resumo tem chave no Apêndice~\ref{ap:json}.
\textit{287. ``Ablações indicam dependência da proxy; $q_w$ reescala''.} As duas teses empíricas do resumo, cada qual com sua tabela (dependência nas diferenças pareadas; reescala na razão $1{:}2{:}3$) e seus limites escritos.
\textit{288. ``Intervalos exploratórios; 20\,000 vetores''.} As duas qualificações numéricas: agrupamento ignorado (resumo exploratório) e amostra finita sem violações (sem certificação universal). Números com ressalva embutida.
\textit{289. ``Sustentam propriedades e preliminares; pendências''.} O fecho em dois tempos (o que os artefatos sustentam; certificação, agrupamento, comparação externa e reprodução pendentes) espelha a Conclusão sem introduzir nada: resumo que fecha como o trabalho fecha.
\section{Trigésima rodada: o casamento semântico}
\textit{290. Por que duas dimensões (64 e 128)?} O incorporador legado opera em 64; o motor semântico, em 128 (``vs 64 do embedder legado'' no código). Módulos distintos, espaços distintos, sem mistura: cada caminho usa sua dimensão declarada.
\textit{291. Qual a hierarquia de provedores?} LiteRT-LM (Gemma/Qwen on-device), Colibri (OLMoE) e hashing SHA-256 stdlib, sondados nessa ordem com recuo silencioso ante exceção. O provedor efetivo fica registrado nas estatísticas --- hierarquia com rastreio, não caixa-preta.
\textit{292. Como a sonda escolhe?} Tentando cliente com checagem de saúde (\texttt{status ok} / \texttt{ready}); falhando tudo, ``hash''. A escolha acontece na construção do motor, uma vez, e vale para as incorporações seguintes.
\textit{293. Para que o cache de 500?} Chave SHA-256 truncada em 32 hex por texto, acertos contados e despejo do mais antigo além de 500 entradas: economia sem vazar memória. Estatísticas de chamadas e acertos expostas por \texttt{get\_stats}.
\textit{294. O que o handbook registra?} Habilidades por agente (id, nome, descrição, tags, confiança inicial $0{,}5$, episteme, categoria, tipo) com execuções anotadas: o cadastro que o casamento consulta. Registro conta habilidades e devolve o número.
\textit{295. Como o casamento responde?} Descrição vs habilidades com \texttt{top\_k} (20 na semântica, 50 na confiança), lista ordenada por escore: o roteador toma o melhor por agente (semântica) ou a média (confiança). Parâmetro e agregação distintos por cabeça, declarados na chamada.
\textit{296. O que é o singleton do módulo?} Instância única construída na importação: o casamento é estado compartilhado, não objeto por decisão. Conhecer o singleton explica por que o roteador o carrega preguiçosamente e recua ante falha.
\textit{297. Qual caminho valeu na bancada?} Os escores medidos são compatíveis com o caminho legado determinístico (hashing, sem modelos externos): o texto reivindica heurística explícita, e os números exibidos correspondem a esse caminho. O encadeamento tentar--recuar cobre os três ramos, de modo que cada escore sai de um ramo nomeado (matcher, parcial ou legado).
\section{Trigésima primeira rodada: dicionário operacional VIII}
\textit{298. Dimensão semântica.} 128 posições do motor de casamento; distinta dos 64 do legado, sem mistura.
\textit{299. Provedor.} Fonte de vetores (LiteRT, Colibri ou hash) escolhida por sonda e registrada nas estatísticas.
\textit{300. Sonda.} Tentativa ordenada com checagem de saúde; recuo silencioso documentado.
\textit{301. Cache.} 500 entradas com chave SHA e despejo do mais antigo; economia rastreada.
\textit{302. Handbook.} Cadastro de habilidades por agente com confiança inicial e execuções; o consultado.
\textit{303. Singleton.} Instância única na importação; estado compartilhado com carga preguiçosa.
\textit{304. Resumo (gênero).} Texto que fecha como o trabalho fecha: cada frase com endereço, cada número com chave.
\textit{305. Endereço.} Par seção-tabela-apêndice que localiza qualquer número citado; a unidade de auditoria deste capítulo.
\textit{306. Suficiência.} Hipótese que basta sem ser necessária (maximizador único, dois candidatos); o estatuto lógico de T10.
\textit{307. Fecho provisório II.} 307 itens em trinta e seis blocos; a contagem prossegue por peças endereçáveis.
\section{Trigésima segunda rodada: a conclusão parágrafo a parágrafo}
\textit{308. Por que a Conclusão retoma a questão em dois níveis?} Porque a pergunta tinha duas cláusulas (formalizar o núcleo; comportamento em avaliação sintética) e o fecho responde nas mesmas, sem introduzir nada novo: onze lemas verificados de um lado, avaliação preliminar no simulador do outro, artefatos locais de reprodução como lastro comum.
\textit{309. O que cada contribuição C1--C5 entrega de verificável?} C1: onze vereditos + auditoria (Tabela~\ref{tab:ax}, Apêndice~\ref{ap:lean}). C2: pipeline lake-env com bloqueio (Apêndice~\ref{ap:sessao}). C3: bancada com intervalos, pareamento e escala (Tabelas~\ref{tab:bench}--\ref{tab:seeds}). C4: fuzz + auditorias (transcrições, 51=51, pinos). C5: trilha com pinos, hashes e sincronia. Cinco itens, cinco endereços.
\textit{310. Como OE1--OE3 se declaram sustentados e OE4 parcial?} Pelos registros locais nos limites explicitados (provas, comparação interna, auditorias) para os três primeiros; pelo depósito público pendente para o quarto. A assimetria $3+1$ está escrita para impedir leitura de cumprimento total.
\textit{311. Para que serve a tabela de futuros na Conclusão?} Transformar pendências em compromissos observáveis: seis pacotes com método e aceite verificável por terceiro (Tabela~\ref{tab:futuro}). Sem ela, o fecho seria promessa; com ela, é contrato.
\textit{312. O que significa não introduzir nada novo no fecho?} Que cada frase da Conclusão remete a seção anterior: questão, níveis, contribuições, objetivos, tabela. A disciplina é verificável por inspeção --- e este capítulo a percorreu item a item nas rodadas por capítulo.
\textit{313. Por que repetir os limites no fecho?} Porque silenciá-los transformaria omissão em overclaim retroativo: OE4 parcial e depósito pendente constam no último parágrafo para que ninguém cite a Conclusão isolada como cumprimento.
\textit{314. ``Otimalidade relativa'' aparece quantas vezes e por quê?} No cabeçalho do código, na delimitação, na Conclusão e no resumo: a definição circula para fixar o teto em todo ponto de entrada do documento. Repetição deliberada contra leitura expansiva.
\textit{315. Qual a tese metodológica em uma frase?} Verificação por máquina, teste de propriedades, ablação e auditoria compõem evidência auditável mesmo sem validação externa --- com o teto declarado. Uma frase, quatro pilares, um limite.
\section{Trigésima terceira rodada: o handbook em código}
\textit{316. O que o registro de habilidade garante contra duplicata?} Idempotência por chave (agente, habilidade): registrado uma vez, devolvido intacto nas seguintes. Sem duplicatas no cadastro, sem dupla contagem no casamento.
\textit{317. Como as tags são enriquecidas?} Expansão de subtokens mais tokens do identificador do agente, deduplicados em ordem: \texttt{Cloud SQL PostgreSQL Specialist} rende \texttt{cloud}, \texttt{sql}, \texttt{postgresql}, \texttt{specialist}. O vetor incorpora nome, descrição, agente e tags --- contexto amplo por desenho.
\textit{318. Override ou inferência na episteme?} Explícita vence; sem ela, heurística determinística; ante falha, \texttt{None} (fail-open). Aqui o desenho diverge do verificador de propósito: episteme é anotação auxiliar, não porta de prova --- falhar fechado onde se anota seria paralisar o cadastro por detalhe.
\textit{319. Como a execução atualiza a confiança?} Média móvel ponderada: $0{,}7$ da taxa recente mais $0{,}3$ da anterior; custo por EMA $0{,}9/0{,}1$ quando positivo. Taxas, não palpites: sucessos sobre execuções, com memória que decai sem apagar.
\textit{320. Quais filtros o casamento aplica?} Cadastro vazio devolve vazio; confiança abaixo de $0{,}3$ exclui; sem interseção de tags, exclui. Só então similaridade, escore e ordenação: filtros primeiro, ranqueamento depois.
\textit{321. Como o escore composto é formado?} $0{,}6$ similaridade mais $0{,}4$ confiança, com peso brando epistêmico de $\pm10\%$ no máximo quando ambos os lados têm episteme identificada; similaridade e escore arredondados em quatro casas, ordenação descendente, corte em \texttt{top\_k}. Fórmula declarada, sem aprendizado escondido.
\textit{322. O que cada resposta do casamento contém?} Agente, habilidade, nome, confiança, similaridade, escore, execuções, episteme e afinidade (ou nulo): nove campos por candidato, todos inspecionáveis. É o dicionário que alimenta as cabeças semântica e de confiança do roteador.
\textit{323. Qual a fronteira do handbook?} Ele cadastra, atualiza e casa; não executa agentes, não valida mérito e não substitui o gate de cobertura. Cadastro e decisão continuam camadas separadas, auditáveis separadamente.
\section{Trigésima quarta rodada: dicionário operacional IX}
\textit{324. Cadastro.} Registro idempotente de habilidade com vetor e confiança inicial $0{,}5$; o que se casa.
\textit{325. Confiança inicial.} $0{,}5$ neutra para habilidade sem histórico; ponto de partida, não avaliação.
\textit{326. Episteme.} Regime epistemológico anotado (explícito ou inferido, \texttt{None} ante falha); auxiliar, não porta.
\textit{327. Afinidade.} Peso brando $\pm10\%$ entre epistemes identificadas; ou nulo sem identificação bilateral.
\textit{328. Filtro.} Confiança mínima $0{,}3$ e interseção de tags; o que precede o ranqueamento.
\textit{329. Composto.} $0{,}6$ similaridade mais $0{,}4$ confiança; fórmula declarada do escore.
\textit{330. Mínimo.} Corte de confiança que exclui habilidades fracas antes de ranquear; declarado na chamada.
\textit{331. Ordenação.} Descendente por escore com corte em \texttt{top\_k}; o ranking que o roteador consome.
\textit{332. Idempotência.} Registrar duas vezes equivale a uma; o que impede duplicatas no cadastro.
\textit{333. Fail-open.} \texttt{None} ante falha de inferência auxiliar; o oposto deliberado do fail-closed das provas, com fronteira declarada.
\section{Trigésima quinta rodada: a episteme em código}
\textit{334. Quais são os seis regimes?} Empírico-analítico, formal-dedutivo, hermenêutico-interpretativo, crítico-reflexivo, pragmático-técnico e regulatório-normativo: cada qual com nome, descrição e léxico de sinais em dois idiomas. Seis regimes, nunca sete.
\textit{335. Como um texto vira regime?} Tokenizando e contando sinais no índice: o dominante é o mais contado, o secundário o segundo (se atingir o mínimo), a confiança é contagem$/4$ limitada a $1{,}0$, e os sinais listados documentam o voto. Empates resolvem-se pela ordem fixa dos regimes --- determinismo até no desempate.
\textit{336. O que é \texttt{MIN\_SINAIS}?} Dois: abaixo disso, \texttt{None} --- ``nunca chuta''. O mínimo é o mecanismo da modéstia: sem evidência lexical suficiente, sem regime atribuído, sem afinidade calculada.
\textit{337. Como a afinidade é calculada?} Matriz entre pares ordenados com neutro $0{,}5$ para desconhecidos e $1{,}0$ para iguais; multiplicada por $0{,}20$ e centrada, produz o peso brando de $\pm10\%$ no máximo. A aritmética do ``brando'' está toda à mostra.
\textit{338. De onde vêm os metadados do agente?} De categoria, tipo, tags, nome e descrição concatenados: inferência sobre metadados já existentes, sem coleta nova. A tarefa contribui só com a descrição.
\textit{339. Por que fail-open aqui e fail-closed lá?} Porque episteme anota e provas decidem: falhar fechado na anotação paralisaria o cadastro por detalhe; falhar aberto devolve \texttt{None} e o escore segue inalterado. A assimetria é deliberada e declarada nos dois cabeçalhos.
\textit{340. O que a cobertura mede?} Função pura sobre definições já carregadas: quantos agentes têm episteme explícita, inferida ou nenhuma --- números medidos, nunca metas. Auditoria do cadastro sem I/O.
\textit{341. Qual a fronteira da episteme?} Léxico, não compreensão: conta sinais, não entende textos. Usá-la além de peso brando seria overclaim do módulo --- por isso o peso é brando e o nulo é comum.
\section{Trigésima sexta rodada: o handbook linha a linha}
\textit{342. O que \texttt{\_expand\_tags} faz?} Quebra identificadores em subtokens para casamento flexível (\texttt{cloudsql-postgres-ad} em partes): sem expansão, só igualdade literal casaria. Flexibilidade declarada com exemplo no código.
\textit{343. Por que tokens do identificador?} Porque o nome do agente carrega semântica (\texttt{Cloud SQL PostgreSQL Specialist} rende quatro tags): enriquecer com o id evita que bons nomes pontuem mal. Deduplicação em ordem preserva a primeira ocorrência.
\textit{344. Como o vetor da habilidade é composto?} Nome, descrição, agente e tags em um texto só, incorporado pelo motor: contexto amplo por desenho, não só descrição isolada.
\textit{345. O que o perfil guarda?} Agente, habilidade, nome, descrição, tags enriquecidas, confiança, vetor, episteme, execuções, sucessos e custo médio: onze campos que sustentam casamento, atualização e auditoria.
\textit{346. Como execuções viram confiança?} Contando sucessos ($0{,}7$ da taxa recente mais $0{,}3$ da anterior) e custo por EMA ($0{,}9/0{,}1$ quando positivo): memória que decai sem apagar, zerada apenas por ausência de histórico ($0{,}5$ inicial).
\textit{347. Qual a ordem do casamento?} Cadastro vazio $\to$ vazio; confiança mínima; interseção de tags; similaridade (ou $0{,}0$ sem vetores); composto; peso brando condicional; arredondamento; ordenação; corte. Nove passos, cada um com condição de saída --- algoritmo legível como contrato.
\textit{348. O que as estatísticas expõem?} Chamadas, acertos de cache e provedor no motor; contagens no handbook: operação observável sem instrumentação externa.
\textit{349. Qual a fronteira do handbook?} Cadastro, atualização e casamento com nove campos por resposta; sem execução de agentes, sem validação de mérito, sem gate de cobertura. Camadas separadas, como sempre.
\section{Trigésima sétima rodada: dicionário operacional X}
\textit{350. Regime.} Um dos seis modos de conhecimento com léxico de sinais; o que se infere ou declara.
\textit{351. Sinal.} Token do léxico que vota num regime; evidência contável, não compreensão.
\textit{352. Mínimo.} Duas ocorrências para atribuir; abaixo, \texttt{None} sem chute.
\textit{353. Secundária.} Segundo regime acima do mínimo; nuance registrada, não decisão.
\textit{354. Confiança epistêmica.} Contagem$/4$ limitada a $1{,}0$; força do voto lexical.
\textit{355. Afinidade.} Matriz com neutro $0{,}5$; $1{,}0$ entre iguais; o que modula o peso brando.
\textit{356. Peso $0{,}20$.} Fator que limita a modulação a $\pm10\%$; a aritmética do ``brando''.
\textit{357. Cobertura.} Contagem de explícitas, inferidas e ausentes; auditoria pura do cadastro.
\textit{358. Puro.} Função sem I/O sobre dados já carregados; o estatuto da medição de cobertura.
\textit{359. Nunca chuta.} \texttt{None} ante sinais insuficientes ou falha; a modéstia em código.
\section{Trigésima oitava rodada: o motor de embeddings}
\textit{360. Como o hash por texto difere do hash por token?} Aqui o texto inteiro é tokenizado e cada token soma $\pm1$ em 8 posições, com normalização L2 ao final: vetor unitário por texto, não por token. Mesma família do legado, outra granularidade --- e dimensão 128, não 64.
\textit{361. Por que o cosseno devolve $0{,}0$ ante vetor nulo?} Porque sem direção não há similaridade: difere do neutro $0{,}5$ do legado, e a diferença está escrita nas duas funções. Dois módulos, duas convenções, ambas declaradas --- confundí-las seria erro de leitura, não do código.
\textit{362. Como o caminho LiteRT produz o vetor?} Cliente com modelo nomeado, vetor da resposta normalizado; dimensão diferente sofre preenchimento com zeros ou corte, depois normalização; ante exceção, hash. Três normalizações possíveis, uma só executada por chamada, todas terminando em vetor unitário.
\textit{363. E o caminho Colibri?} Análogo com limite de 128 tokens, mesma disciplina de dimensão e mesmo recuo ao hash. Simetria entre provedores: trocar de fonte não troca o contrato do vetor.
\textit{364. O que as estatísticas contam?} Chamadas, acertos de cache e provedor efetivo, devolvidas por cópia: observar sem interferir. A cópia impede mutação externa do estado interno.
\textit{365. Quando a sonda escolhe o provedor?} Na construção, uma vez: tenta saúde, fixa o melhor disponível, registra. Escolha estável por instância, não por chamada --- reprodutibilidade do caminho, não só dos números.
\textit{366. O que o cache aproxima?} LRU simples com despejo do mais antigo além de 500: ``aproximação'' escrita no comentário, sem pretensão de otimalidade. Economia declarada como aproximada.
\textit{367. Qual a fronteira do motor?} Produzir vetores unitários determinísticos por texto; não julgar qualidade semântica nem garantir cobertura de vocabulário. Vetor não é compreensão --- e o módulo não alega que seja.
\section{Trigésima nona rodada: funções restantes do handbook}
\textit{368. O que o perfil do agente resume?} Habilidades, confiança média e execuções por agente: leitura agregada do cadastro para inspeção, sem reprocessar o casamento.
\textit{369. O que as estatísticas do handbook contam?} Totais de cadastro e uso: o volume do que foi registrado e consultado, complemento das estatísticas do motor.
\textit{370. Para que serializar em dicionário?} Para persistir e inspecionar o cadastro fora do processo: o handbook vira dado, auditável sem executar o módulo.
\textit{371. O que a cobertura mede, de novo?} Explícitas, inferidas e ausentes sobre definições carregadas, sem I/O: a mesma função pura já descrita, aqui como rotina de auditoria do cadastro.
\textit{372. O que o registro devolve?} O perfil criado (ou o existente, por idempotência) e, no lote, a contagem: cada chamada tem retorno conferível.
\textit{373. O que o casamento devolve vazio?} Lista vazia ante cadastro vazio: sem habilidades, sem candidatos, sem erro. Vazio válido, não exceção.
\textit{374. Onde a episteme aparece na resposta?} Nos campos \texttt{episteme} e \texttt{episteme\_affinity} (ou nulos): anotação por candidato, nunca filtro. Presente para inspecionar, ausente sem punir.
\textit{375. Qual a fronteira do conjunto?} Motor produz vetores, handbook cadastra e casa, singleton compartilha: três papéis, um módulo, nenhuma alegação além de cada docstring.
\section{Quadragésima rodada: dicionário operacional XI}
\textit{376. Unitário.} Vetor de norma 1; saída de toda normalização; o que o cosseno pressupõe.
\textit{377. Granularidade.} Por token no legado, por texto no motor; mesma família, outra soma.
\textit{378. Convenção.} $0{,}0$ ou $0{,}5$ ante degeneração, conforme o módulo; ler a função, não presumir.
\textit{379. Preenchimento.} Zeros até a dimensão; corte além dela; disciplina de dimensão entre provedores.
\textit{380. Cópia.} Devolução por valor das estatísticas; observar sem interferir.
\textit{381. Estável.} Escolha de provedor por instância, não por chamada; reprodutibilidade do caminho.
\textit{382. Aproximação.} LRU simples com despejo do mais antigo; economia declarada como tal.
\textit{383. Vazio válido.} Lista vazia ante cadastro vazio; ausência de candidatos, não erro.
\textit{384. Anotação.} Episteme e afinidade por candidato; inspeção sem punição.
\textit{385. Papéis.} Vetores, cadastro, compartilhamento: três papéis, um módulo, docstrings como fronteira.
\section{Quadragésima primeira rodada: o abstract frase a frase}
\textit{386. Por que o abstract espelha o resumo?} Porque cada sentença inglesa traduz uma portuguesa com os mesmos números e ressalvas: onze lemas, reais exatos sem certificação integral, $0.042\pm0.011$ contra $0.19$--$0.22$, proxy com reescala, intervalos exploratórios, 20\,000 vetores, artefatos com pendências. Espelho, não segundo texto.
\textit{387. Os decimais conferem entre idiomas?} Sim: $0.0421$/$0.0105$/$0.625$/$0.0671$ de um lado, $0{,}0421$/$0{,}0105$/$0{,}625$/$0{,}0671$ do outro; $+0.85$ e $+0{,}85$; $2.2\times10^{-16}$ e $2{,}2\times10^{-16}$. A vírgula muda, os dígitos não --- conferência mecânica, feita aqui uma vez.
\textit{388. As ressalvas sobrevivem à tradução?} Todas: ``exact real arithmetic'', ``do not fully certify'', ``disregard clustering'', ``exploratory'', ``remain pending''. Ressalva traduzida é ressalva mantida; o abstract não alivia o que o resumo restringe.
\textit{389. Por que palavras-chave bilíngues?} Seis termos de cada lado (roteamento/atenção/Lean/verificação/avaliação/reprodutibilidade): indexação nos dois idiomas sem sinonímia criativa. As listas se correspondem termo a termo.
\textit{390. O que o abstract promete a mais?} Nada: mesma questão, mesmos métodos, mesmos números, mesmas pendências. ``A mais'' seria defeito, não virtude, num gênero que resume.
\textit{391. Quem lê só o abstract entende os limites?} Sim, se ler até o fim: a última sentença lista certificação numérica, agrupamento, comparação externa e reprodução como pendentes. Limites no fecho, nos dois idiomas.
\textit{392. Por que manter os dois?} Norma e alcance: resumo em português para o programa, abstract para a comunidade. Dois textos, um conteúdo, zero divergência verificada acima.
\textit{393. Qual a fronteira do gênero?} Resumir sem demonstrar: o abstract aponta para seções, tabelas e apêndices; nada nele se sustenta sozinho. Fronteira respeitada quando cada frase tem endereço --- como percorrido aqui.
\section{Quadragésima segunda rodada: observabilidade}
\textit{394. O que as estatísticas do motor contam?} Chamadas, acertos de cache e provedor efetivo: volume, economia e fonte em três números, devolvidos por cópia a cada consulta.
\textit{395. O que as do handbook contam?} Totais de cadastro e uso: o tamanho do que se registrou e do que se consultou. Motor e handbook observam lados complementares (vetores vs cadastro).
\textit{396. Por que cópia e não referência?} Para observar sem interferir: quem recebe estatísticas não muta o estado interno. Encapsulamento em forma de dicionário copiado.
\textit{397. O que o despejo do mais antigo assume?} Que antiguidade aproxima desuso: LRU simples declarado como aproximação, sem pretensão de otimalidade. Assunção escrita no comentário, não escondida no código.
\textit{398. Sonda única: custo e benefício?} Custa uma checagem de saúde por instância; beneficia com caminho estável e registrado. Paga-se uma vez na construção; colhe-se nas chamadas seguintes --- a economia da decisão estável.
\textit{399. Cache sem invalidação: é seguro?} Textos são imutáveis e o provedor é fixo por instância: mesma entrada, mesma saída, sempre. Sem mutação de fonte, sem necessidade de invalidação --- segurança por imutabilidade, declarada na chave SHA.
\textit{400. O que observar não prova?} Correção semântica, cobertura de vocabulário ou qualidade de casamento: estatísticas contam operação, não mérito. Observar o quanto não diz o quão bem.
\textit{401. Qual a fronteira da observabilidade?} Contar chamadas, acertos e provedores; expor totais de cadastro e uso; nada sobre verdade dos vetores. Números de operação, nunca de mérito.
\section{Quadragésima terceira rodada: dicionário operacional XII}
\textit{402. Espelho.} Abstract que traduz sem aliviar: mesmos dígitos, mesmas ressalvas, outro idioma.
\textit{403. Dígito.} Unidade conferível entre idiomas; vírgula muda, valor não.
\textit{404. Ressalva mantida.} Limite presente nos dois textos; tradução sem alívio.
\textit{405. Chamada.} Incorporação contada nas estatísticas; volume observável.
\textit{406. Acerto.} Entrada encontrada no cache; economia observável.
\textit{407. Provedor efetivo.} Fonte fixada pela sonda e registrada; caminho observável.
\textit{408. Imutabilidade.} Textos que não mudam; o que dispensa invalidação de cache.
\textit{409. Antiguidade.} Critério de despejo do cache; aproximação declarada.
\textit{410. Mérito.} Qualidade semântica não contada por estatísticas; o que observar não prova.
\textit{411. Gênero.} Resumo que aponta sem demonstrar; fronteira do abstract.
\section{Quadragésima quarta rodada: perfil, estatísticas e serialização}
\textit{412. O que o perfil do agente contém?} Identificador, lista de habilidades (id, nome, confiança, execuções, custo médio por habilidade) e confiança média --- ou $0{,}0$ sem habilidades. Cinco campos por habilidade, uma média, um zero declarado: agregado legível sem reprocessar o casamento.
\textit{413. O que as estatísticas do handbook somam?} Total de habilidades, total de agentes, confiança média em quatro casas, provedor do motor, chamadas e acertos de cache: seis chaves que unem cadastro e motor numa só leitura. Vazio devolve zeros, não erro.
\textit{414. Como a serialização nomeia cada perfil?} Chave composta \texttt{agente::habilidade} com nove campos (identidades, nome, descrição, tags, confiança, execuções, sucessos, custo médio), mais o mapa agente-habilidades: persistência completa e legível, sem o vetor (recomputável) e sem o motor.
\textit{415. Por que zeros ante vazio?} Porque ausência de cadastro é estado válido: médias $0{,}0$ e totais zero evitam divisão por zero e exceções. O zero declarado aqui é convenção documentada, não dado.
\textit{416. Ler privados do motor é acoplamento?} Sim, declarado: as estatísticas do handbook leem chamadas e acertos do motor. Acoplamento de leitura para observabilidade unificada, sem escrita cruzada --- fronteira escrita no acesso, não escondida.
\textit{417. Para que expor sucessos e custo médio?} Para auditar a atualização: sucessos sobre execuções conferem a taxa; custo médio confere a EMA. Sem esses campos, a fórmula de atualização seria fé; com eles, é conta conferível.
\textit{418. Serializar sem o vetor perde informação?} Não a essencial: vetor recomputa-se do texto pelo motor; o que se persiste é cadastro, histórico e confiança. Separar recomputável de histórico é desenho, não perda.
\textit{419. Qual a fronteira do trio?} Perfil agrega, estatísticas somam, serialização persiste: leitura, contagem e guarda, sem decidir nada. Nenhuma das três casa, ranqueia ou prova.
\section{Quadragésima quinta rodada: a tese frase a frase}
\textit{420. ``Verificação por máquina''.} Onze teoremas elaborados no núcleo com auditoria por teorema: o primeiro pilar, com endereço na Seção~\ref{cap:provas} e no Apêndice~\ref{ap:lean}.
\textit{421. ``Teste numérico de propriedades''.} Vinte mil vetores em duas sementes contra seis propriedades, com tolerâncias e transcrições: o segundo pilar, tentativa de refutação sobrevivida, nunca prova.
\textit{422. ``Ablação''.} Remoção e perturbação da proxy com diferenças pareadas: o terceiro pilar, que mede dependência sem alegar causalidade.
\textit{423. ``Auditoria de alegações''.} Triagem com decisão humana, citações conferidas, pinos e sincronia: o quarto pilar, que impede a prosa de dizer mais que os artefatos.
\textit{424. ``Padrão de evidência auditável''.} Os quatro pilares amarrados sobre o mesmo objeto com trilha local: composição verificável por terceiro sem executar nada, só lendo.
\textit{425. ``Mesmo sem validação externa''.} A cláusula que dispensa o palco externo sem dispensar o rigor: cada pilar opera dentro do desenho declarado, e as renúncias estão escritas.
\textit{426. ``Teto declarado''.} Otimalidade relativa e evidência preliminar: acima disso, nada se reivindica em nenhum capítulo, resumo ou tabela.
\textit{427. Por que uma frase sustenta uma tese?} Porque cada termo tem seção, tabela ou apêndice: a frase é índice, não argumento. Tese-index, defendida por endereços.
\section{Quadragésima sexta rodada: dicionário operacional XIII}
\textit{428. Agregado.} Perfil com média ou zero; leitura sem reprocessamento.
\textit{429. Serialização.} Dicionário persistível com chave composta; guarda sem motor.
\textit{430. Chave composta.} \texttt{agente::habilidade}; identidade única no cadastro.
\textit{431. Acoplamento.} Leitura de privados para observabilidade unificada; declarado, sem escrita cruzada.
\textit{432. Zero.} Convenção ante vazio (médias, totais); estado válido, não dado.
\textit{433. Média.} Soma sobre contagem com guarda; o que resume sem esconder.
\textit{434. Custo.} EMA por execução positiva; o preço operacional registrado.
\textit{435. Sucesso.} Contagem que sustenta a taxa; o verificável da atualização.
\textit{436. Índice.} Frase que aponta para seções; tese-index, não argumento.
\textit{437. Teto.} Otimalidade relativa e evidência preliminar; acima, nada.
\section{Quadragésima sétima rodada: por verbo da tese}
\textit{438. Verificar.} Elaborar onze teoremas no núcleo com auditoria por teorema; verbo do primeiro pilar, com endereço na Seção~\ref{cap:provas} e revalidação em $191$\,s.
\textit{439. Testar.} Submeter 20\,000 vetores em duas sementes a seis propriedades com tolerâncias; verbo do segundo pilar, tentativa de refutação sobrevivida.
\textit{440. Ablacionar.} Remover e perturbar a proxy com diferenças pareadas; verbo do terceiro pilar, que mede dependência sem alegar causa.
\textit{441. Auditar.} Triar alegações, conferir citações, fixar pinos, amarrar sincronia; verbo do quarto pilar, que impede a prosa de exceder os artefatos.
\textit{442. Compor.} Amarrar os quatro pilares sobre o mesmo objeto com trilha local; verbo da novidade reivindicável, verificável sem executar nada.
\textit{443. Dispensar.} Abrir mão do palco externo sem abrir mão do rigor; verbo da cláusula ``mesmo sem validação externa'', com renúncias escritas.
\textit{444. Declarar.} Fixar o teto (otimalidade relativa, evidência preliminar) em código, delimitação, conclusão e resumo; verbo que circula o documento inteiro.
\textit{445. Indexar.} Apontar cada termo para seção, tabela ou apêndice; verbo deste capítulo, que transforma tese em índice auditável.
\section{Quadragésima oitava rodada: por número}
\textit{446. $0{,}0421$.} Arrependimento médio da atenção em 200 decisões, com margem $0{,}0105$: o numerador da vantagem, chave \texttt{regret\_medio} no JSON.
\textit{447. $0{,}625$.} Taxa de acerto com margem $0{,}0671$: quase dois terços das decisões no ótimo do oráculo, na amostra.
\textit{448. $0{,}85$.} Correlação confiança--qualidade ($0{,}8538$ exatos): a medida da proxy privilegiada, checagem de desenho.
\textit{449. $20\,000$.} Vetores por semente no fuzz, tamanhos e magnitudes declarados: o denominador da tentativa de refutação.
\textit{450. $2{,}2\times10^{-16}$.} Pior desvio de soma nas duas sementes: ordem do épsilon de máquina, abaixo da tolerância $10^{-9}$ por sete ordens.
\textit{451. $0{,}0$.} Pior desvio de deslocamento registrado: invariância exata na amostra, abaixo de $10^{-12}$.
\textit{452. Onze.} Teoremas verificados com três axiomas cada: o numerador formal do trabalho.
\textit{453. Cinco e duzentos.} Piscinas e decisões por coleta: o desenho em dois números, com a limitação do agrupamento declarada.
\textit{454. Quarenta e um.} DOIs com prefixo verificável nas 51 entradas: a medida da integridade referencial.
\textit{455. Sessenta e três, dezesseis, cento e noventa e um.} Testes históricos, reexecutados e segundos da reverificação Lean: três tempos, três escopos, sem extrapolação.
\section{Quadragésima nona rodada: dicionário operacional XIV}
\textit{456. Pilar.} Um dos quatro métodos (verificar, testar, ablacionar, auditar); o que sustenta, com endereço.
\textit{457. Composição.} Amarração dos pilares sobre o mesmo objeto; a novidade reivindicável.
\textit{458. Dispensa.} Renúncia ao palco externo sem renúncia ao rigor; cláusula com renúncias escritas.
\textit{459. Verbo.} Ação da tese (verificar, testar, ablacionar, auditar, compor, dispensar, declarar, indexar); o que se faz, com onde.
\textit{460. Número.} Dígito com chave no apêndice; o que se cita, com fonte.
\textit{461. Rastro.} Caminho seção-tabela-apêndice-código que liga número a origem; auditoria percorrível.
\textit{462. Lastro.} Artefatos locais que sustentam preliminares; o peso por trás da evidência.
\textit{463. Conta.} Cálculo conferível à mão (Ana/Cid, $s_n$, pesos); o que dispensa fé.
\textit{464. Retrato.} Impressão de ambiente datada; onde e quando se mediu.
\textit{465. Guarda.} Persistência com hash e procedência; o que impede deriva silenciosa.
\section{Quinquagésima rodada: por seção do Método}
\textit{466. Provas formais: o que a seção fixa?} Arquivo, ambiente, auditoria e controles do verificador, mais a renúncia ao refinamento completo: o contrato formal do capítulo em um parágrafo.
\textit{467. Bancada e conjuntos: o que distingue?} Três coletas com sementes e sequências próprias (principal, ablações, correlações), pareamento intra-execução e a condição ruidosa binária: o desenho em prosa, sem o qual as tabelas seriam números soltos.
\textit{468. Parâmetro do oráculo: o que demonstra?} A redução entre elegíveis, a invariância do ranking e a fórmula do arrependimento: a seção que impede ler $q_w$ como experimento independente.
\textit{469. Testes e estatística: o que limitam?} Fuzz sem certificação universal, segunda execução não individualizável, intervalos com divisor $n$ sem agrupamento, Pearson sem causalidade: quatro limites que governam toda leitura numérica adiante.
\textit{470. Ameaças e ética: o que renunciam?} Construção arbitrária com proxy privilegiada, confusão residual, validade só interna e sem participantes humanos: a seção que diz o que o desenho não é.
\textit{471. Desenvolvimento: o que documenta?} SPECs, RED--GREEN, axiomas, falhas e a ressalva de não-reexecução: história do fazer, não refazimento.
\textit{472. Protocolo operacional: o que permite?} Reexecutar cada artefato com comandos e resultados esperados dos registros, reportando divergências contra os registros: reprodução como procedimento, não como promessa.
\textit{473. Roteiros: para que tantos?} Para que cada peça (provas, bancada, fuzz, pinos, JSONs, Lean, scripts, roteador, testes, cabeças, incorporador, casamento, handbook, episteme, motor, abstract, perfil, tese, verbos) tenha inspeção guiada sem execução: o Método como manual de auditoria, não só relato.
\section{Quinquagésima primeira rodada: \texttt{route} e \texttt{explain} em código}
\textit{474. Qual a assinatura de \texttt{route}?} Descrição, requisitos, cartões e índice posicional opcional; devolve ranking ordenado do melhor ao pior com pesos somando 1. Contrato em quatro entradas e uma saída ordenada.
\textit{475. Por que o índice é descartado?} Por compatibilidade de API, deliberadamente sem participar do escore: decisões idênticas produzem o mesmo ranking. Descarte documentado, não esquecimento.
\textit{476. Quantas vezes se avalia por chamada?} Uma: \texttt{route} delega a \texttt{\_evaluate} e extrai o ranking. Avaliação única evita divergência entre decidir e explicar.
\textit{477. O que \texttt{explain} acrescenta?} Tarefa e requisitos ecoados mais todo o dicionário da avaliação: a decisão com seu contexto de entrada, pronta para inspeção.
\textit{478. Ecoar a entrada serve para quê?} Para vincular saída a entrada no mesmo objeto: sem o eco, a explicação seria órfã da tarefa que a gerou.
\textit{479. Onde mora o determinismo?} No índice fixo, no desempate por identificador e na ausência de aleatoriedade no caminho: três fontes, zero sorteio. Determinismo triplo, declarado em três pontos.
\textit{480. Para que servem as docstrings?} Como especificação executável legível: ranking ordenado, pesos somando 1, avaliação única, exposição total. Documentação que o código cumpre por construção.
\textit{481. Qual a fronteira da dupla?} Decidir e expor; não aprender, não executar agentes, não validar mérito. Duas funções, um contrato, as mesmas fronteiras do roteador.
\section{Quinquagésima segunda rodada: dicionário operacional XV}
\textit{482. Capítulo.} Unidade de tese (Introdução a Conclusão na matéria); o que se conta em laudas.
\textit{483. Parágrafo.} Unidade de prosa com uma tese e seu endereço; o que se audita por frase.
\textit{484. Legenda.} Texto que diz o que a tabela mostra e como lê-la; leitura correta embutida.
\textit{485. Fonte.} Linha de elaboração própria com ano e base; honestidade tipográfica padronizada.
\textit{486. Nota.} Campo que explica nulos e ressalvas dentro do artefato; metadado que impede misleitura.
\textit{487. Remissão.} \texttt{Seção}/\texttt{Tabela}/\texttt{Apêndice} numerados automaticamente; o fio da rastreabilidade.
\textit{488. Citação.} Chave autor-data com entrada conferida; 51 citadas, 51 entradas.
\textit{489. Verbete.} Entrada de glossário ou dicionário com definição operacional; o átomo deste capítulo.
\textit{490. Bloco.} Seção temática de arguição; cinquenta e quatro neste capítulo e contando.
\textit{491. Item.} Pergunta numerada com resposta endereçada; 465 até aqui, sem nenhum dado novo.
\section{Quinquagésima terceira rodada: por roteiro}
\textit{492. Roteiro de provas.} Guia a leitura do arquivo Lean (ordem, hipóteses, táticas, ausência de atalhos) e da auditoria: o que conferir antes de crer num veredito.
\textit{493. Roteiro de bancada.} Guia sementes, conjuntos, pareamento e condição ruidosa: o desenho em procedimento, para não ler tabelas como mágica.
\textit{494. Roteiro de testes e JSONs.} Guia chaves, amostras, nulos com nota e transcrições: conferência sem execução, número a número.
\textit{495. Roteiro Lean e de scripts.} Guia arquivo de 150 linhas e cinco scripts com linhas e funções: anatomia antes de juízo.
\textit{496. Roteiro do roteador e do pipeline.} Guia gates, avaliação única, ordenação, \texttt{explain}, varredura e motivos de gap: decisão e sincronia inspecionáveis.
\textit{497. Roteiro de testes e transcrições.} Guia oito verificações, 159 linhas e três tempos: o que foi reexecutado, o que é histórico e quanto custou cada.
\textit{498. Roteiro de cabeças e incorporador.} Guia matcher, recuos, neutros, sanitizador, dimensão e pesos textuais: semântica explícita peça por peça.
\textit{499. Para que dezesseis roteiros?} Para que nenhuma peça dependa de fé: cada artefato tem seu guia de inspeção sem execução. Dezesseis portas de entrada, um mesmo padrão.
\section{Quinquagésima quarta rodada: \texttt{explain} em uso e fronteiras}
\textit{500. Quem chama \texttt{explain}?} A suíte de integração (múltiplas chamadas com tarefas e cartões, incluindo teste dedicado que confere excluídos, cabeças, utilidade e ranking na explicação); os scripts da bancada usam \texttt{route} e \texttt{\_evaluate} direto. Uso por testes, não por produção --- honestidade de cobertura.
\textit{501. O que o teste dedicado confere?} Que a explicação inclui excluídos com motivos, escores por cabeça, utilidade e ranking: o dicionário completo sob asserção, não só sob docstring.
\textit{502. Por que a bancada não usa \texttt{explain}?} Porque precisa de linhas tabulares (escolha, arrependimento, acerto, latência, soma), não do dicionário completo: usa \texttt{route} para decidir e campos de \texttt{\_evaluate} para registrar. Economia de formato, mesma avaliação única.
\textit{503. Uma ou duas avaliações por decisão?} Uma: \texttt{route} e \texttt{explain} delegam ao mesmo \texttt{\_evaluate}; decidir e expor nunca divergem por recomputação. A garantia é estrutural (delegação), não disciplinar.
\textit{504. O eco pode mentir?} Não por construção: ecoa os argumentos recebidos junto ao resultado da mesma avaliação. Mentira exigiria adulterar a delegação --- visível em três linhas.
\textit{505. \texttt{explain} aprende algo?} Nada: sem estado, sem escrita, sem chamada a provedor novo. Função pura sobre a avaliação --- a mais auditável do módulo.
\textit{506. Por que docstring como especificação?} Porque ranking ordenado, soma unitária e exposição total são verificáveis por teste (e verificados): a docstring promete, a suíte cobra.
\textit{507. Qual a fronteira final?} Decidir e expor o decidido; jamais executar, aprender ou validar. A fronteira do módulo inteiro, repetida na última função.
\section{Quinquagésima quinta rodada: dicionário operacional XVI}
\textit{508. Roteiro.} Guia de inspeção sem execução por artefato; dezesseis neste Método.
\textit{509. Inspeção.} Leitura guiada com o que conferir; auditoria sem rodar.
\textit{510. Manual.} Método como manual de auditoria, não só relato; o que os roteiros fazem dele.
\textit{511. Auditor.} Leitor com roteiro e apêndices; o papel para o qual tudo foi escrito.
\textit{512. Pergunta.} Itens numerados deste capítulo; 491 até a rodada anterior.
\textit{513. Resposta.} Prosa com endereço em seção, tabela ou apêndice; sem endereço, sem resposta.
\textit{514. Marco.} Contagem redonda com balanço (80, 100, 138...); pausa para conferir a numeração.
\textit{515. Suplemento.} Capítulo que guia e aprofunda sem introduzir dados; o gênero deste.
\textit{516. Travessia.} Percurso pergunta a pergunta por um artefato; o método do suplemento.
\textit{517. Amarração.} Ligação entre texto, código e registro que impede deriva; o que tudo aqui faz.
\section{Quinquagésima sexta rodada: por classe r212}
\textit{518. O que a classe tudo-ou-nada verifica?} Que agente com todas as capacidades recebe a tarefa e com só uma é excluído: o gate de cobertura sob asserção, nos dois sentidos.
\textit{519. E a de confiança fail-closed?} Que permitido é selecionado com negado excluído, e que todos negados deixam a tarefa bloqueada sem atalho: negação que bloqueia em vez de degradar.
\textit{520. O que os extremos normalizados verificam?} Que todas as cabeças ficam em $[0,1]$ nos extremos positivo e negativo da semântica: o intervalo sob estresse, não só no típico.
\textit{521. E os pesos?} Padrão como combinação convexa; inválidos rejeitados ou normalizados (inclusive por subclasse que força pesos ruins): o validador sob ataque de teste.
\textit{522. O que o determinismo verifica?} Repetição com mesma entrada e independência do índice posicional: duas faces da mesma garantia, cada qual com seu teste.
\textit{523. O que ``carga viva'' verifica?} Cartão que reporta carga real de tarefas atribuídas e preferência pelo menos carregado em igualdade nas demais: carga que vive, não constante de enfeite.
\textit{524. E o contrato do \texttt{explain}?} Inclusão de excluídos, cabeças, utilidade e ranking na explicação: o dicionário completo sob asserção dedicada.
\textit{525. Quantos testes e o que cobrem?} Dezenove em seis classes: gate, confiança, normalização, determinismo, carga e explicação. Cobertura por comportamento, não por linha.
\section{Quinquagésima sétima rodada: a partição em detalhe}
\textit{526. Não-dicionário entra?} Não: vira \texttt{invalid\_card} com índice posicional. Tipo errado não quebra a partição --- é classificado.
\textit{527. Sem identificador?} \texttt{missing\_agent\_id} com índice: sem id não há candidato, mas há registro.
\textit{528. Duplicata: qual fica?} Nenhuma: a segunda ocorrência remove a primeira e registra \texttt{duplicate\_agent\_id}. Ambiguidade não se resolve por ordem --- elimina-se.
\textit{529. Espaços no id?} Removidos nas pontas antes de tudo: identidade normalizada antes de comparada.
\textit{530. Por que copiar o cartão?} Para normalizar o id sem mutar a entrada: o elegível carrega id limpo, o original segue intacto. Higiene sem efeito colateral.
\textit{531. O que volta quando nada é elegível?} Listas vazias com pesos declarados: o caminho do gate vazio, sem exceção. Vazio válido de novo, agora no roteador.
\textit{532. Para que o mapa de excluídos?} Auditoria por decisão: cada fora com seus motivos nominais. Sem o mapa, o gate seria caixa-preta; com ele, é relatório.
\textit{533. Qual a fronteira da partição?} Separar elegíveis de excluídos com motivos; não pontuar, não ordenar, não decidir. Uma função, um trabalho.
\section{Quinquagésima oitava rodada: dicionário operacional XVII}
\textit{534. Tudo-ou-nada.} Gate de cobertura nos dois sentidos; recebe com todas, exclui sem uma.
\textit{535. Negação.} Agente negado que bloqueia sem atalho; fail-closed na seleção.
\textit{536. Extremo.} Entrada de estresse que confina cabeças em $[0,1]$; intervalo sob pressão.
\textit{537. Combinação.} Pesos não negativos somando 1; o validado e o atacado.
\textit{538. Repetição.} Mesma entrada, mesmo ranking; determinismo testado.
\textit{539. Carga viva.} Carga reportada de tarefas reais; preferência em igualdade.
\textit{540. Contrato.} Dicionário completo sob asserção; explicação cobrada por teste.
\textit{541. Classificação.} Tipo errado que vira motivo, não exceção; partição que não quebra.
\textit{542. Normalizado.} Id limpo em cópia, original intacto; higiene sem efeito.
\textit{543. Uma função.} Partição que só separa com motivos; fronteira de uma frase.
\section{Sexagésima rodada: por motivo de gap}
\textit{544. \texttt{arquivo\_ausente}: quando?} Leitura que devolve \texttt{None} (inexistente ou ilegível), no alvo ou na referência: sem texto, sem checagem, com gap registrado e detalhe do caminho.
\textit{545. \texttt{faltante}: quando?} Esperado fora dos achados: o nome que o plano exige e o artefato não contém. Falta cobrada uma a uma, com o nome no detalhe.
\textit{546. \texttt{excedente}: quando?} Achado fora dos esperados: o nome que o artefato contém e o plano não prevê. Adição silenciosa vira gap nominal --- a anti-deriva por acréscimo.
\textit{547. \texttt{proibido}: quando?} Padrão banido presente no texto, por regex própria: o veto executável (atalhos de prova, por exemplo). Presença que reprova, mesmo com todo o resto certo.
\textit{548. \texttt{cobertura}: quando?} Nome sem correspondente na referência cruzada: existe aqui, falta lá. A checagem relacional entre dois artefatos, não dentro de um só.
\textit{549. Como gaps se agregam?} Por \texttt{id}: checagens com qualquer gap contam como falhadas; o veredito soma passadas e falhadas sobre o total. Agregação por identidade, sem dupla contagem.
\textit{550. E o plano inválido?} \texttt{ValueError} antes de qualquer leitura: não-objeto, conjuntos vazios, tipo fora do par, arquivo ou regex ausentes, referência malformada. Plano inspecionado antes dos artefatos --- meta-auditoria embutida.
\textit{551. Qual a fronteira dos motivos?} Cinco classes fechadas, cada qual com \texttt{id} e detalhe; fora delas, nada se registra. Taxonomia que não admite ``outros''.
\section{Sexagésima primeira rodada: motivos de gate em detalhe}
\textit{552. Status indisponível: o que acontece?} Motivo \texttt{status\_not\_available} e exclusão: só \texttt{available} compete. Disponibilidade como pré-condição, não como escore.
\textit{553. Capacidades em string única?} Envolvidas em tupla antes do conjunto: robustez a formato sem afrouxar a exigência. Formato tolerado, conteúdo cobrado.
\textit{554. Capacidades inválidas?} Conjunto vazio mais motivo \texttt{invalid\_capabilities}: tipo quebrado vira exclusão explicada, nunca exceção.
\textit{555. Faltantes listadas como?} Ordenadas e nominais em \texttt{missing\_capabilities:...}: o candidato sabe exatamente o que lhe falta. Exclusão didática.
\textit{556. Por que a duplicata elimina ambas?} Porque identidade ambígua não se resolve por ordem: remover as duas evita que posição decida elegibilidade. Regra anti-ambiguidade, não punição.
\textit{557. Onde vão os motivos?} No mapa de excluídos por identificador: cada fora com suas razões nominais. Auditoria por decisão, preservada no \texttt{explain} e conferida por teste.
\textit{558. Cópia normalizada: por quê?} Id limpo sem mutar a entrada: higiene com o original intacto. O elegível carrega identidade tratada; a fonte segue como veio.
\textit{559. Qual a fronteira dos motivos?} Explicar exclusões, uma a uma, com nomes; nunca pontuar, ordenar ou decidir. Gate que relata, não que julga.
\section{Sexagésima segunda rodada: dicionário operacional XVIII}
\textit{560. Ausência.} Arquivo ilegível ou inexistente; gap sem checagem, com caminho no detalhe.
\textit{561. Falta.} Esperado não achado; cobrança nominal uma a uma.
\textit{562. Excesso.} Achado não esperado; anti-deriva por acréscimo.
\textit{563. Veto.} Padrão banido presente; presença que reprova.
\textit{564. Cruzada.} Checagem relacional entre dois artefatos; cobertura, não conteúdo.
\textit{565. Agregação.} Soma por \texttt{id} sem dupla contagem; veredito sobre o total.
\textit{566. Meta.} Plano inspecionado antes dos artefatos; auditoria da auditoria.
\textit{567. Nominal.} Motivo com nomes (faltantes, caminhos, padrões); explicação endereçável.
\textit{568. Indisponível.} Status fora de \texttt{available}; pré-condição, não escore.
\textit{569. Relato.} Gate que explica exclusões; julgar, nunca.
\section{Sexagésima terceira rodada: \texttt{\_checar\_nomes} linha a linha}
\textit{570. O que a função valida primeiro?} Que esperados e proibidos sejam listas: contrato de entrada checado antes de qualquer leitura. Tipo errado aborta com \texttt{ValueError} nomeando a checagem.
\textit{571. De onde vêm os achados?} De \texttt{\_nomes} com a regex do item: nomes únicos na ordem de aparição. Uma só extração serve aos três testes seguintes --- faltante, excedente e proibido bebem da mesma fonte.
\textit{572. Como o faltante é cobrado?} Esperado fora dos achados vira gap com o nome no detalhe: cobrança nominal uma a uma, sem agregação que esconda qual faltou.
\textit{573. E o excedente?} Achado fora dos esperados (com coerção a texto na comparação): adição silenciosa vira gap nominal. A coerção evita falsos excedentes por tipo (número vs texto).
\textit{574. Como o proibido é testado?} Compilando cada padrão (regex inválida aborta nomeando a checagem) e buscando no texto: presença vira gap com o padrão no detalhe. Veto com mensagem que diz qual padrão mordeu.
\textit{575. Qual o formato do gap?} Dicionário com \texttt{id}, motivo e detalhe, sempre: três chaves, sem variação por motivo. Formato fechado facilita agregação e leitura tabular.
\textit{576. Onde a função pode abortar?} Em listas malformadas e regex proibida inválida: dois \texttt{ValueError} nomeando \texttt{cid}. Aborto com endereço, nunca silencioso.
\textit{577. Qual a fronteira da função?} Comparar nomes contra listas dentro de um texto; não ler arquivos (recebe texto), não julgar mérito, não cruzar referências (isso é a cobertura). Uma função, um texto, três testes.
\section{Sexagésima quarta rodada: cobertura, nomes e leitura}
\textit{578. O que a cobertura valida primeiro?} A referência como dicionário com arquivo e regex textuais: sem referência válida, \texttt{ValueError} nomeando a checagem. Cruzar exige dois lados bem formados.
\textit{579. E se a referência não abre?} Gap \texttt{arquivo\_ausente} com o caminho: ausência do outro lado também é gap, não aborto. Simetria de tratamento entre alvo e referência.
\textit{580. Como cada nome é testado?} Extraído do alvo e procurado nos nomes da referência: sem correspondente, gap de cobertura com o nome. Relacional um a um, sem agregação que esconda qual.
\textit{581. Por que a regex exige grupo?} Porque extração sem grupo capturado não nomeia nada: \texttt{\_nomes} recusa padrão sem ao menos um grupo e padrão inválido, ambos com \texttt{ValueError}. Contrato da extração declarado na entrada.
\textit{582. Por que deduplicar em ordem?} Porque ``excedente'' compara conjuntos e o relatório lê melhor na ordem de aparição: unicidade para a lógica, ordem para o humano. Uma decisão, dois beneficiários.
\textit{583. Quando a leitura devolve \texttt{None}?} Ante erro de sistema ou de decodificação: arquivo ausente ou ilegível viram \texttt{None}, nunca exceção. O chamador decide (gap), o leitor não julga.
\textit{584. Como o veredito é somado?} Total de checagens, passadas menos falhadas por \texttt{id}, lista de gaps, rótulo e bandeira: aritmética por identidade, sem dupla contagem. Cinco chaves mais rótulo e bandeira, sempre.
\textit{585. Qual a fronteira do trio?} Ler sem julgar, extrair com contrato, cruzar com referência: três funções, cada qual com sua validação de entrada e sua forma de falha. Falha nomeada em todo caminho.
\section{Sexagésima quinta rodada: dicionário operacional XIX}
\textit{586. Lista.} Contrato de entrada (esperados, proibidos); tipo checado antes de ler.
\textit{587. Achado.} Nome extraído único em ordem; a fonte dos três testes.
\textit{588. Coerção.} Conversão a texto na comparação; sem falsos excedentes por tipo.
\textit{589. Veto.} Padrão banido com mensagem do padrão; presença que reprova.
\textit{590. Relacional.} Teste entre dois artefatos; cobertura, não conteúdo.
\textit{591. Detalhe.} Nome, caminho ou padrão no gap; o que permite agir.
\textit{592. Identidade.} \texttt{id} que agrega vereditos; sem dupla contagem.
\textit{593. Soma.} Total, passadas e falhadas; aritmética do veredito.
\textit{594. Bandeira.} \texttt{exige\_validacao\_externa} sempre verdadeira; inspeção requerida.
\textit{595. Caminho.} Falha nomeada em toda validação; aborto com endereço, nunca silêncio.

% Módulo 07 — Conclusão (R762+)
\chapter{Conclusão}\label{cap:conclusao}
Retomando a questão da Seção~\ref{cap:intro}, registramos o que foi demonstrado, o que foi medido e o que permanece aberto.
A questão da pesquisa foi respondida em dois níveis: os registros documentam onze lemas do modelo matemático formalizados em Lean 4 e uma avaliação preliminar da implementação no simulador, com artefatos locais de reprodução. Reivindica-se otimalidade relativa e evidência preliminar. Contribuições numeradas: (C1) onze lemas com provas verificadas por máquina e axiomas auditados; (C2) pipeline de verificação lake-env com bloqueio a atalhos; (C3) bancada com intervalos exploratórios, pareamento e variação de escala do oráculo; (C4) testes aleatórios de propriedades numéricas e auditoria de alegações e referências; (C5) trilha de reprodução com pinos, hashes e protocolo de sincronia. Trabalhos futuros: Q/K/V formal completo, erro numérico certificado à Flocq, análise com agrupamento por piscina, temperatura com limites, benchmark externo com suíte compartilhada e orçamento próprio.
Os objetivos de formalização (OE1), comparação interna (OE2) e auditoria documental (OE3) são sustentados pelos registros locais nos limites explicitados. O objetivo de disponibilização pública (OE4) permanece parcialmente atendido: os artefatos estão organizados localmente, mas o depósito público e sua identificação persistente estão pendentes.
A Tabela~\ref{tab:futuro} sintetiza os requisitos de continuidade da pesquisa.
\begin{table}[htbp]
\fontsize{10}{12}\selectfont\SingleSpacing
\centering
\caption{Programa de trabalhos futuros (pacotes com aceite).}
\label{tab:futuro}
\begin{tabularx}{\textwidth}{lXX}
\toprule
Pacote & Método & Aceite \\
\midrule
Q/K/V formal & Estruturas + refinamento & machine-checked \\
Erro certificado & Intervalos + Flocq & cotas verificadas \\
Temperatura limite & Análise assintótica & teoremas de limite \\
Benchmark externo & Suíte + orçamento & head-to-head executado \\
Reprodução independente & Terceiros + pins & mesmos números \\
Seleção parcial & Heap + benchmark & mesmos top-\textit{k}, menor latência \\
\bottomrule
\end{tabularx}
\fonte{Elaboração própria (2026), a partir dos artefatos e das fontes identificados no texto.}
\end{table}

\postextual
\begin{SingleSpacing}
\raggedright
% Referencias: sistema autor-data; ABNT NBR 6023:2025.
\begin{thebibliography}{99}

\bibitem[Bagnara et al.(2019)]{bagnara2019} BAGNARA, R. et al. \textit{Correct Approximation of IEEE 754 Floating-Point Arithmetic for Program Verification}. [S. l.]: arXiv, 2019. Preprint arXiv:1903.06119. DOI: 10.48550/arXiv.1903.06119. Disponível em: \url{https://arxiv.org/abs/1903.06119}. Acesso em: 10 out. 2026.

\bibitem[Basili, Selby e Hutchens(1986)]{basili1986} BASILI, V. R.; SELBY, R. W.; HUTCHENS, D. H. Experimentation in software engineering. \textit{IEEE Transactions on Software Engineering}, [s. l.], v. SE-12, n. 7, p. 733--743, 1986. DOI: 10.1109/TSE.1986.6312975. Disponível em: \url{https://doi.org/10.1109/TSE.1986.6312975}. Acesso em: 10 out. 2026.

\bibitem[Beckenbauer et al.(2025)]{orchestrator2025} BECKENBAUER, L. et al. \textit{Orchestrator: Active Inference for Multi-Agent Systems in Long-Horizon Tasks}. [S. l.]: arXiv, 2025. Preprint arXiv:2509.05651. DOI: 10.48550/arXiv.2509.05651. Disponível em: \url{https://arxiv.org/abs/2509.05651}. Acesso em: 10 out. 2026.

\bibitem[Bridle(1990)]{bridle1990} BRIDLE, J. S. Probabilistic Interpretation of Feedforward Classification Network Outputs, with Relationships to Statistical Pattern Recognition. In: SOULIÉ, F. F.; HÉRAULT, J. (ed.). \textit{Neurocomputing}. Berlin; Heidelberg: Springer, 1990. p. 227--236. (NATO ASI Series, v. 68). DOI: 10.1007/978-3-642-76153-9\_28. Disponível em: \url{https://doi.org/10.1007/978-3-642-76153-9_28}. Acesso em: 10 out. 2026.

\bibitem[Chidamber e Kemerer(1994)]{chidamber1994} CHIDAMBER, S. R.; KEMERER, C. F. A metrics suite for object oriented design. \textit{IEEE Transactions on Software Engineering}, [s. l.], v. 20, n. 6, p. 476--493, 1994. DOI: 10.1109/32.295895. Disponível em: \url{https://doi.org/10.1109/32.295895}. Acesso em: 10 out. 2026.

\bibitem[Collberg e Proebsting(2016)]{collberg2016} COLLBERG, C.; PROEBSTING, T. A. Repeatability in computer systems research. \textit{Communications of the ACM}, [s. l.], v. 59, n. 3, p. 62--69, 2016. DOI: 10.1145/2812803. Disponível em: \url{https://doi.org/10.1145/2812803}. Acesso em: 10 out. 2026.

\bibitem[Cumming(2014)]{cumming2014} CUMMING, G. The New Statistics: Why and How. \textit{Psychological Science}, [s. l.], v. 25, n. 1, p. 7--29, 2014. DOI: 10.1177/0956797613504966. Disponível em: \url{https://doi.org/10.1177/0956797613504966}. Acesso em: 10 out. 2026.

\bibitem[Dai et al.(2019)]{dai2019} DAI, Z. et al. Transformer-XL: Attentive Language Models beyond a Fixed-Length Context. In: ANNUAL MEETING OF THE ASSOCIATION FOR COMPUTATIONAL LINGUISTICS, 57., 2019, Florence. \textit{Proceedings of the 57th Annual Meeting of the Association for Computational Linguistics}. Florence: Association for Computational Linguistics, 2019. p. 2978--2988. DOI: 10.18653/v1/P19-1285. Disponível em: \url{https://doi.org/10.18653/v1/P19-1285}. Acesso em: 10 out. 2026.

\bibitem[De Moura e Bjørner(2008)]{demoura2008} DE MOURA, L.; BJØRNER, N. Z3: An Efficient SMT Solver. In: INTERNATIONAL CONFERENCE ON TOOLS AND ALGORITHMS FOR THE CONSTRUCTION AND ANALYSIS OF SYSTEMS, 14., 2008, Budapest. \textit{Tools and Algorithms for the Construction and Analysis of Systems}. Berlin; Heidelberg: Springer, 2008. Lecture Notes in Computer Science, v. 4963. p. 337--340. DOI: 10.1007/978-3-540-78800-3\_24. Disponível em: \url{https://doi.org/10.1007/978-3-540-78800-3_24}. Acesso em: 10 out. 2026.

\bibitem[De Moura e Ullrich(2021)]{moura2021} DE MOURA, L.; ULLRICH, S. The Lean 4 Theorem Prover and Programming Language. In: INTERNATIONAL CONFERENCE ON AUTOMATED DEDUCTION, 28., 2021, evento virtual. \textit{Automated Deduction -- CADE 28}. Cham: Springer, 2021. Lecture Notes in Computer Science, v. 12699. p. 625--635. DOI: 10.1007/978-3-030-79876-5\_37. Disponível em: \url{https://doi.org/10.1007/978-3-030-79876-5_37}. Acesso em: 10 out. 2026.

\bibitem[Deng et al.(2009)]{deng2009} DENG, J. et al. ImageNet: A large-scale hierarchical image database. In: IEEE CONFERENCE ON COMPUTER VISION AND PATTERN RECOGNITION, 2009, Miami. \textit{2009 IEEE Conference on Computer Vision and Pattern Recognition}. [S. l.]: IEEE, 2009. p. 248--255. DOI: 10.1109/CVPR.2009.5206848. Disponível em: \url{https://doi.org/10.1109/CVPR.2009.5206848}. Acesso em: 10 out. 2026.

\bibitem[Devlin et al.(2019)]{devlin2019} DEVLIN, J. et al. BERT: Pre-training of Deep Bidirectional Transformers for Language Understanding. In: CONFERENCE OF THE NORTH AMERICAN CHAPTER OF THE ASSOCIATION FOR COMPUTATIONAL LINGUISTICS: HUMAN LANGUAGE TECHNOLOGIES, 2019, Minneapolis. \textit{Proceedings of the 2019 Conference of the North American Chapter of the Association for Computational Linguistics: Human Language Technologies, Volume 1 (Long and Short Papers)}. Minneapolis: Association for Computational Linguistics, 2019. p. 4171--4186. DOI: 10.18653/v1/N19-1423. Disponível em: \url{https://doi.org/10.18653/v1/N19-1423}. Acesso em: 10 out. 2026.

\bibitem[Efron(1979)]{efron1979} EFRON, B. Bootstrap Methods: Another Look at the Jackknife. \textit{The Annals of Statistics}, [s. l.], v. 7, n. 1, p. 1--26, 1979. DOI: 10.1214/aos/1176344552. Disponível em: \url{https://doi.org/10.1214/aos/1176344552}. Acesso em: 10 out. 2026.

\bibitem[Erman et al.(1980)]{erman1980} ERMAN, L. D. et al. The Hearsay-II Speech-Understanding System: Integrating Knowledge to Resolve Uncertainty. \textit{ACM Computing Surveys}, [s. l.], v. 12, n. 2, p. 213--253, 1980. DOI: 10.1145/356810.356816. Disponível em: \url{https://doi.org/10.1145/356810.356816}. Acesso em: 10 out. 2026.

\bibitem[Fagan(1976)]{fagan1976} FAGAN, M. E. Design and code inspections to reduce errors in program development. \textit{IBM Systems Journal}, [s. l.], v. 15, n. 3, p. 182--211, 1976. DOI: 10.1147/sj.153.0182. Disponível em: \url{https://doi.org/10.1147/sj.153.0182}. Acesso em: 10 out. 2026.

\bibitem[Gebru et al.(2021)]{gebru2021} GEBRU, T. et al. Datasheets for datasets. \textit{Communications of the ACM}, [s. l.], v. 64, n. 12, p. 86--92, 2021. DOI: 10.1145/3458723. Disponível em: \url{https://doi.org/10.1145/3458723}. Acesso em: 10 out. 2026.

\bibitem[Gneiting e Raftery(2007)]{gneiting2007} GNEITING, T.; RAFTERY, A. E. Strictly Proper Scoring Rules, Prediction, and Estimation. \textit{Journal of the American Statistical Association}, [s. l.], v. 102, n. 477, p. 359--378, 2007. DOI: 10.1198/016214506000001437. Disponível em: \url{https://doi.org/10.1198/016214506000001437}. Acesso em: 10 out. 2026.

\bibitem[Goldberg(1991)]{goldberg1991} GOLDBERG, D. What every computer scientist should know about floating-point arithmetic. \textit{ACM Computing Surveys}, [s. l.], v. 23, n. 1, p. 5--48, 1991. DOI: 10.1145/103162.103163. Disponível em: \url{https://doi.org/10.1145/103162.103163}. Acesso em: 10 out. 2026.

\bibitem[Harris et al.(2020)]{harris2020} HARRIS, C. R. et al. Array programming with NumPy. \textit{Nature}, [s. l.], v. 585, n. 7825, p. 357--362, 2020. DOI: 10.1038/s41586-020-2649-2. Disponível em: \url{https://doi.org/10.1038/s41586-020-2649-2}. Acesso em: 10 out. 2026.

\bibitem[Hoare(1969)]{hoare1969} HOARE, C. A. R. An axiomatic basis for computer programming. \textit{Communications of the ACM}, [s. l.], v. 12, n. 10, p. 576--580, 1969. DOI: 10.1145/363235.363259. Disponível em: \url{https://doi.org/10.1145/363235.363259}. Acesso em: 10 out. 2026.

\bibitem[Ioannidis(2005)]{ioannidis2005} IOANNIDIS, J. P. A. Why Most Published Research Findings Are False. \textit{PLoS Medicine}, [s. l.], v. 2, n. 8, e124, 2005. DOI: 10.1371/journal.pmed.0020124. Disponível em: \url{https://doi.org/10.1371/journal.pmed.0020124}. Acesso em: 10 out. 2026.

\bibitem[Jaynes(1957)]{jaynes1957} JAYNES, E. T. Information Theory and Statistical Mechanics. \textit{Physical Review}, [s. l.], v. 106, n. 4, p. 620--630, 1957. DOI: 10.1103/PhysRev.106.620. Disponível em: \url{https://doi.org/10.1103/PhysRev.106.620}. Acesso em: 10 out. 2026.

\bibitem[Komisarenko e Kull(2024)]{komisarenko2024} KOMISARENKO, V.; KULL, M. \textit{Improving Calibration by Relating Focal Loss, Temperature Scaling, and Properness}. [S. l.]: arXiv, 2024. Preprint arXiv:2408.11598. DOI: 10.48550/arXiv.2408.11598. Disponível em: \url{https://arxiv.org/abs/2408.11598}. Acesso em: 10 out. 2026.

\bibitem[Kullback e Leibler(1951)]{kullback1951} KULLBACK, S.; LEIBLER, R. A. On Information and Sufficiency. \textit{The Annals of Mathematical Statistics}, [s. l.], v. 22, n. 1, p. 79--86, 1951. DOI: 10.1214/aoms/1177729694. Disponível em: \url{https://doi.org/10.1214/aoms/1177729694}. Acesso em: 10 out. 2026.

\bibitem[Liu et al.(2023)]{agentbench2023} LIU, X. et al. \textit{AgentBench: Evaluating LLMs as Agents}. [S. l.]: arXiv, 2023. Preprint arXiv:2308.03688. DOI: 10.48550/arXiv.2308.03688. Disponível em: \url{https://arxiv.org/abs/2308.03688}. Acesso em: 10 out. 2026.

\bibitem[Luong, Pham e Manning(2015)]{luong2015} LUONG, T.; PHAM, H.; MANNING, C. D. Effective Approaches to Attention-based Neural Machine Translation. In: CONFERENCE ON EMPIRICAL METHODS IN NATURAL LANGUAGE PROCESSING, 2015, Lisbon. \textit{Proceedings of the 2015 Conference on Empirical Methods in Natural Language Processing}. Lisbon: Association for Computational Linguistics, 2015. p. 1412--1421. DOI: 10.18653/v1/D15-1166. Disponível em: \url{https://doi.org/10.18653/v1/D15-1166}. Acesso em: 10 out. 2026.

\bibitem[McCabe(1976)]{mccabe1976} MCCABE, T. J. A Complexity Measure. \textit{IEEE Transactions on Software Engineering}, [s. l.], v. SE-2, n. 4, p. 308--320, 1976. DOI: 10.1109/TSE.1976.233837. Disponível em: \url{https://doi.org/10.1109/TSE.1976.233837}. Acesso em: 10 out. 2026.

\bibitem[Meurer et al.(2017)]{meurer2017} MEURER, A. et al. SymPy: symbolic computing in Python. \textit{PeerJ Computer Science}, [s. l.], v. 3, e103, 2017. DOI: 10.7717/peerj-cs.103. Disponível em: \url{https://doi.org/10.7717/peerj-cs.103}. Acesso em: 10 out. 2026.

\bibitem[Miller, Fredriksen e So(1990)]{miller1990} MILLER, B. P.; FREDRIKSEN, L.; SO, B. An empirical study of the reliability of UNIX utilities. \textit{Communications of the ACM}, [s. l.], v. 33, n. 12, p. 32--44, 1990. DOI: 10.1145/96267.96279. Disponível em: \url{https://doi.org/10.1145/96267.96279}. Acesso em: 10 out. 2026.

\bibitem[Mitchell et al.(2019)]{mitchell2019} MITCHELL, M. et al. Model Cards for Model Reporting. In: CONFERENCE ON FAIRNESS, ACCOUNTABILITY, AND TRANSPARENCY, 2019, Atlanta. \textit{Proceedings of the Conference on Fairness, Accountability, and Transparency}. New York: ACM, 2019. p. 220--229. DOI: 10.1145/3287560.3287596. Disponível em: \url{https://doi.org/10.1145/3287560.3287596}. Acesso em: 10 out. 2026.

\bibitem[Muller et al.(2018)]{muller2018} MULLER, J.-M. et al. \textit{Handbook of Floating-Point Arithmetic}. 2. ed. Cham: Birkhäuser, 2018. DOI: 10.1007/978-3-319-76526-6. Disponível em: \url{https://doi.org/10.1007/978-3-319-76526-6}. Acesso em: 10 out. 2026.

\bibitem[Murray(2026)]{murray2026} MURRAY, T. \textit{Lipschitz-Based Robustness Certification Under Floating-Point Execution}. [S. l.]: arXiv, 2026. Preprint arXiv:2603.13334. DOI: 10.48550/arXiv.2603.13334. Disponível em: \url{https://arxiv.org/abs/2603.13334}. Acesso em: 10 out. 2026.

\bibitem[Open Science Collaboration(2015)]{osc2015} OPEN SCIENCE COLLABORATION. Estimating the reproducibility of psychological science. \textit{Science}, [s. l.], v. 349, n. 6251, aac4716, 2015. DOI: 10.1126/science.aac4716. Disponível em: \url{https://doi.org/10.1126/science.aac4716}. Acesso em: 10 out. 2026.

\bibitem[Page et al.(2021)]{page2021} PAGE, M. J. et al. The PRISMA 2020 statement: an updated guideline for reporting systematic reviews. \textit{BMJ}, [s. l.], v. 372, n71, 2021. DOI: 10.1136/bmj.n71. Disponível em: \url{https://doi.org/10.1136/bmj.n71}. Acesso em: 10 out. 2026.

\bibitem[Parnas e Clements(1986)]{parnas1986} PARNAS, D. L.; CLEMENTS, P. C. A rational design process: How and why to fake it. \textit{IEEE Transactions on Software Engineering}, [s. l.], v. SE-12, n. 2, p. 251--257, 1986. DOI: 10.1109/TSE.1986.6312940. Disponível em: \url{https://doi.org/10.1109/TSE.1986.6312940}. Acesso em: 10 out. 2026.

\bibitem[Rajpurkar et al.(2016)]{rajpurkar2016} RAJPURKAR, P. et al. SQuAD: 100,000+ Questions for Machine Comprehension of Text. In: CONFERENCE ON EMPIRICAL METHODS IN NATURAL LANGUAGE PROCESSING, 2016, Austin. \textit{Proceedings of the 2016 Conference on Empirical Methods in Natural Language Processing}. Austin: Association for Computational Linguistics, 2016. p. 2383--2392. DOI: 10.18653/v1/D16-1264. Disponível em: \url{https://doi.org/10.18653/v1/D16-1264}. Acesso em: 10 out. 2026.

\bibitem[Runeson e Höst(2009)]{runeson2009} RUNESON, P.; HÖST, M. Guidelines for conducting and reporting case study research in software engineering. \textit{Empirical Software Engineering}, [s. l.], v. 14, n. 2, p. 131--164, 2009. DOI: 10.1007/s10664-008-9102-8. Disponível em: \url{https://doi.org/10.1007/s10664-008-9102-8}. Acesso em: 10 out. 2026.

\bibitem[Shannon(1948)]{shannon1948} SHANNON, C. E. A Mathematical Theory of Communication. \textit{Bell System Technical Journal}, [s. l.], v. 27, n. 3, p. 379--423, 1948. DOI: 10.1002/j.1538-7305.1948.tb01338.x. Disponível em: \url{https://doi.org/10.1002/j.1538-7305.1948.tb01338.x}. Acesso em: 10 out. 2026.

\bibitem[Smith(1980)]{smith1980} SMITH, R. G. The Contract Net Protocol: High-Level Communication and Control in a Distributed Problem Solver. \textit{IEEE Transactions on Computers}, [s. l.], v. C-29, n. 12, p. 1104--1113, 1980. DOI: 10.1109/TC.1980.1675516. Disponível em: \url{https://doi.org/10.1109/TC.1980.1675516}. Acesso em: 10 out. 2026.

\bibitem[Sun et al.(2026)]{perspectivegap2026} SUN, Y. et al. \textit{PerspectiveGap: A Benchmark for Multi-Agent Orchestration Prompting}. [S. l.]: arXiv, 2026. Preprint arXiv:2606.08878. DOI: 10.48550/arXiv.2606.08878. Disponível em: \url{https://arxiv.org/abs/2606.08878}. Acesso em: 10 out. 2026.

\bibitem[The Mathlib Community(2020)]{mathlib2020} THE MATHLIB COMMUNITY. The lean mathematical library. In: ACM SIGPLAN INTERNATIONAL CONFERENCE ON CERTIFIED PROGRAMS AND PROOFS, 9., 2020, New Orleans. \textit{Proceedings of the 9th ACM SIGPLAN International Conference on Certified Programs and Proofs}. New York: ACM, 2020. p. 367--381. DOI: 10.1145/3372885.3373824. Disponível em: \url{https://doi.org/10.1145/3372885.3373824}. Acesso em: 10 out. 2026.

\bibitem[Tsang et al.(2026)]{orchrm2026} TSANG, K. Y. et al. \textit{Reward Modeling for Multi-Agent Orchestration}. [S. l.]: arXiv, 2026. Preprint arXiv:2606.13598. DOI: 10.48550/arXiv.2606.13598. Disponível em: \url{https://arxiv.org/abs/2606.13598}. Acesso em: 10 out. 2026.

\bibitem[Vaswani et al.(2017)]{vaswani2017} VASWANI, A. et al. Attention is all you need. In: ANNUAL CONFERENCE ON NEURAL INFORMATION PROCESSING SYSTEMS, 31., 2017, Long Beach. \textit{Advances in Neural Information Processing Systems}. [S. l.]: Curran Associates, 2017. v. 30. p. 5998--6008. Disponível em: \url{https://proceedings.neurips.cc/paper_files/paper/2017/hash/3f5ee243547dee91fbd053c1c4a845aa-Abstract.html}. Acesso em: 10 out. 2026.

\bibitem[Virtanen et al.(2020)]{virtanen2020} VIRTANEN, P. et al. SciPy 1.0: fundamental algorithms for scientific computing in Python. \textit{Nature Methods}, [s. l.], v. 17, n. 3, p. 261--272, 2020. DOI: 10.1038/s41592-019-0686-2. Disponível em: \url{https://doi.org/10.1038/s41592-019-0686-2}. Acesso em: 10 out. 2026.

\bibitem[Wang et al.(2018)]{wang2018} WANG, A. et al. GLUE: A Multi-Task Benchmark and Analysis Platform for Natural Language Understanding. In: EMNLP WORKSHOP BLACKBOXNLP: ANALYZING AND INTERPRETING NEURAL NETWORKS FOR NLP, 2018, Brussels. \textit{Proceedings of the 2018 EMNLP Workshop BlackboxNLP: Analyzing and Interpreting Neural Networks for NLP}. Brussels: Association for Computational Linguistics, 2018. p. 353--355. DOI: 10.18653/v1/W18-5446. Disponível em: \url{https://doi.org/10.18653/v1/W18-5446}. Acesso em: 10 out. 2026.

\bibitem[Wasserstein e Lazar(2016)]{wasserstein2016} WASSERSTEIN, R. L.; LAZAR, N. A. The ASA Statement on <i>p</i> -Values: Context, Process, and Purpose. \textit{The American Statistician}, [s. l.], v. 70, n. 2, p. 129--133, 2016. DOI: 10.1080/00031305.2016.1154108. Disponível em: \url{https://doi.org/10.1080/00031305.2016.1154108}. Acesso em: 10 out. 2026.

\bibitem[Wilkinson et al.(2016)]{wilkinson2016} WILKINSON, M. D. et al. The FAIR Guiding Principles for scientific data management and stewardship. \textit{Scientific Data}, [s. l.], v. 3, n. 1, 160018, 2016. DOI: 10.1038/sdata.2016.18. Disponível em: \url{https://doi.org/10.1038/sdata.2016.18}. Acesso em: 10 out. 2026.

\bibitem[Wohlin et al.(2012)]{wohlin2012} WOHLIN, C. et al. \textit{Experimentation in Software Engineering}. Berlin; Heidelberg: Springer, 2012. DOI: 10.1007/978-3-642-29044-2. Disponível em: \url{https://doi.org/10.1007/978-3-642-29044-2}. Acesso em: 10 out. 2026.

\bibitem[Wolf et al.(2020)]{wolf2020} WOLF, T. et al. Transformers: State-of-the-Art Natural Language Processing. In: CONFERENCE ON EMPIRICAL METHODS IN NATURAL LANGUAGE PROCESSING: SYSTEM DEMONSTRATIONS, 2020, evento virtual. \textit{Proceedings of the 2020 Conference on Empirical Methods in Natural Language Processing: System Demonstrations}. [S. l.]: Association for Computational Linguistics, 2020. p. 38--45. DOI: 10.18653/v1/2020.emnlp-demos.6. Disponível em: \url{https://doi.org/10.18653/v1/2020.emnlp-demos.6}. Acesso em: 10 out. 2026.

\bibitem[Xie et al.(2024)]{xie2024} XIE, J. et al. \textit{Calibrating Language Models with Adaptive Temperature Scaling}. [S. l.]: arXiv, 2024. Preprint arXiv:2409.19817. DOI: 10.48550/arXiv.2409.19817. Disponível em: \url{https://arxiv.org/abs/2409.19817}. Acesso em: 10 out. 2026.

\bibitem[Xuan, Yang e Li(2025)]{xuan2025} XUAN, H.; YANG, B.; LI, X. \textit{Exploring the Impact of Temperature Scaling in Softmax for Classification and Adversarial Robustness}. [S. l.]: arXiv, 2025. Preprint arXiv:2502.20604. DOI: 10.48550/arXiv.2502.20604. Disponível em: \url{https://arxiv.org/abs/2502.20604}. Acesso em: 10 out. 2026.

\end{thebibliography}

\end{SingleSpacing}
\chapter*{Glossário}\addcontentsline{toc}{chapter}{GLOSSÁRIO}\label{ap:glossario}
\begin{description}[style=nextline,leftmargin=0pt,labelindent=0pt]
\item[Ablação.] Alteração ou remoção de um componente para comparar seu efeito dentro de um desenho definido.
\item[Agente.] Programa ao qual uma tarefa pode ser atribuída; o estudo com cartões calcula rankings e não executa os agentes.
\item[Arrependimento.] Diferença entre a melhor utilidade do oráculo e a utilidade da escolha realizada.
\item[Axioma.] Proposição admitida pelo sistema formal. A auditoria registra \texttt{propext}, \texttt{Classical.choice} e \texttt{Quot.sound}, sem demonstrar minimalidade desse conjunto.
\item[Calibração probabilística.] Correspondência entre probabilidades previstas e frequências observadas; não é garantida apenas pela soma unitária do softmax.
\item[Combinação convexa.] Soma ponderada com pesos não negativos de soma unitária.
\item[Comparador.] Estratégia utilizada como referência para avaliação no mesmo desenho.
\item[Intervalo exploratório.] Faixa calculada pela aproximação normal registrada, sem correção para agrupamento nas piscinas deste estudo.
\item[Maximizador.] Elemento de conjunto finito não vazio que atinge o maior valor de uma função; T3 exige finitude e não-vazio como hipóteses.
\item[Oráculo.] Critério sintético usado para avaliar as escolhas; não representa qualidade observada em tarefas reais.
\item[Overflow.] Excesso da faixa representável em ponto flutuante. T11 prova um limite da exponencial em reais exatos e não certifica todas as operações da implementação.
\item[Reprodução.] Execução dos procedimentos por outro executor a partir dos artefatos e do ambiente especificado; reprodução independente permanece pendente.
\item[Roteador.] Mecanismo que filtra candidatos e ordena suas utilidades para selecionar um agente.
\item[Semente.] Valor que inicializa uma sequência pseudoaleatória.
\item[Softmax.] Normalização exponencial de uma família finita não vazia de escores. Em reais exatos, seus pesos somam um e são positivos; em ponto flutuante, underflow pode produzir zeros.
\item[Temperatura.] Parâmetro positivo que controla a concentração dos pesos. T10 utiliza pelo menos dois candidatos e maximizador único como hipóteses suficientes.
\item[Top-\textit{k}.] Prefixo de tamanho $k$ de um ranking; T9 pressupõe a ordenação decrescente da lista.
\item[Underflow.] Resultado muito pequeno para ser representado com a precisão ou a faixa disponíveis.
\item[Verificação por máquina.] Checagem de um termo de prova pelo núcleo do assistente, relativa ao enunciado e aos axiomas utilizados; não garante que o enunciado modele integralmente o sistema concreto.
\end{description}

% Módulo 09 — Apêndices (R762+; código integral + trilhas com hashes)
\apendices
\chapter{Código Lean integral}
\label{ap:lean}
\lstinputlisting{/home/marceloclaro/opencode-ecosystem-core/plugins/lean-nucleo-verificacao/teoremas/roteamento_atencao.lean}
\chapter{Trilha de reprodução}
\label{ap:reproducao}
Lean 4 v4.34.1, Lake 5.0.0, elan 4.2.4; Mathlib rev.\ v4.34.1 (4760 oleans). Bancada principal: sementes 11, 22, 33, 44 e 55; 40 tarefas por semente. Ablações: sementes efetivas 10011, 10022, 10033, 10044 e 10055. Os códigos e transcrições reproduzidos são artefatos históricos; a execução integral da fonte atual não foi revalidada nesta revisão editorial. Comandos: \nolinkurl{lake env lean RoteamentoAtencao.lean} (saída 0); \texttt{pytest} 9 suítes (63 aprovados). Sincronia automatizada entre teoremas, testes, SPEC e axiomas: \nolinkurl{scanners/diagnostic_pipeline} (SPEC-935-R767, 8 testes herméticos).
A Tabela~\ref{tab:pins} reproduz os identificadores históricos dos repositórios e das licenças.
\begin{table}[htbp]
\fontsize{10}{12}\selectfont\SingleSpacing
\centering
\caption{Pinos verificados por \texttt{git ls-remote} (2026-10-09) e SHA-256 das licenças.}
\label{tab:pins}
\begin{tabularx}{\textwidth}{lXX}
\toprule
Repositório & Commit HEAD (8) & SHA-256 licença (8) \\
\midrule
leanprover/lean4 & 5e96ee11 & 8b28515f \\
mathlib4 & a02a63c2 & b40930bb \\
batteries & b71aaf60 & c71d239d \\
aesop & 04879b1e & c71d239d \\
Z3Prover/z3 & 0d0e28c3 & e617cad2 \\
sympy/sympy & 04a7da3a & 07a5e981 \\
\bottomrule
\end{tabularx}
\fonte{Elaboração própria (2026), a partir dos artefatos e das fontes identificados no texto.}
\end{table}
Depósito público dos artefatos pendente do autor. ORCID: \url{https://orcid.org/0000-0001-8996-2887}. Contato: \texttt{marceloclaro@gmail.com}, \url{https://bit.ly/geomaker}.
\chapter{Código dos experimentos e do pipeline}
\label{ap:codigo}
\lstinputlisting{/home/marceloclaro/opencode-ecosystem-core/plugins/lean-nucleo-verificacao/bancada/bancada_routing.py}
\lstinputlisting{/home/marceloclaro/opencode-ecosystem-core/plugins/lean-nucleo-verificacao/bancada/bancada_ablacao.py}
\lstinputlisting{/home/marceloclaro/opencode-ecosystem-core/plugins/lean-nucleo-verificacao/bancada/correlacoes.py}
\lstinputlisting{/home/marceloclaro/opencode-ecosystem-core/plugins/lean-nucleo-verificacao/bancada/contraprova_float.py}
\lstinputlisting{/home/marceloclaro/opencode-ecosystem-core/scanners/diagnostic_pipeline.py}
\lstinputlisting{/home/marceloclaro/opencode-ecosystem-core/transformer/attention.py}
\lstinputlisting{/home/marceloclaro/opencode-ecosystem-core/integrations/lean_proof_mcp.py}
\lstinputlisting{/home/marceloclaro/opencode-ecosystem-core/tests/test_r747_nucleo_lean.py}
\chapter{Resultados em JSON}
\label{ap:json}
\lstinputlisting{/home/marceloclaro/opencode-ecosystem-core/plugins/lean-nucleo-verificacao/bancada/resultado_bancada.json}
\lstinputlisting{/home/marceloclaro/opencode-ecosystem-core/plugins/lean-nucleo-verificacao/bancada/resultado_ablacao_q0p3.json}
\lstinputlisting{/home/marceloclaro/opencode-ecosystem-core/plugins/lean-nucleo-verificacao/bancada/resultado_ablacao_q0p6.json}
\lstinputlisting{/home/marceloclaro/opencode-ecosystem-core/plugins/lean-nucleo-verificacao/bancada/resultado_ablacao_q0p9.json}
\lstinputlisting{/home/marceloclaro/opencode-ecosystem-core/plugins/lean-nucleo-verificacao/bancada/correlacoes.json}
\lstinputlisting{/home/marceloclaro/opencode-ecosystem-core/plugins/lean-nucleo-verificacao/bancada/per_seed.json}
\chapter{Transcrição da suíte de testes}
\label{ap:testes}
\lstinputlisting{/home/marceloclaro/opencode-ecosystem-core/artigos/pucrs-roteamento-atencao/transcricao_testes.txt}
\lstinputlisting{/home/marceloclaro/opencode-ecosystem-core/artigos/pucrs-roteamento-atencao/transcricao_pytest_verbose.txt}
\lstinputlisting{/home/marceloclaro/opencode-ecosystem-core/artigos/pucrs-roteamento-atencao/impressao_ambiente.txt}
\chapter{Sessão do verificador (tabelas)}
\label{ap:sessao}
Transcrição real de \texttt{lean\_verify} (compilador presente): \texttt{sorry} $\to$ inválido (\texttt{incomplete\_sorry}); \texttt{admit} $\to$ inválido; \texttt{axiom} declarado $\to$ inválido (\texttt{axiom\_nao\_declarado}); \texttt{theorem t : 2 + 2 = 4 := by decide} $\to$ válido (\texttt{machine\_checked}); código com \texttt{import Mathlib.Tactic} fora de projeto lake $\to$ válido (\texttt{syntax\_verified}, sem alegar compilação). A ausência de \texttt{decide}/\texttt{trivial}/\texttt{rfl}/\texttt{assumption} no vocabulário de táticas, detectada nesta revisão, foi corrigida: sem o compilador no \texttt{PATH}, o mesmo teorema caía em \texttt{syntax\_error} em vez do correto \texttt{syntax\_verified}. Segunda sessão (R763): \texttt{div\_lt\_div\_iff} inexistente na Mathlib pinada (forma vigente \texttt{div\_lt\_div\_iff₀}, consultada no fonte) e \texttt{Nontrivial} exigido pelo compilador em T10, pois o singleton falsifica a nitidez --- ambas corrigidas em rascunho antes de incorporar.
\begin{table}[htbp]
\fontsize{10}{12}\selectfont\SingleSpacing
\centering
\caption{Micro-provas por tática, todas verificadas (núcleo via verificador direto; aritmética via projeto lake).}
\label{tab:microprovas-ap}
\begin{tabularx}{\textwidth}{lXX}
\toprule
Táticas & Miniatura provada & Veredito \\
\midrule
\texttt{intro} + \texttt{exact} & simetria da igualdade & machine-checked \\
\texttt{rw} & reescrita por hipótese & machine-checked \\
\texttt{simp} & identidade aditiva & machine-checked \\
\texttt{omega} & sucessor maior que o número & machine-checked \\
\texttt{decide} & $2+2=4$ por computação & machine-checked \\
\texttt{cases} & booleanos esgotados & machine-checked \\
\texttt{trivial} & verdade imediata & machine-checked \\
\texttt{rfl} & igualdade por reflexividade & machine-checked \\
\texttt{assumption} & hipótese do contexto & machine-checked \\
\texttt{apply} & regra com unificação & machine-checked \\
\texttt{induction} & indução nos naturais & machine-checked \\
\texttt{contradiction} & absurdo fecha tudo & machine-checked \\
\texttt{have} & lema local intermediário & machine-checked \\
\texttt{obtain} & extração de testemunha & machine-checked \\
\texttt{norm\_num} & aritmética real (via projeto) & machine-checked \\
\bottomrule
\end{tabularx}
\fonte{Elaboração própria (2026), a partir dos artefatos e das fontes identificados no texto.}
\end{table}
\begin{table}[htbp]
\fontsize{10}{12}\selectfont\SingleSpacing
\centering
\caption{Diário de erros reais e correções (mensagens verídicas do compilador).}
\label{tab:erros-ap}
\begin{tabularx}{\textwidth}{lXX}
\toprule
Item & Mensagem ou sintoma & Correção \\
\midrule
Notação & \texttt{unexpected token 'in'} na v4.34 & reescrito com \texttt{Finset.sum} explícito \\
Assinatura & \nolinkurl{Application type mismatch} em soma unitária & elemento explícito antes da pertinência \\
Vocabulário & quatro táticas do núcleo ausentes; \texttt{syntax\_error} indevido & nomes acrescentados \\
Hipótese & \nolinkurl{failed to synthesize Nontrivial} em T10 & hipótese adicionada (singleton falsificaria a tese) \\
Nome & \texttt{Unknown identifier `div\_lt\_div\_iff'} & forma vigente consultada no fonte pinado \\
Silencioso & saída de auditoria com \texttt{sorryAx} em T7 & erro de aplicação corrigido; pego só pela auditoria \\
\bottomrule
\end{tabularx}
\fonte{Elaboração própria (2026), a partir dos artefatos e das fontes identificados no texto.}
\end{table}
\chapter{Histórico de auditoria}
\label{ap:historico}
Ciclos desta linhagem (registro compartilhado; homônimos legados com escores $\approx0{,}9$ pertencem a outra linhagem e foram excluídos pelo filtro de escore e tema):
A Tabela~\ref{tab:historico} apresenta os ciclos do registro evolutivo, cujos escores são internos.
\begin{SingleSpacing}
\fontsize{10}{12}\selectfont
\begin{longtable}{p{1.1cm}p{1.2cm}p{12.2cm}}
\caption{Ciclos de evolução da linhagem Lean-roteamento (registros históricos).}\label{tab:historico}\\
\toprule
Ciclo & Escore & Objeto registrado \\
\midrule
\endfirsthead
\toprule
Ciclo & Escore & Objeto registrado \\
\midrule
\endhead
R747 & 8.5 & \parbox[t]{12cm}{\raggedright\footnotesize Núcleo Lean R747: 5 agentes, 2 skills, MCP lean\_proof\_mcp e plugin verificação total} \\
R748 & 7.5 & \parbox[t]{12cm}{\raggedright\footnotesize Sequência polimata Lean: manifesto + auditoria clones + orientação pinagem viva Z3/SymPy} \\
R749 & 8.0 & \parbox[t]{12cm}{\raggedright\footnotesize Autorização e validação R749: dossiê extensão R711 Lean sem federar} \\
R750 & 8.2 & \parbox[t]{12cm}{\raggedright\footnotesize Validação viva somente-leitura R750: API+LICENSE dos 4 Lean} \\
R751 & 8.8 & \parbox[t]{12cm}{\raggedright\footnotesize Execução R751: ls-remote 6 commits + sha256 6 LICENSE + pins+propor} \\
R752 & 9.0 & \parbox[t]{12cm}{\raggedright\footnotesize Aprovação R752: extensão allowlist 21→25 Lean + proposta 6/0} \\
R753 & 8.7 & \parbox[t]{12cm}{\raggedright\footnotesize Install/build R753: elan+stable v4.34.1 e fallback Mathlib-ausente honesto} \\
R754 & 9.2 & \parbox[t]{12cm}{\raggedright\footnotesize R754: lake new math + cache Mathlib + lake-env machine\_checked + fallback tática-Mathlib} \\
R755 & 9.3 & \parbox[t]{12cm}{\raggedright\footnotesize R755: teorema eficiência roteamento-atenção 3/3 machine\_checked} \\
R756 & 9.4 & \parbox[t]{12cm}{\raggedright\footnotesize R756: pacote 1+2+3 — invariância, Ioo, bilateral, 6/6 machine\_checked} \\
R757 & 8.5 & \parbox[t]{12cm}{\raggedright\footnotesize R757: bancada interna + 2 revisões internas automatizadas + relatório honesto} \\
R758 & 9.5 & \parbox[t]{12cm}{\raggedright\footnotesize R758: máxima eficiência — T7+axiomas, fuzz 20k, ablações qw, sensibilidade} \\
R759 & 9.0 & \parbox[t]{12cm}{\raggedright\footnotesize Artigo PUCRS: revalidação exaustiva + PDF 5p compilado} \\
R760 & 9.0 & \parbox[t]{12cm}{\raggedright\footnotesize Artigo PUCRS ampliado: autoria+teoria+impactos, claim limpo, PDF 6p} \\
R761 & 9.4 & \parbox[t]{12cm}{\raggedright\footnotesize R761 T8 monotonicidade mais T9 top-k mais artigo 9 lemas} \\
R762 & 9.1 & \parbox[t]{12cm}{\raggedright\footnotesize Dissertacao ABNT 35p: 7 capitulos + QKV + 9 teoremas + estatistica total + Lean integral} \\
R763 & 9.4 & \parbox[t]{12cm}{\raggedright\footnotesize R763: T10 temperatura + T11 limite da exponencial deslocada + correlacoes Pearson + benchmarks externos mapeados} \\
R764 & 9.3 & \parbox[t]{12cm}{\raggedright\footnotesize R764: correcao minuciosa total + sincronia 11 + recalculos} \\
R765 & 9.2 & \parbox[t]{12cm}{\raggedright\footnotesize R765: dissertacao modular 9 modulos + 12 refs + teoria ampliada} \\
R766 & 8.8 & \parbox[t]{12cm}{\raggedright\footnotesize R766: revisao por gaps — pipeline inexistente + varredura total sem derivas} \\
R767 & 9.2 & \parbox[t]{12cm}{\raggedright\footnotesize R767: diagnostic\_pipeline + gap real pego em estreia} \\
R768 & 9.3 & \parbox[t]{12cm}{\raggedright\footnotesize R768: pesquisa refeita — bateria fresca + docs 63-9 + PDFs} \\
R769 & 8.9 & \parbox[t]{12cm}{\raggedright\footnotesize R769: apendice didatico leigo + maswos\_pipeline inexistente} \\
R770 & 9.2 & \parbox[t]{12cm}{\raggedright\footnotesize R770: modulos expandidos — enunciados integrais + pareada + ABNT vigente} \\
R771 & 9.0 & \parbox[t]{12cm}{\raggedright\footnotesize R771: versao arXiv — 11 lemas com provas + contexto didatico} \\
R772 & 8.9 & \parbox[t]{12cm}{\raggedright\footnotesize R772: preprint integralmente em PT-BR} \\
R773 & 9.4 & \parbox[t]{12cm}{\raggedright\footnotesize R773: dissertacao 80 laudas — 82p modular ABNT vigente} \\
R774 & 9.0 & \parbox[t]{12cm}{\raggedright\footnotesize R774: sessao Lean demonstrativa + vocabulario de taticas corrigido + micro-provas verificadas} \\
R775 & 9.1 & \parbox[t]{12cm}{\raggedright\footnotesize R775: preliminares matematicos + 04b tratado por teorema + T10/T11 com hipoteses e limites} \\
R776 & 9.0 & \parbox[t]{12cm}{\raggedright\footnotesize R776: consolidacoes 02d+02h+02e e 04c em tabela no apendice + demonstracao sobre cartoes} \\
R777 & 9.2 & \parbox[t]{12cm}{\raggedright\footnotesize R777: revisao de 10 itens + medicao honesta (100p totais; materia pp.\ 12--38) + 41 DOIs verificados} \\
R778 & 8.8 & \parbox[t]{12cm}{\raggedright\footnotesize R778: capitulo de arguicoes aprofundado (40 itens + leitura guiada, sem dados novos) + materia pp.\ 12--46} \\
R779 & 8.9 & \parbox[t]{12cm}{\raggedright\footnotesize R779: rodadas Q41--Q58 + protocolo operacional no metodo + materia pp.\ 12--49} \\
R780 & 9.0 & \parbox[t]{12cm}{\raggedright\footnotesize R780: rodadas Q59--Q80 (registros brutos, pacotes, familias, tabelas) + roteiro JSON + 116p, materia pp.\ 13--53} \\
R781 & 9.0 & \parbox[t]{12cm}{\raggedright\footnotesize R781: rodadas Q81--Q110 (por capitulo + dicionario) + roteiro Lean + 119p, materia pp.\ 13--56} \\
R782 & 9.0 & \parbox[t]{12cm}{\raggedright\footnotesize R782: rodadas Q111--Q138 (scripts, apendices, dicionario II) + roteiro Python + 122p, materia pp.\ 13--59} \\
R783 & 9.0 & \parbox[t]{12cm}{\raggedright\footnotesize R783: rodadas Q139--Q166 (roteador/MCP/pipeline, SPECs, dicionario III) + 124p, materia pp.\ 13--61} \\
R784 & 9.0 & \parbox[t]{12cm}{\raggedright\footnotesize R784: rodadas Q167--Q196 (estudo/pipeline, MCP/testes, dicionario IV) + 127p, materia pp.\ 13--64} \\
R785 & 9.0 & \parbox[t]{12cm}{\raggedright\footnotesize R785: rodadas Q197--Q222 (testes, transcricoes, dicionario V) + pior shift zero + 129p, materia pp.\ 13--66} \\
R786 & 9.0 & \parbox[t]{12cm}{\raggedright\footnotesize R786: rodadas Q223--Q250 (cabecas, ablacoes, dicionario VI) + roteiro cabecas + 131p, materia pp.\ 13--68} \\
R787 & 9.1 & \parbox[t]{12cm}{\raggedright\footnotesize R787: reverificacao Lean 4 (4/4 em 191s) + rodadas Q251--Q281 + 133p, materia pp.\ 13--70} \\
R788 & 9.0 & \parbox[t]{12cm}{\raggedright\footnotesize R788: rodadas Q282--Q307 (resumo, matcher, dicionario VIII) + 134p, materia pp.\ 13--71} \\
R789 & 9.0 & \parbox[t]{12cm}{\raggedright\footnotesize R789: rodadas Q308--Q333 (conclusao, handbook, dicionario IX) + 137p, materia pp.\ 14--74} \\
R790 & 9.0 & \parbox[t]{12cm}{\raggedright\footnotesize R790: rodadas Q334--Q359 (episteme, handbook, dicionario X) + 139p, materia pp.\ 14--76} \\
R791 & 9.0 & \parbox[t]{12cm}{\raggedright\footnotesize R791: rodadas Q360--Q385 (engine, funcoes, dicionario XI) + 141p, materia pp.\ 14--78} \\
R792 & 9.0 & \parbox[t]{12cm}{\raggedright\footnotesize R792: rodadas Q386--Q411 (abstract, observabilidade, dicionario XII) + 142p, materia pp.\ 14--79} \\
R793 & 9.0 & \parbox[t]{12cm}{\raggedright\footnotesize R793: rodadas Q412--Q437 (perfil/stats/serializacao, tese, dicionario XIII) + 143p, materia pp.\ 14--80} \\
R794 & 9.0 & \parbox[t]{12cm}{\raggedright\footnotesize R794: rodadas Q438--Q465 (verbos, numeros, dicionario XIV) + 144p, materia pp.\ 14--81} \\
R795 & 9.0 & \parbox[t]{12cm}{\raggedright\footnotesize R795: rodadas Q466--Q491 (metodo, route/explain, dicionario XV) + 146p, materia pp.\ 14--83} \\
R796 & 9.0 & \parbox[t]{12cm}{\raggedright\footnotesize R796: rodadas Q492--Q517 (roteiros, explain em uso, dicionario XVI) + 148p, materia pp.\ 14--85} \\
R797 & 9.0 & \parbox[t]{12cm}{\raggedright\footnotesize R797: rodadas Q518--Q543 (r212, particao, dicionario XVII) + 149p, materia pp.\ 14--86} \\
R798 & 9.0 & \parbox[t]{12cm}{\raggedright\footnotesize R798: rodadas Q544--Q569 (gaps, gates, dicionario XVIII) + 150p, materia pp.\ 14--87} \\
R799 & 9.0 & \parbox[t]{12cm}{\raggedright\footnotesize R799: rodadas Q570--Q595 (checagens, cobertura, dicionario XIX) + 153p, materia pp.\ 15--90} \\
\bottomrule
\end{longtable}
\noindent\fonte{Elaboração própria (2026), a partir do registro evolutivo local. Os escores históricos são autoavaliações internas e não equivalem a validação externa.}
\end{SingleSpacing}
\chapter{Especificações formais}
\label{ap:specs}
\lstinputlisting{/home/marceloclaro/opencode-ecosystem-core/specs/SPEC-935-R755-teorema-roteamento-atencao.md}
\lstinputlisting{/home/marceloclaro/opencode-ecosystem-core/specs/SPEC-935-R754-lake-env-machine-checked.md}
\lstinputlisting{/home/marceloclaro/opencode-ecosystem-core/specs/SPEC-935-R747-nucleo-lean-verificacao-total.md}
\lstinputlisting{/home/marceloclaro/opencode-ecosystem-core/specs/SPEC-935-R767-diagnostic-pipeline.md}

\end{document}
